% Learn with the assistance of computers — ICT Upper Sixth — Cours signé OnBuch+
% Official MINESEC syllabus. Compiler avec : tectonic ICT14-learn-with-the-assistance-of-computers.tex
\newcommand{\DOCMATIERE}{ICT}
\newcommand{\DOCNIVEAU}{Upper Sixth}
\newcommand{\DOCMODULE}{Upper Sixth — ICT}
\newcommand{\DOCLECON}{14}
\newcommand{\DOCTITRE}{Learn with the assistance of computers}
% OnBuch+ — préambule « cours signés » ENRICHI (compilé avec tectonic / XeLaTeX)
% Système visuel « Pop » d'OnBuch+ : fond crème (jamais blanc), bordures encre
% épaisses, ombres « dures » décalées, titres Archivo Black en capitales.
% Version MATHÉMATIQUES : graphes (pgfplots/tikz), boîtes pédagogiques
% (définition, propriété, méthode, attention, le savais-tu, exercice résolu,
% exercice, TP, bilan), page de garde et signature OnBuch+ ancrée en bas.
%
% Macros attendues dans main.tex AVANT \input{preamble} :
%   \DOCMATIERE  \DOCNIVEAU  \DOCMODULE  \DOCLECON (numéro)  \DOCTITRE
\documentclass[11pt]{article}

\usepackage{fontspec}
\usepackage[english]{babel}
\usepackage[a4paper,margin=2cm,top=3.4cm,bottom=2.8cm,headheight=1.4cm,footskip=1.3cm]{geometry}
\usepackage{amsmath,amssymb,amsfonts}
\usepackage{mathtools} % \xmapsto, \coloneqq... souvent utilisés par les modèles
\usepackage{cancel} % simplifications barrées (bilans de forces)
% \abs{z} (module, valeur absolue) et \norm{u} (norme), utilisés par les modèles
\providecommand{\abs}{}\let\abs\relax\DeclarePairedDelimiter{\abs}{\lvert}{\rvert}
\providecommand{\norm}{}\let\norm\relax\DeclarePairedDelimiter{\norm}{\lVert}{\rVert}
\usepackage[table]{xcolor} % [table] : fournit aussi \rowcolors (alternance de lignes)
\usepackage{pagecolor}
\usepackage{tikz}
\usetikzlibrary{shapes.symbols,circuits.logic.US,circuits.logic.IEC,arrows.meta,positioning,calc,shapes.geometric,shapes.misc,shapes.arrows,decorations.pathmorphing,decorations.pathreplacing,decorations.markings,patterns,shadows,mindmap,trees,fit,backgrounds,matrix,intersections,angles,quotes}
\usepackage{pgfplots}
\usepackage[version=4]{mhchem} % équations nucléaires \ce{^{235}_{92}U}
\usepackage[european,siunitx]{circuitikz} % schémas électriques
\pgfplotsset{compat=1.18}
\usepgfplotslibrary{fillbetween,groupplots,polar,statistics,dateplot} % aires entre courbes (calcul intégral)
\usepackage{enumitem}
\usepackage{fancyhdr}
% [most] charge toutes les bibliothèques tcolorbox (listings, minted, raster,
% documentation...) et gonfle inutilement la mémoire TeX sur un document long
% (~300 boîtes) : on ne charge que ce qui est réellement utilisé.
\usepackage[skins,breakable]{tcolorbox}
\usepackage{graphicx}
\usepackage{colortbl}
\usepackage{booktabs}
\usepackage{tabularx}
\usepackage{diagbox} % case d'en-tête diagonale des tableaux à double entrée
% Colonnes extensibles centrée / gauche / droite (C, L, R), souvent utilisées par les modèles
\newcolumntype{C}{>{\centering\arraybackslash}X}
\newcolumntype{L}{>{\raggedright\arraybackslash}X}
\newcolumntype{R}{>{\raggedleft\arraybackslash}X}
\usepackage{array}
\usepackage{multicol}
\usepackage{siunitx}
% parse-numbers=false : accepte aussi \SI{2,5 \times 10^{-3}}{...}
\sisetup{locale=UK,output-decimal-marker={.},group-minimum-digits=4,parse-numbers=false}
\DeclareSIUnit\ohm{\ensuremath{\symup{\Omega}}} % U+2126 absent de la police
\DeclareSIUnit\division{div} % réglages d'oscilloscope (ms/div, V/div)
\DeclareSIUnit\year{yr}\DeclareSIUnit\annum{yr}\DeclareSIUnit\Ma{Ma}\DeclareSIUnit\tr{tr} % tours (vitesse de rotation, ex. \SI{1500}{\tr\per\minute})
\usepackage{qrcode}
% \degree : commande fréquemment utilisée par les modèles pour un angle ou une
% température en dehors de \SI{}{} (non fournie par les paquets chargés).
\providecommand{\degree}{^{\circ}}
% \qed (fin de démonstration), utilisé par les modèles sans amsthm
\providecommand{\qed}{\hfill$\square$}
% \barre : double barre des valeurs interdites dans un tableau de variations (tkz-tab)
\providecommand{\barre}{\ensuremath{\|}}
% environnement proof (amsthm non chargé) utilisé par les modèles
\ifdefined\proof\else\newenvironment{proof}[1][Proof]{\par\noindent\textit{#1.}\ }{\qed\par}\fi
\usepackage{lastpage}
\usepackage{hyperref}
\hypersetup{hidelinks}

% ============================================================
% Palette « Pop » OnBuch+ — mêmes hex que la classe P (plus_ui.dart)
% ============================================================
\definecolor{poporange}{HTML}{F59321}
\definecolor{popdark}{HTML}{E07A0C}
\definecolor{popcream}{HTML}{FFF6E8}
\definecolor{popink}{HTML}{1C1714}
\definecolor{poppurple}{HTML}{7A5AE0}
\definecolor{popgreen}{HTML}{2E9E6A}
\definecolor{poppink}{HTML}{E0457B}
\definecolor{popblue}{HTML}{2F7DE1}
\definecolor{popgold}{HTML}{C98A12}
\definecolor{popmuted}{HTML}{8A7F76}
\definecolor{popline}{HTML}{EADFD2}
\definecolor{popgray}{HTML}{8A7F76}
\colorlet{popgrey}{popgray}
% Alias tolérants (le modèle nomme parfois les couleurs différemment)
\colorlet{popwhite}{white}
\colorlet{popblack}{popink}
\colorlet{popred}{poppink}
\colorlet{popyellow}{popgold}
% Teintes claires pour fonds de tableaux / zones
\colorlet{popblueL}{popblue!10!white}
\colorlet{poporangeL}{poporange!14!white}
\colorlet{popgreenL}{popgreen!12!white}
\colorlet{poppurpleL}{poppurple!12!white}
\colorlet{poppinkL}{poppink!12!white}
\colorlet{popgoldL}{popgold!14!white}

\pagecolor{popcream}

% ============================================================
% Typographie
% ============================================================
\defaultfontfeatures{Ligatures=TeX}
\setmainfont{PlusJakartaSans-Medium.ttf}[Path=../../../fonts/,BoldFont=PlusJakartaSans-Bold.ttf]
% Maths en sans-serif (Fira Math) avec chiffres et lettres droites de Plus
% Jakarta, pour que les formules se fondent dans le texte.
\usepackage[math-style=ISO,bold-style=ISO]{unicode-math}
\setmathfont{FiraMath-Regular.otf}[Path=../../../fonts/]
\setmathfont{PlusJakartaSans-Medium.ttf}[Path=../../../fonts/,range={up/num,up/{Latin,latin},\mathup}]
% (icomma retiré : décimales avec point en anglais)
% Caractères mathématiques Unicode absents des polices de texte (Plus Jakarta,
% Archivo Black) : rendus par leur équivalent mathématique, en texte comme en
% formule — sinon ils s'affichent en carré vide (ex. « Arithmétique dans ℤ »).
\usepackage{newunicodechar}
\newunicodechar{ℝ}{\ensuremath{\mathbb{R}}}
\newunicodechar{ℤ}{\ensuremath{\mathbb{Z}}}
\newunicodechar{ℂ}{\ensuremath{\mathbb{C}}}
\newunicodechar{ℕ}{\ensuremath{\mathbb{N}}}
\newunicodechar{ℚ}{\ensuremath{\mathbb{Q}}}
\newunicodechar{𝔻}{\ensuremath{\mathbb{D}}}
\newunicodechar{α}{\ensuremath{\alpha}}
\newunicodechar{β}{\ensuremath{\beta}}
\newunicodechar{γ}{\ensuremath{\gamma}}
\newunicodechar{δ}{\ensuremath{\delta}}
\newunicodechar{ε}{\ensuremath{\varepsilon}}
\newunicodechar{θ}{\ensuremath{\theta}}
\newunicodechar{λ}{\ensuremath{\lambda}}
\newunicodechar{μ}{\ensuremath{\mu}}
\newunicodechar{σ}{\ensuremath{\sigma}}
\newunicodechar{τ}{\ensuremath{\tau}}
\newunicodechar{φ}{\ensuremath{\varphi}}
\newunicodechar{ψ}{\ensuremath{\psi}}
\newunicodechar{ω}{\ensuremath{\omega}}
\newunicodechar{Σ}{\ensuremath{\Sigma}}
% majuscules produites par \MakeUppercase dans les titres (α-aminés -> Α)
\newunicodechar{Α}{\ensuremath{\alpha}}
\newunicodechar{Β}{\ensuremath{\beta}}
\newunicodechar{Γ}{\ensuremath{\Gamma}}
\newunicodechar{Δ}{\ensuremath{\Delta}}
\newunicodechar{Ε}{\ensuremath{\varepsilon}}
\newunicodechar{Θ}{\ensuremath{\Theta}}
\newunicodechar{Λ}{\ensuremath{\Lambda}}
\newunicodechar{Μ}{\ensuremath{\mu}}
\newunicodechar{Τ}{\ensuremath{\tau}}
\newunicodechar{Φ}{\ensuremath{\Phi}}
\newunicodechar{Ψ}{\ensuremath{\Psi}}
\newunicodechar{Ω}{\ensuremath{\Omega}}
\newunicodechar{↦}{\ensuremath{\mapsto}}
\newunicodechar{∈}{\ensuremath{\in}}
\newunicodechar{∉}{\ensuremath{\notin}}
\newunicodechar{∧}{\ensuremath{\wedge}}
\newunicodechar{⇔}{\ensuremath{\Leftrightarrow}}
\newunicodechar{⟺}{\ensuremath{\Longleftrightarrow}}
\newunicodechar{⇒}{\ensuremath{\Rightarrow}}
\newunicodechar{⟹}{\ensuremath{\Longrightarrow}}
\newunicodechar{∘}{\ensuremath{\circ}}
\newunicodechar{∪}{\ensuremath{\cup}}
\newunicodechar{∩}{\ensuremath{\cap}}
\newunicodechar{≡}{\ensuremath{\equiv}}
\newunicodechar{⊂}{\ensuremath{\subset}}
\newunicodechar{⊥}{\ensuremath{\perp}}
\newunicodechar{∃}{\ensuremath{\exists}}
\newunicodechar{∀}{\ensuremath{\forall}}
\newunicodechar{∛}{\ensuremath{\sqrt[3]{\,}}}
\newunicodechar{∜}{\ensuremath{\sqrt[4]{\,}}}
\newunicodechar{⃗}{\ensuremath{{}^{\to}}}
\newunicodechar{ⁿ}{\ifmmode^{n}\else\textsuperscript{\ensuremath{n}}\fi}
\newunicodechar{ˣ}{\ifmmode^{x}\else\textsuperscript{\ensuremath{x}}\fi}
\newunicodechar{ᵇ}{\ifmmode^{b}\else\textsuperscript{\ensuremath{b}}\fi}
\newunicodechar{ʳ}{\ifmmode^{r}\else\textsuperscript{\ensuremath{r}}\fi}
\newunicodechar{ᵘ}{\ifmmode^{u}\else\textsuperscript{\ensuremath{u}}\fi}
\newunicodechar{ʸ}{\ifmmode^{y}\else\textsuperscript{\ensuremath{y}}\fi}
\newunicodechar{ᵃ}{\ifmmode^{a}\else\textsuperscript{\ensuremath{a}}\fi}
\newunicodechar{ᵉ}{\ifmmode^{e}\else\textsuperscript{\ensuremath{e}}\fi}
\newunicodechar{ᵏ}{\ifmmode^{k}\else\textsuperscript{\ensuremath{k}}\fi}
\newunicodechar{ᵖ}{\ifmmode^{p}\else\textsuperscript{\ensuremath{p}}\fi}
\newunicodechar{ᴺ}{\ifmmode^{N}\else\textsuperscript{\ensuremath{N}}\fi}
\newunicodechar{ᶜ}{\ifmmode^{c}\else\textsuperscript{\ensuremath{c}}\fi}
\newunicodechar{ʰ}{\ifmmode^{h}\else\textsuperscript{\ensuremath{h}}\fi}
\newunicodechar{ᵠ}{\ifmmode^{q}\else\textsuperscript{\ensuremath{q}}\fi}
\newunicodechar{ₙ}{\ifmmode_{n}\else\textsubscript{\ensuremath{n}}\fi}
\newunicodechar{ₐ}{\ifmmode_{a}\else\textsubscript{\ensuremath{a}}\fi}
\newunicodechar{ᵢ}{\ifmmode_{i}\else\textsubscript{\ensuremath{i}}\fi}
\newunicodechar{ᵣ}{\ifmmode_{r}\else\textsubscript{\ensuremath{r}}\fi}
\newunicodechar{ₖ}{\ifmmode_{k}\else\textsubscript{\ensuremath{k}}\fi}
\newunicodechar{ᵦ}{\ifmmode_{\beta}\else\textsubscript{\ensuremath{\beta}}\fi}
\newunicodechar{ₓ}{\ifmmode_{x}\else\textsubscript{\ensuremath{x}}\fi}
\newfontfamily\popdisplay{ArchivoBlack-Regular.ttf}[Path=../../../fonts/]
\newfontfamily\poplogo{SpaceGrotesk-Bold.ttf}[Path=../../../fonts/]
\newfontfamily\popsemi{PlusJakartaSans-SemiBold.ttf}[Path=../../../fonts/]
\newfontfamily\popxbold{PlusJakartaSans-ExtraBold.ttf}[Path=../../../fonts/]
\newfontfamily\popxboldls{PlusJakartaSans-ExtraBold.ttf}[Path=../../../fonts/,LetterSpace=3.0]

\setlist[enumerate]{leftmargin=1.7em,itemsep=5pt,topsep=4pt}
\setlist[itemize]{leftmargin=1.7em,itemsep=4pt,topsep=4pt,label={\color{poporange}\textbullet}}
\setlength{\parindent}{0pt}
\setlength{\parskip}{5pt}
\renewcommand{\arraystretch}{1.25}

% pgfplots : style Pop par défaut
\pgfplotsset{
  every axis/.append style={axis line style={popink,line width=1pt},tick style={popink},
    grid style={popline,line width=0.5pt},label style={font=\small},tick label style={font=\footnotesize}},
  popcurve/.style={poporange,line width=1.8pt,no markers,smooth},
}
% Même style disponible dans un \draw TikZ simple (hors pgfplots)
\tikzset{popcurve/.style={poporange,line width=1.8pt}}
\tikzset{popgrid/.style={popline,very thin}}
\colorlet{popcurve}{poporange} % le nom du style est parfois utilisé comme couleur

% ============================================================
% Pill / squiggle
% ============================================================
\newcommand{\poppill}[3]{%
  \tikz[baseline=-0.55ex]\node[draw=popink,line width=1.2pt,rounded corners=2.6pt,%
    fill=#2,inner xsep=6.5pt,inner ysep=3pt]{%
    \popxboldls\fontsize{7}{7}\selectfont\color{#3}\MakeUppercase{#1}};%
}
\newcommand{\popsquiggle}[1]{%
  \tikz[baseline=-0.4ex]{
    \draw[#1,line width=2.6pt,line cap=round]
      plot [smooth] coordinates {(0,0.09) (0.34,0.01) (0.68,0.21) (1.02,0.01) (1.36,0.10)};
  }%
}
% Mot-clé surligné (marqueur orange sous le texte)
\newcommand{\cle}[1]{{\popxbold\color{popdark}#1}}

% ============================================================
% En-tête / pied
% ============================================================
\pagestyle{fancy}
\fancyhf{}
\renewcommand{\headrulewidth}{0pt}
\renewcommand{\footrulewidth}{0pt}
\lhead{\raisebox{-0.15ex}{\poplogo\fontsize{13}{13}\selectfont\color{popink}OnBuch\color{poporange}+}}
\rhead{\poppill{\DOCMATIERE\ \textbullet\ \DOCNIVEAU\ \textbullet\ Chapter \DOCLECON}{white}{popink}}
\lfoot{\popxbold\fontsize{7.6}{7.6}\selectfont\color{popmuted}\MakeUppercase{Signed course OnBuch+}\ \textbullet\ {\popsemi\DOCTITRE}}
\rfoot{\tikz[baseline=-0.6ex]\node[rounded corners=8.5pt,fill=popink,minimum height=17pt,minimum width=17pt,inner xsep=5pt,inner sep=0]{\popxbold\fontsize{8}{8}\selectfont\color{popcream}\thepage\,/\,\pageref*{LastPage}};}

% ============================================================
% Titres
% ============================================================
\newcommand{\chaptitre}[2]{%
  \begin{tcolorbox}[enhanced,colback=poporange,colframe=popink,boxrule=2.6pt,arc=11pt,
      drop shadow=popink,left=15pt,right=15pt,top=12pt,bottom=11pt,before skip=3pt,after skip=18pt]
    \poppill{#2}{popink}{popcream}\par\vspace{9pt}
    {\raggedright\hyphenpenalty=10000\exhyphenpenalty=10000\popdisplay\fontsize{22}{24}\selectfont\color{popcream}\MakeUppercase{#1}\par}
  \end{tcolorbox}
}
% Section numérotée : pastille orange avec le numéro + titre Archivo
\newcounter{popsec}
\newcommand{\coursec}[1]{%
  \stepcounter{popsec}\setcounter{popsub}{0}%
  \par\vspace{16pt}\noindent
  \begin{minipage}{\linewidth}
  \tikz[baseline=-0.7ex]\node[draw=popink,line width=1.6pt,fill=poporange,rounded corners=4pt,minimum size=22pt,inner sep=0]{\popdisplay\fontsize{12}{12}\selectfont\color{popcream}\thepopsec};%
  \hspace{8pt}{\popdisplay\fontsize{15}{17}\selectfont\color{popink}\MakeUppercase{#1}}\par
  \vspace{4pt}\noindent{\color{poporange}\rule{36pt}{3.6pt}}
  \end{minipage}\par\nopagebreak\vspace{8pt}
}
\newcounter{popsub}
\newcommand{\courssub}[1]{%
  \stepcounter{popsub}%
  \par\vspace{9pt}\noindent{\popxbold\fontsize{11.5}{13}\selectfont\color{popink}{\color{poporange}\thepopsec.\thepopsub}\hspace{6pt}#1}\par\nopagebreak\vspace{4pt}
}

% ============================================================
% Boîtes pédagogiques — même recette : fond blanc, bordure épaisse colorée,
% ombre dure encre, badge plein. Titre optionnel [#1].
% ============================================================
\tcbset{popbase/.style={breakable,enhanced,colback=white,boxrule=2.3pt,arc=9pt,
  shadow={3pt}{-3pt}{0pt}{popink},left=14pt,right=14pt,top=11pt,bottom=11pt,before skip=11pt,after skip=15pt,
  fontupper=\popsemi\fontsize{10.6}{14.5}\selectfont\color{popink},
  fontlower=\popsemi\fontsize{10.6}{14.5}\selectfont\color{popink}}}
% Coche dessinée (le glyphe ✓ n'existe pas dans les polices du document)
\newcommand{\popcheck}[1]{\tikz[baseline=-0.5ex]\draw[#1,line width=1.6pt,line cap=round,line join=round](-0.13,0.0)--(-0.04,-0.09)--(0.13,0.1);}
\providecommand{\coche}{\popcheck{popgreen}}
\providecommand{\resultat}[1]{\par\smallskip\noindent\fcolorbox{popgreen}{popgreenL}{\parbox{\dimexpr\linewidth-2\fboxsep-2\fboxrule\relax}{\textbf{Result.} #1}}\par\smallskip}
\newcommand{\popboxhead}[3]{\poppill{#1}{#2}{white}\ifx\relax#3\relax\else\hspace{6pt}{\popxbold\fontsize{10.6}{12}\selectfont\color{popink}#3}\fi\par\vspace{7pt}}

\newtcolorbox{aretenir}[1][]{popbase,colframe=popgreen,before upper={\popboxhead{Key points}{popgreen}{#1}}}
\newtcolorbox{exemplebox}[1][]{popbase,colframe=poporange,before upper={\popboxhead{Exemple}{poporange}{#1}}}
\newtcolorbox{definition}[1][]{popbase,colframe=popblue,before upper={\popboxhead{Definition}{popblue}{#1}}}
\newtcolorbox{propriete}[1][]{popbase,colframe=poppurple,before upper={\popboxhead{Property}{poppurple}{#1}}}
\newtcolorbox{methode}[1][]{popbase,colframe=popgold,before upper={\popboxhead{Method}{popgold}{#1}}}
\newtcolorbox{attention}[1][]{popbase,colframe=poppink,before upper={\popboxhead{Warning}{poppink}{#1}}}
\newtcolorbox{experience}[1][]{popbase,colframe=popblue,colback=popblueL,before upper={\popboxhead{Experiment}{popblue}{#1}}}
% « Le savais-tu ? » : carte violette pleine (contexte, culture, Cameroun)
\newtcolorbox{savaistu}[1][]{popbase,colback=poppurple,colframe=popink,
  fontupper=\popsemi\fontsize{10.4}{14.2}\selectfont\color{white},
  before upper={\poppill{Did you know?}{white}{poppurple}\ifx\relax#1\relax\else\hspace{6pt}{\popxbold\color{white}#1}\fi\par\vspace{7pt}}}
% Exercice résolu : énoncé en haut, \tcblower puis la solution sur fond crème
\newcounter{popexo}\newcounter{popexr}
\newtcolorbox{exoresolu}[1][]{popbase,colframe=poporange,colbacklower=popcream,
  segmentation style={popink,line width=1.2pt,dashed},
  before upper={\stepcounter{popexr}\popboxhead{Worked example \thepopexr}{poporange}{#1}},
  before lower={\poppill{Solution}{popink}{popcream}\par\vspace{6pt}}}
% Exercice (non résolu en place) — la correction est dans la partie Corrigés
\newtcolorbox{exercice}[1][]{popbase,colframe=popink,
  before upper={\stepcounter{popexo}\popboxhead{Exercise \thepopexo}{popink}{#1}}}
% Corrigé
\newtcolorbox{corrige}[1][]{popbase,colframe=popgreen,colback=popgreenL,
  before upper={\popboxhead{Solution}{popgreen}{#1}}}
% Objectifs / prérequis / bilan
\newtcolorbox{objectifs}{popbase,colframe=popink,colback=poporangeL,
  before upper={\popboxhead{Chapter objectives}{popink}{}}}
\newtcolorbox{prerequis}{popbase,colframe=popink,colback=popblueL,
  before upper={\popboxhead{Prerequisites}{popblue}{}}}
\newtcolorbox{bilan}{popbase,colframe=popink,colback=white,boxrule=2.8pt,
  before upper={\popboxhead{Fiche bilan}{poporange}{L'essentiel à connaître pour l'examen}}}
% Figure encadrée avec légende
\newtcolorbox{popfigure}[1][]{enhanced,breakable=false,colback=white,colframe=popline,boxrule=1.4pt,arc=8pt,
  left=10pt,right=10pt,top=10pt,bottom=8pt,before skip=10pt,after skip=12pt,halign=center,
  lower separated=false,halign lower=center,
  fontlower=\popsemi\fontsize{9}{11.5}\selectfont\color{popmuted},
  #1}
\newcommand{\legende}[1]{\par\vspace{6pt}{\popsemi\fontsize{9}{11.5}\selectfont\color{popmuted}#1}}

% ============================================================
% Signature OnBuch+
% ============================================================
\newcommand{\poplogoinline}[1]{{\poplogo\fontsize{#1}{#1}\selectfont\color{popink}OnBuch\color{poporange}+}}
\newcommand{\signatureonbuch}{%
  \begin{tcolorbox}[enhanced,colback=popink,colframe=popink,boxrule=2.2pt,arc=10pt,
      drop shadow=poporange,left=14pt,right=14pt,top=10pt,bottom=10pt,before skip=0pt,after skip=0pt]
    \tikz[baseline=-0.7ex]\node[circle,fill=popgreen,draw=popcream,line width=1.3pt,minimum size=22pt,inner sep=0]{\popcheck{white}};%
    \hspace{9pt}\begin{minipage}[c]{0.62\linewidth}
      {\popxbold\fontsize{12}{13}\selectfont\color{popcream}Signed\ \textbullet\ {\poplogo OnBuch\color{poporange}+}}\par\vspace{2pt}
      {\raggedright\popsemi\fontsize{8}{10.5}\selectfont\color{popline}\MakeUppercase{OnBuch+ teaching team}\\\MakeUppercase{Follows the official MINESEC syllabus}\par}
    \end{minipage}\hfill
    {\poplogo\fontsize{22}{22}\selectfont\color{popcream}OnBuch\color{poporange}+}
  \end{tcolorbox}%
}

% ============================================================
% Page de garde
% #1 titre · #2 matière · #3 niveau · #4 module · #5 n° leçon · #6 durée ·
% #7 description / objectifs (texte ou itemize)
% ============================================================
\newcommand{\pagegarde}[7]{%
  \thispagestyle{empty}
  \newgeometry{a4paper,margin=1.7cm,top=1.6cm,bottom=1.4cm}%
  \begin{tikzpicture}[remember picture,overlay]
    \fill[poporange!18] ([xshift=-3.2cm,yshift=-3.6cm]current page.north east) circle (4.6cm);
    \fill[poppurple!14] ([xshift=2.4cm,yshift=7.6cm]current page.south west) circle (3.8cm);
  \end{tikzpicture}%
  {\poplogo\fontsize{30}{30}\selectfont\color{popink}OnBuch\color{poporange}+}\hfill
  \raisebox{0.6ex}{\poppill{Signed course \textbullet\ Premium}{popink}{popcream}}\par
  \vspace{1pt}\popsquiggle{poppurple}\par
  \vspace{8pt}
  \begin{tcolorbox}[enhanced,colback=poporange,colframe=popink,boxrule=2.8pt,arc=13pt,
      drop shadow=popink,left=18pt,right=18pt,top=12pt,bottom=13pt,after skip=10pt]
    \poppill{#2 \textbullet\ #3}{popink}{popcream}\hspace{5pt}\poppill{#4}{white}{popink}\par\vspace{8pt}
    {\popdisplay\fontsize{14}{14}\selectfont\color{popink}CHAPTER #5}\par\vspace{4pt}
    {\raggedright\hyphenpenalty=10000\exhyphenpenalty=10000\popdisplay\fontsize{25}{27}\selectfont\color{popcream}\MakeUppercase{#1}\par}
    \vspace{8pt}
    {\popxbold\fontsize{9.5}{9.5}\selectfont\color{popink}\MakeUppercase{Suggested time: #6}}
  \end{tcolorbox}
  \begin{tcolorbox}[enhanced,colback=white,colframe=popink,boxrule=2.2pt,arc=10pt,
      drop shadow=popink,left=16pt,right=16pt,top=10pt,bottom=10pt,after skip=8pt]
    \poppill{In this chapter}{poppurple}{white}\par\vspace{5pt}
    \popsemi\fontsize{9.3}{12.6}\selectfont\color{popink}#7
  \end{tcolorbox}
  \noindent\begin{minipage}[t]{0.487\linewidth}
    \begin{tcolorbox}[enhanced,colback=popgreen,colframe=popink,boxrule=2.2pt,arc=10pt,
        drop shadow=popink,left=11pt,right=11pt,top=10pt,bottom=10pt,height=3.1cm,valign=center]
      \tcbox[colback=white,colframe=popink,boxrule=1.3pt,arc=5pt,boxsep=0pt,nobeforeafter,
        left=3pt,right=3pt,top=3pt,bottom=3pt,valign=center]{\qrcode[height=1.9cm]{https://play.google.com/store/apps/details?id=com.luvvix.onbuch&hl=en}}%
      \hspace{8pt}\begin{minipage}[c]{0.5\linewidth}
        {\popdisplay\fontsize{11}{12.5}\selectfont\color{white}\MakeUppercase{Download OnBuch}}\par\vspace{4pt}
        {\raggedright\popsemi\fontsize{8.4}{11}\selectfont\color{white}Free \textbullet\ Google Play\\AI tutor, past papers, results\par}
      \end{minipage}
    \end{tcolorbox}
  \end{minipage}\hfill
  \begin{minipage}[t]{0.487\linewidth}
    \begin{tcolorbox}[enhanced,colback=poppurple,colframe=popink,boxrule=2.2pt,arc=10pt,
        drop shadow=popink,left=11pt,right=11pt,top=10pt,bottom=10pt,height=3.1cm,valign=center]
      \tikz[baseline=-0.7ex]\node[circle,fill=white,draw=popink,line width=1.3pt,minimum size=1.6cm,inner sep=0]{\popdisplay\fontsize{24}{24}\selectfont\color{poppurple}+};%
      \hspace{8pt}\begin{minipage}[c]{0.6\linewidth}
        {\popdisplay\fontsize{11}{12.5}\selectfont\color{white}\MakeUppercase{Go OnBuch+}}\par\vspace{4pt}
        {\raggedright\popsemi\fontsize{8.4}{11}\selectfont\color{white}All signed courses, unlimited AI tutor, PDF sheets — OnBuch+ tab in the app\par}
      \end{minipage}
    \end{tcolorbox}
  \end{minipage}
  \vspace{8pt}
  \popcontenu
  \vfill
  \signatureonbuch
  \restoregeometry
  \newpage
}

% Rangée « Ce cours contient » (page de garde)
\newcommand{\poptile}[3]{% #1 couleur #2 gros texte #3 légende
  \begin{tcolorbox}[enhanced,colback=#1,colframe=popink,boxrule=2pt,arc=9pt,shadow={2.5pt}{-2.5pt}{0pt}{popink},
      left=6pt,right=6pt,top=9pt,bottom=9pt,halign=center,valign=center,height=1.75cm,width=\linewidth,nobeforeafter]
    {\popdisplay\fontsize{12.5}{13}\selectfont\color{white}\MakeUppercase{#2}}\par\vspace{3pt}
    {\centering\popsemi\fontsize{7.8}{9.5}\selectfont\color{white}#3\par}
  \end{tcolorbox}}
\newcommand{\popcontenu}{%
  \par\noindent{\popxboldls\fontsize{7.5}{7.5}\selectfont\color{popmuted}THIS SIGNED COURSE CONTAINS}\par\vspace{7pt}
  \noindent\begin{minipage}[t]{0.235\linewidth}\poptile{popblue}{Course}{complete, illustrated, step by step}\end{minipage}\hfill
  \begin{minipage}[t]{0.235\linewidth}\poptile{poporange}{Methods}{and worked examples}\end{minipage}\hfill
  \begin{minipage}[t]{0.235\linewidth}\poptile{poppink}{Exercises}{exam style + solutions}\end{minipage}\hfill
  \begin{minipage}[t]{0.235\linewidth}\poptile{popgreen}{Summary sheet}{the essentials to remember}\end{minipage}\par}

% Sommaire
\newtcolorbox{sommaire}{enhanced,colback=white,colframe=popink,boxrule=2.2pt,arc=10pt,
  drop shadow=popink,left=18pt,right=18pt,top=13pt,bottom=13pt,before skip=6pt,after skip=20pt,
  before upper={\poppill{Contents}{poppurple}{white}\par\vspace{9pt}\popsemi\fontsize{10.6}{14.5}\selectfont\color{popink}}}

% Signature de fin de cours — poussée tout en bas de la dernière page
\newcommand{\findecours}{%
  \par\vfill\nopagebreak
  \begin{center}\popsquiggle{poporange}\end{center}\vspace{4pt}
  \signatureonbuch
}
\AtBeginDocument{\def\llbracket{\mathopen{[\mkern-3.5mu[}}\def\rrbracket{\mathclose{]\mkern-3.5mu]}}} % intervalles d'entiers (glyphe ⟦ absent de la police)
% Glyphes absents de Fira Math : redéfinis à partir de glyphes disponibles
\AtBeginDocument{%
  \def\setminus{\mathbin{\backslash}}\let\smallsetminus\setminus
  \def\vdots{\mathord{\vbox{\baselineskip3.2pt\lineskiplimit0pt\kern4pt\hbox{.}\hbox{.}\hbox{.}}}}%
  \def\ddots{\mathinner{\mkern1mu\raise6.5pt\hbox{.}\mkern2mu\raise3.6pt\hbox{.}\mkern2mu\raise0.7pt\hbox{.}\mkern1mu}}%
  \def\triangle{\mathord{\scalebox{0.9}{$\symup{\Delta}$}}}%
  \def\nabla{\mathord{\raisebox{\depth}{\scalebox{1}[-1]{$\symup{\Delta}$}}}}%
  \def\checkmark{\popcheck{popdark}}%
  \def\star{\mathbin{\ast}}\def\bigstar{\mathord{\ast}}%
  \def\diamond{\mathbin{\scalebox{0.6}{$\Diamond$}}}%
  \def\bigcup{\mathop{\mathchoice{\raisebox{-0.3em}{\scalebox{1.7}{$\cup$}}}{\scalebox{1.25}{$\cup$}}{\cup}{\cup}}\displaylimits}%
  \def\bigcap{\mathop{\mathchoice{\raisebox{-0.3em}{\scalebox{1.7}{$\cap$}}}{\scalebox{1.25}{$\cap$}}{\cap}{\cap}}\displaylimits}%
  \def\lesseqqgtr{\mathrel{\lessgtr}}\def\bigoplus{\mathop{\oplus}\displaylimits}%
}
\providecommand{\ssi}{if and only if} % abréviation fréquente
\AtBeginDocument{\providecommand{\wideparen}{\overparen}} % arc de cercle

% Fonctions usuelles que les modèles utilisent sans que amsmath les fournisse
\providecommand{\arsinh}{\operatorname{arsinh}}\providecommand{\arcosh}{\operatorname{arcosh}}
\providecommand{\artanh}{\operatorname{artanh}}\providecommand{\arsech}{\operatorname{arsech}}
\providecommand{\arcsch}{\operatorname{arcsch}}\providecommand{\arcoth}{\operatorname{arcoth}}
\providecommand{\arcsinh}{\operatorname{arsinh}}\providecommand{\arccosh}{\operatorname{arcosh}}
\providecommand{\arctanh}{\operatorname{artanh}}\providecommand{\sech}{\operatorname{sech}}
\providecommand{\csch}{\operatorname{csch}}\providecommand{\arcsec}{\operatorname{arcsec}}
\providecommand{\arccsc}{\operatorname{arccsc}}\providecommand{\arccot}{\operatorname{arccot}}
\providecommand{\sgn}{\operatorname{sgn}}\providecommand{\lcm}{\operatorname{lcm}}
\providecommand{\Var}{\operatorname{Var}}\providecommand{\Cov}{\operatorname{Cov}}
\providecommand{\permil}{\ensuremath{{}^{0}\!/_{00}}}\providecommand{\textperthousand}{\ensuremath{{}^{0}\!/_{00}}}

%% FIGFIX-BEGIN v3 — correctifs globaux des figures (docs/onbuch-generation/tools/figfix)
\usepackage{adjustbox}\usepackage{etoolbox}
% toute figure TikZ/pgfplots plus large que la ligne est réduite à la largeur disponible
\BeforeBeginEnvironment{tikzpicture}{\begin{adjustbox}{max width=\linewidth}}
\AfterEndEnvironment{tikzpicture}{\end{adjustbox}}
% légendes pgfplots lisibles par-dessus les courbes
\pgfplotsset{every axis/.append style={legend style={fill=white,fill opacity=0.92,text opacity=1,draw=black!15,inner sep=3pt}}}
% étiquettes d'arêtes (syntaxe edge["..."]) avec fond blanc
\tikzset{every edge quotes/.append style={fill=white,inner sep=1.2pt}}
% bulles de cartes mentales : texte à la ligne dans la bulle, bulle un peu plus grande
\tikzset{every concept/.append style={text width=2.0cm,align=center,minimum size=2.3cm,inner sep=2pt}}
%% FIGFIX-END
\begin{document}

\pagegarde{Learn with the assistance of computers}{ICT}{Upper Sixth}{Upper Sixth — ICT}{14}{3 periods}{This chapter explores the main electronic services — e‑commerce, e‑health, e‑government — their applications, benefits and challenges, and then delves into Computer‑Assisted Learning (CAL) and e‑learning models, tools and evaluation methods. It equips learners to analyse, design and critically assess ICT‑driven services in the Cameroonian context.
\vspace{4pt}
\begin{multicols}{2}\small
\begin{itemize}
\item Identify and describe the three major categories of electronic services (e‑commerce, e‑health, e‑government) with Cameroonian examples.
\item Explain the technical architecture and key stakeholders of each electronic service.
\item Analyse the benefits (efficiency, accessibility, transparency) and challenges (infrastructure, security, digital divide) of electronic services.
\item Define Computer‑Assisted Learning (CAL) and distinguish it from traditional instruction.
\item Compare major e‑learning models (synchronous, asynchronous, blended, MOOC) and their suitability for Cameroonian schools.
\item Design a simple CAL/e‑learning scenario integrating multimedia, assessment and learner tracking.
\end{itemize}\end{multicols}}

\begin{sommaire}
\begin{multicols}{2}
\begin{enumerate}
\item Electronic Services Landscape: E‑Commerce, E‑Health, E‑Government
\item Applications, Benefits and Challenges of Electronic Services
\item Computer‑Assisted Learning (CAL): Concepts, Tools and Design
\item E‑Learning Models, Implementation and Evaluation in the Cameroonian Context
\item Integration activity
\item Methods and worked examples
\item Exercises
\item Solutions to the exercises
\item Summary sheet
\end{enumerate}
\end{multicols}
\end{sommaire}

\begin{objectifs}
\begin{itemize}
\item Identify and describe the three major categories of electronic services (e‑commerce, e‑health, e‑government) with Cameroonian examples.
\item Explain the technical architecture and key stakeholders of each electronic service.
\item Analyse the benefits (efficiency, accessibility, transparency) and challenges (infrastructure, security, digital divide) of electronic services.
\item Define Computer‑Assisted Learning (CAL) and distinguish it from traditional instruction.
\item Compare major e‑learning models (synchronous, asynchronous, blended, MOOC) and their suitability for Cameroonian schools.
\item Design a simple CAL/e‑learning scenario integrating multimedia, assessment and learner tracking.
\item Evaluate the effectiveness of an e‑learning intervention using Kirkpatrick’s four‑level model.
\item Discuss legal, ethical and policy issues (data protection, intellectual property) relevant to electronic services in Cameroon.
\end{itemize}
\end{objectifs}

\begin{prerequis}
\begin{itemize}
\item Basic Internet concepts (URL, HTTP, browser) from Lower Sixth ICT.
\item Fundamentals of database management (tables, queries) covered in ICT13.
\item Understanding of computer networks (LAN, WAN, mobile networks) from ICT11.
\item General knowledge of Cameroonian administrative structure (ministries, local councils).
\item Familiarity with common office software (word processor, spreadsheet) for content creation.
\end{itemize}
\end{prerequis}

\newpage

\coursec{Electronic Services Landscape: E‑Commerce, E‑Health, E‑Government}

In a bustling market in Douala, a vendor accepts payment via Mobile Money while a nurse in a rural health centre consults a specialist through a telemedicine app. Meanwhile, a student in Buea logs into an e‑learning platform to revise for the GCE exams, and a citizen files a tax return on the government portal. These everyday scenes illustrate how electronic services and computer‑assisted learning are reshaping Cameroonian life. But what exactly makes a service ``electronic''? How do these systems actually work behind the screen, who builds and regulates them, and why do some succeed while others fail? This section answers these questions by mapping the landscape of the three great families of electronic services: \cle{e‑commerce}, \cle{e‑health} and \cle{e‑government}.

\courssub{What is an Electronic Service?}

\textbf{Intuition.}\par
Think of buying a textbook. Traditionally, you travel to a bookshop in Yaoundé, queue, pay cash and carry the book home. Now imagine doing the same thing from your phone in Bamenda: you browse a catalogue, pay with MTN Mobile Money, and the book is dispatched. The \emph{service} (selling a book) has not changed; what has changed is the \emph{channel} through which it is delivered — an electronic network instead of a physical counter. That is the heart of the notion.

\begin{definition}[Electronic Service]
An \cle{electronic service} (e‑service) is a service delivered over electronic networks, typically the Internet or mobile networks, in which the interaction between the provider and the user is carried out through digital devices (computers, smartphones, USSD handsets) rather than exclusively through face‑to‑face contact.
\end{definition}

The \emph{scope} of electronic services is very broad. It covers commercial transactions (e‑commerce), health care delivery (e‑health), public administration (e‑government), education (e‑learning, studied later in this chapter), banking (e‑banking) and many others. Three features characterise any e‑service:

\begin{itemize}
\item \textbf{A digital channel:} a website, a mobile application, an SMS/USSD menu or a call centre backed by software.
\item \textbf{A back‑end information system:} databases and application servers that process requests automatically.
\item \textbf{A physical or legal back‑office:} goods still have to be delivered, doctors still diagnose, and civil status documents still have legal value. The electronic layer \emph{mediates} the service; it rarely replaces the whole of it.
\end{itemize}

\begin{attention}[E‑commerce is not e‑marketing]
A very common mistake in examination answers is to confuse \cle{e‑commerce} with \cle{e‑marketing}. E‑commerce is \emph{transactional}: it covers the actual buying, selling and paying online. E‑marketing is \emph{promotional}: it covers advertising, social‑media campaigns and customer‑relationship techniques used to attract buyers. A Facebook page advertising dresses is e‑marketing; it only becomes e‑commerce when the customer can order and pay through an electronic channel. When you describe the \emph{benefits} of e‑commerce, do not list purely promotional advantages (``the shop becomes visible'') as if they were transactional ones.
\end{attention}

\begin{savaistu}[Cameroon's digital leap]
Cameroon's \emph{Stratégie Nationale de Développement du Numérique (SNDN 2020-2030)} aims to make the country a digital hub in Central Africa. Key pillars include: extending broadband coverage (the \emph{Backbone} fibre‑optic network linking major cities), promoting e‑government services (the \emph{e‑Gouv} platform), securing cyberspace via ANTIC, and fostering digital entrepreneurship. Mobile money, launched in 2011, now processes billions of CFA francs monthly — far exceeding traditional bank transactions in volume — and is the primary enabler of e‑commerce for the unbanked majority.
\end{savaistu}

\courssub{E‑Commerce: Models and the Cameroonian Reality}

\textbf{Intuition.}\par
Not all online trade looks the same. When Camrail buys spare parts from a supplier through a private online procurement system, when you order a phone on Jumia, and when a student resells her used scientific calculator through a Facebook group, three different \emph{business models} are at work. Classifying them is a classic examination question.

\begin{definition}[E‑commerce models]
\begin{itemize}
\item \cle{B2B} (Business‑to‑Business): electronic transactions between two businesses, e.g.\ a wholesaler supplying retailers through an online ordering portal.
\item \cle{B2C} (Business‑to‑Consumer): a business sells directly to individual customers online, e.g.\ Jumia Cameroon selling a blender to a student.
\item \cle{C2C} (Consumer‑to‑Consumer): individuals trade with each other, usually through a platform that connects them, e.g.\ second‑hand goods exchanged via Facebook Marketplace or classified‑ads groups.
\end{itemize}
(One also meets B2G — business‑to‑government, e.g.\ electronic public procurement — and C2G, e.g.\ paying taxes online, which overlaps with e‑government.)
\end{definition}

\textbf{Cameroonian platforms and payment gateways.}\par
In Cameroon, the best‑known B2C platform has been \textbf{Jumia CM}, offering electronics, fashion and household goods with payment on delivery or by mobile money. \textbf{Kaymu} operated earlier as a C2C/B2C marketplace before its model was absorbed into Jumia's marketplace. A very large share of real Cameroonian e‑commerce, however, is \emph{informal}: thousands of vendors run shops through \textbf{Facebook pages and Marketplace}, WhatsApp Business and Instagram, agreeing on price by chat and delivering by motorbike or travel agency (``envoyer par car''). Payment is dominated not by bank cards — card penetration remains low — but by \cle{mobile money} gateways: \textbf{MTN Mobile Money} and \textbf{Orange Money}, which allow any subscriber with a basic phone to send and receive funds identified by a telephone number and a PIN code. Emerging fintech solutions (e.g.\ \emph{Djangui}, \emph{Maviance}) now offer aggregated payment APIs for merchants.

\begin{exemplebox}[Classifying transactions]
State the model for each case:
\begin{enumerate}
\item A Douala bakery orders 50 bags of flour from a miller's online portal. $\rightarrow$ B2B.
\item A student buys earphones on Jumia CM and pays with Orange Money. $\rightarrow$ B2C.
\item A civil servant sells his used generator to a neighbour found on Facebook Marketplace. $\rightarrow$ C2C.
\end{enumerate}
\end{exemplebox}

\courssub{E‑Health: Telemedicine, EHR and mHealth}

\textbf{Intuition.}\par
Cameroon has a severe shortage of specialist doctors, and most of them are concentrated in Yaoundé and Douala. A patient in a remote health area cannot easily travel 300 km for a cardiology opinion. E‑health uses ICT to move the \emph{information} instead of moving the \emph{patient}.

\begin{definition}[E‑health]
\cle{E‑health} is the use of information and communication technologies to support health services, health surveillance, health education and the management of health information.
\end{definition}

Its main components are:

\begin{itemize}
\item \textbf{\cle{Telemedicine}:} remote delivery of clinical care — a nurse in a rural health centre transmits a patient's symptoms, photos or ECG to a specialist in a reference hospital, who replies with a diagnosis and prescription (teleconsultation, telediagnosis, telemonitoring).
\item \textbf{\cle{Electronic Health Records} (EHR):} the patient's medical history (consultations, prescriptions, lab results, vaccinations) stored in a structured database rather than in a paper booklet, so that any authorised facility can retrieve it. Standards such as HL7/FHIR enable interoperability between systems.
\item \textbf{\cle{mHealth} (mobile health):} health services delivered through mobile phones — SMS reminders for vaccination or antenatal appointments, apps for tracking epidemics (e.g.\ cholera, COVID‑19), USSD menus for health information (e.g.\ \#123\# for Ministry of Health alerts).
\end{itemize}

\begin{experience}[Case study: the ``e‑Santé'' pilot in the North‑West region]
In the ``e‑Santé'' pilot initiative conducted in the North‑West region, selected rural health centres were equipped with computers, Internet connectivity and a shared patient‑record application linked to a referral hospital. Health workers could register patients, enter consultation data and request remote advice from doctors at the referral level. Observed effects included faster referral decisions, better follow‑up of chronic patients and reduced unnecessary travel. The pilot also revealed typical difficulties: unstable electricity (solved partly with solar panels and generators), intermittent connectivity, the need to train health workers in basic computer skills, and the obligation to protect the confidentiality of medical data. The lesson: technology alone is not an e‑health service — it must be embedded in training, infrastructure and data‑protection rules.
\end{experience}

\courssub{E‑Government: Bringing the Administration Online}

\textbf{Intuition.}\par
Anyone who has queued for hours to obtain a document understands the promise of e‑government: let the citizen submit the request online, let the administration process it digitally, and let the citizen come only once — or not at all.

\begin{definition}[E‑government]
\cle{E‑government} is the use of ICT, especially the Internet, by public administrations to deliver information and services to citizens (G2C), businesses (G2B) and other administrations (G2G), and to improve internal efficiency and transparency.
\end{definition}

Concrete Cameroonian examples include:

\begin{itemize}
\item \textbf{Online tax filing:} the Directorate General of Taxation (DGI) provides an online portal through which registered taxpayers can declare and pay certain taxes electronically, reducing queues at tax centres.
\item \textbf{Business registration:} the \textbf{CFCE} (Centre de Formalités de Création d'Entreprises) one‑stop shops, supported by computerised procedures, allow an entrepreneur to create a company through a single window in a few days instead of weeks.
\item \textbf{Civil status requests:} computerised procedures for requesting birth certificates and other civil status documents, progressively being digitised in councils.
\item \textbf{e‑Procurement:} the \emph{Marchés Publics En Ligne} platform for transparent public tendering.
\end{itemize}

The \textbf{National Agency for Information and Communication Technologies (\cle{ANTIC})} plays a central role: it promotes and monitors the Government's ICT policy, secures State information systems, manages the national public key infrastructure and electronic certification, and supports the interconnection of administrations. The \textbf{Telecommunications Regulatory Board (\cle{ART})} regulates the telecom operators and tariffs on which all e‑services depend.

\begin{propriete}[Interoperability]
\cle{Interoperability} is the ability of different e‑government (and e‑service) systems to exchange data and work together through \emph{open standards} such as XML and JSON for data formats and SOAP/REST for web‑service interfaces (APIs). Without interoperability, each ministry builds an isolated ``silos'' system and the citizen must supply the same information again and again. (The statement of this property, with examples of standards, is examinable; no formal proof is required.)
\end{propriete}

\courssub{Technical Architecture of an Electronic Service}

\textbf{Intuition.}\par
Whether you consult Jumia, a telemedicine app or the DGI portal, the same three questions arise behind the screen: \emph{How is the information presented to me? Where is the request processed? Where is the data stored?} This leads to the universal \cle{three‑layer architecture}.

\begin{definition}[Three‑layer architecture]
\begin{itemize}
\item \textbf{Presentation layer:} the user interface — web pages served by a \emph{web server}, a mobile app, or a USSD/SMS menu for basic phones.
\item \textbf{Logic (application) layer:} the \emph{application server} runs the business rules (compute a tax amount, check stock, validate a prescription) and exposes \emph{APIs} so that other systems (e.g.\ a mobile money gateway) can connect.
\item \textbf{Data layer:} the \emph{database management system} stores products, patients, taxpayers, payments.
\end{itemize}
\end{definition}

The USSD/SMS channel is crucial in Cameroon: it allows users with ordinary feature phones and no Internet bundle to access e‑services (checking a balance, paying a bill, receiving an appointment reminder) through short codes such as \#126\#.

\begin{popfigure}
\begin{center}
\begin{tikzpicture}[font=\small,
layer/.style={draw=popink, fill=popblueL, rounded corners, minimum width=10.5cm, minimum height=1.5cm, align=center},
box/.style={draw=popdark, fill=white, rounded corners, minimum width=3.1cm, minimum height=0.95cm, align=center},
gw/.style={draw=poporange, fill=popgreenL, rounded corners, minimum width=4.6cm, minimum height=0.95cm, align=center},
lbl/.style={font=\bfseries\small, text=popdark}]
% Layers
\node[layer] (pres) at (0,4.6) {};
\node[layer] (log) at (0,2.3) {};
\node[layer] (dat) at (0,0) {};
\node[lbl] at (-4.7,4.6) {Presentation};
\node[lbl] at (-4.7,2.3) {Logic};
\node[lbl] at (-4.7,0) {Data};
% Presentation boxes
\node[box, align=center] at (-2.6,4.6){Web server\\(Jumia, DGI portal)};
\node[box, align=center] at (0.9,4.6){Mobile app\\(mHealth, e-learning)};
\node[box, align=center] at (4.4,4.6){USSD / SMS\\menu (\#126\#)};
% Logic boxes
\node[box, align=center] at (-2.6,2.3){Application server\\e-commerce rules};
\node[box, align=center] at (0.9,2.3){Application server\\e-health / e-gov};
\node[box, align=center] at (4.4,2.3){APIs (REST/SOAP)\\XML, JSON};
% Data boxes
\node[box, align=center] at (-2.6,0){Products \&\\orders DB};
\node[box, align=center] at (0.9,0){EHR / taxpayers\\DB};
\node[box, align=center] at (4.4,0){Payments\\log DB};
% Gateways
\node[gw, align=center] (mm) at (-2.3,-2.0){Mobile Money gateway\\(MTN MoMo, Orange Money)};
\node[gw, align=center] (ussd) at (3.1,-2.0){USSD / SMS gateway\\(CAMTEL, MTN, Orange)};
% Arrows vertical
\draw[<->, thick, popblue] (-2.6,4.1) -- (-2.6,2.85);
\draw[<->, thick, popblue] (0.9,4.1) -- (0.9,2.85);
\draw[<->, thick, popblue] (4.4,4.1) -- (4.4,2.85);
\draw[<->, thick, popblue] (-2.6,1.8) -- (-2.6,0.55);
\draw[<->, thick, popblue] (0.9,1.8) -- (0.9,0.55);
\draw[<->, thick, popblue] (4.4,1.8) -- (4.4,0.55);
% Arrows to gateways
\draw[<->, thick, poporange] (-2.6,1.8) .. controls (-3.4,0.9) and (-3.2,-0.9) .. (-2.3,-1.5);
\draw[<->, thick, poporange] (4.4,4.05) .. controls (5.6,2.6) and (5.0,-0.6) .. (3.6,-1.5);
\draw[<->, thick, poporange] (0.9,1.8) .. controls (1.8,0.8) and (2.6,-0.8) .. (3.0,-1.5);
\end{tikzpicture}
\end{center}
\legende{Three‑layer architecture of Cameroonian e‑services: presentation (web, app, USSD), logic (application servers and APIs) and data (databases), connected to mobile money and USSD/SMS gateways.}
\end{popfigure}

\courssub{Stakeholder Map}

An e‑service succeeds only if many actors cooperate. The main stakeholders are:

\begin{center}
\begin{tabularx}{\linewidth}{|l|X|}
\hline
\rowcolor{popblueL} \textbf{Stakeholder} & \textbf{Role} \\
\hline
Government ministries (Public Health, Finance, MINPOSTEL, MINESEC) & Define policy, fund and operate e‑health, e‑tax and e‑learning services. \\
\hline
ANTIC & Secures State systems, electronic certification, coordinates e‑government. \\
\hline
ART & Regulates telecom operators, quality of service and tariffs. \\
\hline
Private operators (MTN, Orange, CAMTEL, banks, startups) & Provide networks, mobile money, platforms and hosting. \\
\hline
Civil society and consumer associations & Defend users' rights, digital inclusion, data protection awareness. \\
\hline
End‑users (citizens, patients, students, traders) & Adopt (or reject) the services; their trust determines success. \\
\hline
\end{tabularx}
\end{center}

\courssub{Worked Example and Method}

\begin{exoresolu}[Analysing an e‑commerce scenario]
\emph{Statement.} A Douala shoe seller creates a WhatsApp Business catalogue, posts prices on Facebook, receives orders by chat, is paid by MTN Mobile Money and delivers through a travel agency. (a) Identify the e‑commerce model(s). (b) Identify the presentation, logic and data elements. (c) Give one benefit and one risk.\\[4pt]
\tcblower
\emph{Solution.} (a) The seller is a business selling to individuals: \textbf{B2C} (informal/social commerce). If she later resells a personal phone on Marketplace, that would be C2C. (b) \textbf{Presentation:} WhatsApp catalogue and Facebook page; \textbf{logic:} the seller's own order handling plus the MTN MoMo payment gateway (there is no dedicated application server — a weakness); \textbf{data:} the chat history and MoMo transaction records act as an informal database. (c) \textbf{Benefit:} very low start‑up cost and nationwide reach. \textbf{Risk:} no formal contract or receipt, so disputes over non‑delivery are hard to resolve, and customer data is not protected.
\end{exoresolu}

\begin{methode}[Conducting a SWOT analysis of a Cameroonian e‑commerce platform]
\begin{enumerate}
\item \textbf{Choose and describe} the platform precisely (e.g.\ Jumia CM or a Facebook‑based shop): model (B2C/C2C), products, payment and delivery channels.
\item \textbf{Strengths (internal):} list what the platform does well — brand, logistics, mobile‑money integration, prices.
\item \textbf{Weaknesses (internal):} high delivery costs, limited product range, poor after‑sales service, dependence on one operator.
\item \textbf{Opportunities (external):} growth of smartphone use, expanding 4G coverage, youth entrepreneurship, e‑government digitisation building trust.
\item \textbf{Threats (external):} poor road infrastructure, low purchasing power, fraud and distrust of online payment, competition from informal sellers.
\item \textbf{Conclude} with a strategy: e.g.\ use a Strength to seize an Opportunity (leverage MoMo integration to reach unbanked customers).
\end{enumerate}
\end{methode}

\begin{attention}[Describing benefits rigorously]
When asked for the \emph{benefits of e‑commerce}, give transactional and organisational benefits: 24/7 availability, wider geographic market, lower transaction costs, electronic payment and record‑keeping. Do not present a promotional poster (e‑marketing) as an e‑commerce service, and do not claim a platform ``eliminates all middlemen'' — delivery, payment gateways and platforms are themselves intermediaries.
\end{attention}

\begin{aretenir}[Key points]
\begin{itemize}
\item An \cle{electronic service} is delivered over electronic networks (Internet or mobile); the physical back‑office usually remains.
\item E‑commerce models: \cle{B2B}, \cle{B2C}, \cle{C2C}; in Cameroon, Jumia CM, Facebook/WhatsApp commerce, paid via MTN Mobile Money and Orange Money.
\item \cle{E‑health} = telemedicine + electronic health records + mHealth, as illustrated by the North‑West ``e‑Santé'' pilot.
\item \cle{E‑government} (DGI online tax filing, CFCE business registration, civil status) is coordinated and secured by \cle{ANTIC}; telecoms are regulated by \cle{ART}.
\item All e‑services share a \cle{three‑layer architecture} (presentation, logic, data) linked to payment and USSD/SMS gateways; \cle{interoperability} via XML/JSON and SOAP/REST APIs is essential.
\end{itemize}
\end{aretenir}

\begin{exercice}[Check your understanding]
\begin{enumerate}
\item Define an electronic service and give two Cameroonian examples not mentioned in this course.
\item Classify: (a) a brewery ordering bottles from a manufacturer's portal; (b) a student buying data bundles on an operator's app; (c) two farmers exchanging a goat for a bicycle via a Facebook group.
\item Explain the difference between telemedicine and mHealth, with one example each.
\item Why is interoperability important for e‑government? Name two open standards that enable it.
\item Draw and label the three‑layer architecture of the DGI online tax portal, including the mobile money gateway.
\item A rural health centre wants to implement a teleconsultation service. List the minimum hardware, software and connectivity requirements, and name one Cameroonian regulation that applies to patient data.
\end{enumerate}
\end{exercice}

\begin{corrige}[Exercise n°1]
(1) A service delivered over electronic networks; examples: e‑banking (Afriland mobile app), online university registration, CAMTEL bill payment by USSD. (2) (a) B2B; (b) B2C; (c) C2C. (3) Telemedicine = remote clinical care between health professionals/patient (a nurse sending an ECG to a cardiologist); mHealth = health services via mobile phones (SMS vaccination reminders). (4) It allows ministries' systems to exchange data so citizens supply information once; standards: XML/JSON, SOAP/REST. (5) Presentation: DGI web portal and USSD; Logic: application server computing tax, API to MoMo gateway; Data: taxpayer and payments databases. (6) Hardware: computer/tablet, webcam, scanner, printer, UPS/solar backup. Software: teleconsultation app (video, chat, file sharing), EHR module. Connectivity: stable Internet (4G/fibre) or at least 3G for low‑bandwidth mode. Regulation: Law No. 2010/012 on cybersecurity and cybercrime (data protection provisions) and ANTIC guidelines on health data confidentiality.
\end{corrige}

\coursec{Applications, Benefits and Challenges of Electronic Services}

In the previous section we explored the landscape of electronic services — e‑commerce, e‑health and e‑government — and saw how they are reshaping interactions between citizens, businesses and the state. We now turn to a deeper analysis: why do governments and organisations invest in these services, what tangible benefits do they bring, and what obstacles must be overcome for them to thrive in Cameroon? Understanding this balance is essential for any ICT professional who will design, implement or evaluate digital services.

\courssub{Key Benefits of Electronic Services}

Electronic services are not merely a technological upgrade; they transform the economics and experience of service delivery. The benefits can be grouped into five pillars.

\textbf{1. 24/7 Availability and Convenience.} Unlike physical offices with fixed opening hours, a well‑designed e‑service platform operates continuously. A trader in Douala can file tax returns at midnight; a patient in a remote village can book a teleconsultation on a Sunday. This \emph{time‑shift} capability is especially valuable in a country where travel to administrative centres can take hours.

\textbf{2. Reduced Transaction Costs.} Moving processes online eliminates paper, printing, physical storage and manual data entry. For the administration, the cost per transaction drops dramatically. For the user, indirect costs — transport, lost wages while queuing, informal “facilitation” payments — are slashed. Studies on African e‑government initiatives estimate that digital tax filing reduces the compliance cost for small businesses by 30–50\%.

\textbf{3. Broader Market Reach (E‑Commerce).} An artisan in Foumban selling traditional crafts through Jumia or a dedicated website instantly accesses buyers in Yaoundé, Douala, and even abroad. The geographic barrier dissolves, enabling micro‑enterprises to scale without a physical storefront in every city.

\textbf{4. Improved Health Outcomes (E‑Health).} Telemedicine platforms, electronic medical records (EMR) and mobile health (mHealth) apps enable early diagnosis, chronic disease monitoring and specialist referrals without the patient travelling hundreds of kilometres. Pilot mHealth projects for maternal health in the East Region have shown reductions in antenatal appointment no‑shows through SMS reminders.

\textbf{5. Transparency and Anti‑Corruption.} When processes are digitised, every step leaves an audit trail. E‑procurement systems, for instance, publish tenders, evaluation criteria and award decisions online, making it harder for officials to manipulate bids. The \emph{Open Contracting Data Standard} adopted by some Cameroonian ministries allows citizens and watchdogs to track public spending in real time.

\begin{aretenir}[Benefits at a Glance]
\begin{itemize}
\item \textbf{Availability:} Services accessible anytime, anywhere with an Internet connection.
\item \textbf{Cost reduction:} Lower operational costs for providers; lower access costs for users.
\item \textbf{Reach:} National and international markets opened to local producers.
\item \textbf{Health impact:} Faster diagnosis, better follow‑up, reduced mortality in pilot zones.
\item \textbf{Governance:} Audit trails, open data and reduced discretionary power curb corruption.
\end{itemize}
\end{aretenir}

\courssub{Infrastructure Challenges: The Foundation That Must Hold}

No e‑service can function without reliable infrastructure. Cameroon faces three intertwined gaps.

\begin{popfigure}
\centering
\begin{tikzpicture}
\begin{axis}[
    ybar,
    bar width=0.35cm,
    width=\linewidth,
    height=10cm,
    enlargelimits=0.15,
    legend style={at={(0.5,-0.25)},anchor=north,legend columns=-1,font=\small},
    ylabel={Percentage (\% of population)},
    symbolic x coords={Adamawa, Centre, East, Far North, Littoral, North, Northwest, West, South, Southwest},
    xtick=data,
    x tick label style={rotate=45,anchor=east,font=\small},
    ymin=0,ymax=100,
    ymajorgrids=true,
    grid style=dashed,
]
\addplot[fill=popblue] coordinates {
    (Adamawa,28) (Centre,45) (East,22) (Far North,18) (Littoral,58)
    (North,20) (Northwest,35) (West,40) (South,30) (Southwest,38)
};
\addplot[fill=poporange] coordinates {
    (Adamawa,12) (Centre,28) (East,8) (Far North,5) (Littoral,35)
    (North,7) (Northwest,15) (West,18) (South,10) (Southwest,20)
};
\legend{Internet penetration (est.), E‑service adoption (est.)}
\end{axis}
\end{tikzpicture}
\legende{Illustrative comparison of estimated Internet penetration vs. estimated e‑service adoption across Cameroon's ten regions (based on 2022–2023 reports). Note the adoption gap even in relatively connected regions.}
\end{popfigure}

\textbf{Broadband Penetration.} While mobile Internet subscriptions exceed 15 million (2023 estimates), fixed broadband remains below 2\% of households. 4G coverage is concentrated in urban corridors; 3G/Edge dominates rural areas, limiting bandwidth for video consultations or large file transfers in e‑procurement.

\textbf{Electricity Reliability.} Frequent outages and voltage fluctuations damage equipment and interrupt services. Many telecentres and schools rely on generators, increasing operational costs. The \emph{Rural Electrification Master Plan} aims to raise rural access from approximately 25\% to 60\% by 2030, but progress is gradual.

\textbf{Last‑Mile Connectivity.} The “last mile” — the final link to the user — is often the most expensive. In the Far North and East regions, sparse population density makes fibre deployment economically unattractive for private operators. Universal Service Fund projects have piloted VSAT and TV white‑space solutions, but sustainability remains a challenge.

\begin{attention}[Infrastructure \textbackslash= Usage]
High mobile penetration (SIM cards per capita) is often mistaken for high e‑service readiness. A farmer may own a smartphone but lack the digital literacy to navigate an e‑tax portal, or distrust the platform due to past fraud. Infrastructure is necessary but not sufficient; human factors (literacy, trust, relevance) are equally critical.
\end{attention}

\courssub{Security, Privacy and the Legal Framework}

As e‑services handle sensitive personal, medical and financial data, security is not optional — it is a legal requirement and a trust prerequisite.

\begin{definition}[Digital Divide]
The \cle{digital divide} is the gap between individuals, households, businesses and geographic areas at different socio‑economic levels with regard to their opportunities to access information and communication technologies (ICTs) and to their use of the Internet for a wide variety of activities.
\end{definition}

\begin{propriete}[Law No. 2010/012 on Cybersecurity and Cybercrime]
\textbf{Key provisions relevant to e‑services:}
\begin{itemize}
\item \textbf{Article 5:} Criminalises unauthorised access to computer systems, data interception and system interference.
\item \textbf{Article 12:} Obliges service providers to implement technical and organisational measures to secure personal data.
\item \textbf{Article 23:} Establishes the \cle{National Agency for Information and Communication Technologies (ANTIC)} as the national Computer Emergency Response Team (CERT) and the authority for cryptography and electronic certification.
\item \textbf{Article 31:} Requires critical infrastructure operators (including e‑health and e‑government platforms) to report security incidents to ANTIC within 72 hours.
\end{itemize}
Compliance is mandatory for all e‑service providers operating in Cameroon.
\end{propriete}

\textbf{Threat Landscape.} Cameroon records rising cases of phishing (fake MTN Mobile Money or Orange Money pages), ransomware targeting hospital EMR systems, and business email compromise (BEC) scams against companies using e‑procurement. The lack of a widely deployed \cle{Public Key Infrastructure (PKI)} means most authentication still relies on passwords, which are easily phished.

\textbf{Data Protection.} While Law No. 2010/012 addresses cybercrime, a comprehensive \emph{Data Protection Law} is still under drafting. ANTIC has issued guidelines aligning with the African Union Convention on Cyber Security (Malabo Convention) and GDPR principles: purpose limitation, data minimisation, consent, and the right to access/correct one’s data. E‑health platforms must be especially vigilant: medical data is classified as \emph{sensitive} and requires explicit consent and enhanced encryption.

\begin{methode}[Risk Assessment Matrix for E‑Health Data Breaches]
\textbf{Step 1: Identify assets.} Patient records, prescription data, appointment schedules, billing info.
\textbf{Step 2: Identify threats.} Insider misuse, ransomware, phishing of staff credentials, lost unencrypted laptop.
\textbf{Step 3: Assess likelihood (1–5) and impact (1–5).} Example: Phishing of nurse credentials — Likelihood 4, Impact 4 (access to full EMR).
\textbf{Step 4: Compute risk score = Likelihood $\times$ Impact.} Score 16 $\rightarrow$ \textbf{High}. Prioritise mitigation: mandatory 2FA, phishing simulation training, encrypted disks.
\textbf{Step 5: Define controls and residual risk.} After 2FA, likelihood drops to 2 $\rightarrow$ residual score 8 (Medium). Document acceptance or further action.
\textbf{Step 6: Review quarterly or after any incident.}
\end{methode}

\courssub{The Digital Divide: Inclusion Strategies}

The digital divide in Cameroon has three pronounced dimensions:

\begin{itemize}
\item \textbf{Urban–Rural:} Internet penetration in Littoral (Douala) is estimated ~58\%; in Far North ~18\%. Rural users often rely on expensive, low‑bandwidth mobile data.
\item \textbf{Gender:} GSMA 2023 data indicates women in Cameroon are approximately 23\% less likely than men to use mobile Internet, due to affordability, literacy and social norms.
\item \textbf{Age:} Youth (15–35) are the primary users; older adults, who often need e‑health and e‑government most, are the least connected.
\end{itemize}

\textbf{Inclusion Strategies.}
\begin{enumerate}
\item \textbf{Community Telecentres (Multimedia Community Centres — CMC).} Supported by the Universal Service Access Fund, these provide shared Internet, training and e‑service assistance. The CMC in Ngaoundéré (Adamawa) helps farmers access market prices and e‑extension services.
\item \textbf{Subsidised Smartphones.} Public‑private partnership initiatives offer entry‑level 4G devices payable in monthly instalments via Mobile Money.
\item \textbf{Digital Literacy Curricula.} Integrating ICT skills into primary and secondary education (the new curriculum includes coding and digital citizenship) and adult literacy programmes.
\item \textbf{Local Language Interfaces.} E‑service portals in Fulfulde, Bamileke, Duala, etc., lower the language barrier. The \emph{e‑Santé} pilot in the North Region added a Fulfulde voice menu for appointment booking.
\end{enumerate}

\begin{propriete}[Trust Framework]
A \cle{Trust Framework} for e‑government combines four pillars:
\begin{enumerate}
\item \textbf{Strong Authentication:} Multi‑factor (password + OTP / digital certificate) verifying the user’s identity.
\item \textbf{Authorisation:} Role‑based access control (RBAC) ensuring users only reach data/functions relevant to their role.
\item \textbf{Audit Trails:} Immutable logs of every access, modification and transaction, timestamped and signed.
\item \textbf{Legal Recourse:} Clear liability rules, dispute resolution mechanisms and enforcement of Law No. 2010/012.
\end{enumerate}
Without all four, citizens and businesses will not trust the system, and adoption will stall.
\end{propriete}

\courssub{Case Analysis: Lessons from an E‑Procurement Pilot (2018–2020)}

Between 2018 and 2020, the Ministry of Public Contracts piloted an e‑procurement platform for tenders above 50 million FCFA. The technical solution (open‑source, hosted on government cloud) was sound. Yet, by 2021, less than 15\% of eligible tenders used it. A post‑mortem identified two root causes:

\begin{enumerate}
\item \textbf{Change‑Management Resistance.} Procurement officers were accustomed to paper‑based workflows and informal “negotiations”. The platform’s transparency removed discretionary power. No incentives or mandates were attached; usage remained voluntary. Middle managers silently discouraged staff by claiming “the system is too slow” (often due to their own lack of training).
\item \textbf{Inadequate Training.} A single two‑day workshop was organised for 200 officers. No follow‑up, no helpdesk, no user manuals in French/English. Officers forgot procedures, made errors, blamed the tool, and reverted to paper.
\end{enumerate}

\textbf{Lessons Learned.}
\begin{itemize}
\item Technical deployment is only 20\% of the project; change management, training, incentives and leadership sponsorship are 80\%.
\item A \cle{phased mandatory adoption} (e.g., “all tenders > 100M FCFA must use the platform from Jan 202X”) works better than voluntary pilots.
\item Continuous support (helpdesk, super‑users in each ministry) is non‑negotiable.
\end{itemize}

\begin{exoresolu}[Designing a Risk Mitigation Plan for a New E‑Health Platform]
\textbf{Statement:} The Ministry of Public Health plans to launch a national telemedicine platform connecting 50 district hospitals to 5 referral hospitals. You are the ICT consultant tasked with producing a risk mitigation plan focusing on data security and service continuity. Use the risk assessment matrix method.

\textbf{Solution:}
\begin{enumerate}
\item \textbf{Asset Inventory:} Patient video consultations (real‑time), stored EMR (PDF, images), prescription database, user credentials (doctors, nurses, admins).
\item \textbf{Threat Identification:}
\begin{itemize}
\item T1: Ransomware encrypting EMR database (Likelihood 3, Impact 5 $\rightarrow$ Score 15).
\item T2: Unauthorised access to video consultations via stolen doctor credentials (L 4, I 4 $\rightarrow$ 16).
\item T3: Network outage at district hospital during consultation (L 4, I 3 $\rightarrow$ 12).
\item T4: Insider (admin) exporting patient data (L 2, I 5 $\rightarrow$ 10).
\end{itemize}
\item \textbf{Mitigation Controls:}
\begin{itemize}
\item For T1: Offline encrypted backups daily; network segmentation; endpoint EDR on all hospital PCs. Residual L 1, I 5 $\rightarrow$ 5 (Medium).
\item For T2: Enforce hardware token 2FA (YubiKey) for all clinical staff; session timeout 5 min; audit log review weekly. Residual L 1, I 4 $\rightarrow$ 4 (Low).
\item For T3: Deploy dual SIM 4G routers (MTN + Orange) with automatic failover; store‑and‑forward queue for consultations. Residual L 2, I 2 $\rightarrow$ 4 (Low).
\item For T4: Role separation (no single admin has export + delete rights); data loss prevention (DLP) rules blocking bulk export to USB/email. Residual L 1, I 4 $\rightarrow$ 4 (Low).
\end{itemize}
\item \textbf{Governance:} Monthly risk review committee (MoH, ANTIC, hospital IT leads); quarterly penetration test; annual disaster recovery drill.
\item \textbf{Compliance Check:} All controls mapped to Law No. 2010/012 Articles 12, 23, 31 and ANTIC guidelines.
\end{enumerate}
\end{exoresolu}

\courssub{Legal and Regulatory Framework Summary}

\begin{center}
\begin{tabularx}{\linewidth}{|l|X|}\hline
\rowcolor{popblueL}
\textbf{Instrument} & \textbf{Relevance to E‑Services} \\ \hline
Law No. 2010/012 (Cybersecurity \& Cybercrime) & Criminalises attacks; mandates security measures; designates ANTIC as CERT/PKI authority. \\ \hline
Telecommunications Act (Law No. 98/014, amended 2010) & Licensing of operators; Universal Service Fund financing rural connectivity; quality of service obligations. \\ \hline
Electronic Transactions Framework (Decree No. 2015/0831/PM) & Legal recognition of electronic signatures, contracts, timestamps; defines trusted service providers. \\ \hline
ANTIC Data Protection Guidelines (2021) & Operationalises consent, data subject rights, breach notification, cross‑border transfer rules. \\ \hline
Public Procurement Code (Decree No. 2018/366) & Mandates e‑procurement for thresholds above 50M FCFA; sets transparency rules. \\ \hline
\end{tabularx}
\end{center}

\begin{savaistu}[ANTIC’s Role in Trust]
ANTIC (Agence Nationale des Technologies de l'Information et de la Communication) operates the national \cle{Root Certification Authority (RCA)}. It issues digital certificates to government entities (e.g., \texttt{*.gov.cm} TLS certificates, qualified electronic seals for e‑procurement). Since 2022, ANTIC also runs the \cle{National Cybersecurity Observatory}, publishing monthly threat bulletins and coordinating incident response for critical e‑services.
\end{savaistu}

\courssub{Consolidation Exercises}

\begin{exercice}[Exercise 1: Benefit–Challenge Mapping]
For each benefit below, identify the most critical challenge that could undermine it in the Cameroonian context, and propose one concrete mitigation.
\begin{enumerate}
\item 24/7 availability of e‑tax filing.
\item Broader market reach for a rural cooperative selling cocoa via an e‑commerce platform.
\item Transparency in e‑procurement.
\end{enumerate}
\end{exercice}

\begin{exercice}[Exercise 2: Risk Matrix Application]
A new e‑government portal will allow citizens to apply for birth certificates online. The portal stores scanned copies of parents’ IDs and the child’s hospital birth record. Construct a 3$\times$3 risk matrix (Likelihood: Low/Medium/High; Impact: Low/Medium/High) for the following threats:
\begin{itemize}
\item SQL injection on the application form.
\item Phishing campaign mimicking the portal’s login page.
\item Insider (civil registrar) selling bulk data to identity thieves.
\end{itemize}
For each, state the risk level (Low/Medium/High/Critical) and one technical + one organisational control.
\end{exercice}

\begin{exercice}[Exercise 3: Digital Divide Policy Brief]
Write a 200‑word policy brief for the Minister of Posts and Telecommunications proposing three priority actions to close the gender digital divide in the Far North Region by 2027. Reference specific barriers (cost, literacy, cultural norms) and measurable indicators.
\end{exercice}

\begin{corrige}[Exercise 1]
\begin{enumerate}
\item \textbf{Benefit:} 24/7 e‑tax filing. \textbf{Critical challenge:} Electricity/internet outages at cybercafés where most SMEs file. \textbf{Mitigation:} Offline-capable mobile app with background sync when connectivity returns; solar‑powered filing kiosks at tax centres.
\item \textbf{Benefit:} Market reach for cocoa cooperative. \textbf{Critical challenge:} Low digital literacy and trust in online payments among farmers. \textbf{Mitigation:} Cooperative‑level digital champions trained to manage the platform; integration with Mobile Money (MTN/Orange) with USSD fallback for feature phones.
\item \textbf{Benefit:} Transparency in e‑procurement. \textbf{Critical challenge:} Change resistance from procurement officers losing discretionary power. \textbf{Mitigation:} Mandatory adoption decree with phased timeline; performance bonuses tied to platform usage; independent audit of procurement integrity published annually.
\end{enumerate}
\end{corrige}

\begin{corrige}[Exercise 2]
\begin{center}
\begin{tabularx}{\linewidth}{|l|c|c|c|X|X|}\hline
\rowcolor{popblueL}
Threat & Likelihood & Impact & Risk Level & Technical Control & Organisational Control \\ \hline
SQL Injection & High & High & Critical & Parameterised queries, WAF, quarterly pen‑test & Secure coding training for devs; bug bounty programme \\ \hline
Phishing & High & High & Critical & DMARC/DKIM/SPF; hardware 2FA for staff; phishing‑resistant auth (WebAuthn) & Quarterly simulated phishing; clear reporting channel; revoke compromised creds in <1h \\ \hline
Insider Data Sale & Medium & Critical & High & DLP (block bulk export); RBAC (no single user has full export); encrypted DB with audit logs & Background checks; separation of duties; mandatory annual ethics training; whistle‑blower protection \\ \hline
\end{tabularx}
\end{center}
\end{corrige}

\begin{corrige}[Exercise 3 — Outline]
\textbf{Title:} Closing the Gender Digital Divide in Far North — 2024–2027 Priority Actions. \\
\textbf{Context:} Women in Far North are estimated 35\% less likely to use mobile Internet (GSMA 2023). Barriers: handset cost (avg. 45,000 FCFA), low female literacy (42\% vs 68\% male), cultural norms restricting women’s public mobility and technology use. \\
\textbf{Action 1 — Affordable Access:} Subsidise 4G smartphones for women’s groups via Universal Service Fund; target 50,000 devices by 2027. Indicator: \% women owning smartphone (baseline 18\% $\rightarrow$ target 45\%). \\
\textbf{Action 2 — Literacy \& Skills:} Deploy mobile digital literacy caravans (partner with MINPROFF, NGOs) teaching USSD, Mobile Money, e‑health appointment booking in local languages (Fulfulde, Kanuri). Indicator: 30,000 women certified “Digital Citizen Level 1” by 2027. \\
\textbf{Action 3 — Safe \& Relevant Content:} Fund development of e‑services women actually need: maternal health teleconsultation, market price info for petty trade, GBV reporting chatbot. Ensure female moderators. Indicator: Monthly active female users on priority e‑services > 20,000. \\
\textbf{Governance:} Inter‑ministerial steering committee (MINPOSTEL, MINPROFF, MINSANTE, ANTIC) reporting quarterly to Parliament.
\end{corrige}

\coursec{Computer‑Assisted Learning (CAL): Concepts, Tools and Design}

In the previous section we saw how electronic services transform commerce, health and government. Education is no exception: the same digital technologies that allow a patient in Bamenda to consult a doctor remotely also allow a student in a remote college to learn with the assistance of a computer. This section studies \cle{Computer‑Assisted Learning} (CAL): what it is, the learning theories behind it, the tools used to build it, and how to design a CAL lesson properly.

\courssub{What is Computer‑Assisted Learning?}

\textbf{Intuition first.} Imagine a student revising for the GCE using a tablet app that presents a lesson on spreadsheets, asks questions, marks the answers instantly, and gives extra practice where the student is weak. The computer is not replacing the teacher; it is \emph{assisting} the learning. That is the essence of CAL.

\begin{definition}[Computer‑Assisted Learning (CAL)]
\cle{Computer‑Assisted Learning} (also called Computer‑Assisted Instruction, CAI) is the use of computers and digital media to present instructional content, provide practice, give feedback and assess learners, in support of (or sometimes in place of) a human teacher.
\end{definition}

Three related terms must be carefully distinguished, because examiners like to test the difference:

\begin{center}
\begin{tabularx}{\linewidth}{|l|X|X|}
\hline
\rowcolor{popblueL} \textbf{Term} & \textbf{Description} & \textbf{Example} \\
\hline
\cle{CAL} & General term: any use of a computer to assist learning (drills, tutorials, simulations). & A typing tutor; a physics simulation. \\
\hline
\cle{CBT} & \emph{Computer‑Based Training}: structured, self‑contained courseware delivering a complete training sequence, often offline (CD, kiosk, local network). & Airline staff safety training module on a standalone kiosk. \\
\hline
\cle{ITS} & \emph{Intelligent Tutoring System}: CAL that uses artificial intelligence to model the learner and adapt instruction individually. & A maths tutor that diagnoses misconceptions and chooses the next exercise. \\
\hline
\end{tabularx}
\end{center}

\begin{definition}[Intelligent Tutoring System]
An \cle{Intelligent Tutoring System} (ITS) is a CAL system that uses \cle{artificial intelligence} techniques — typically a \emph{domain model} (what to teach), a \emph{learner model} (what the student knows) and a \emph{tutoring model} (how to teach) — to provide personalised instruction, hints and feedback adapted to each individual learner.
\end{definition}

\begin{attention}[Do not confuse the terms]
CBT is a \emph{delivery format} (a self‑contained training package); CAL is the \emph{general concept}; an ITS is a \emph{smart kind of CAL}. Every ITS is CAL, but most CAL is not intelligent. Writing "CAL and CBT are synonyms" is a common exam error.
\end{attention}

\courssub{Pedagogical Foundations of CAL}

Good CAL is not just technology; it is built on \cle{learning theories}. Three foundations matter at this level.

\textbf{1. Behaviourist drill‑and‑practice.} Rooted in behaviourism (learning as stimulus–response–reinforcement), the computer presents a question (stimulus), the learner answers (response), and the system gives immediate feedback (reinforcement). This is ideal for memorisation tasks: multiplication tables, French vocabulary, keyboard shortcuts in a spreadsheet.

\textbf{2. Constructivist simulations.} Constructivism holds that learners \emph{build} knowledge by exploring and experimenting. A simulation — for example, a virtual electric circuit or a model of an ecosystem in the Congo Basin — lets students change variables, observe consequences and discover principles themselves.

\textbf{3. Cognitive load theory.} Working memory is very limited. If a lesson floods the learner with decorative animations, background music and dense text simultaneously, the \emph{extraneous} load crowds out actual learning. Good design presents one idea at a time, synchronises narration with visuals, and removes anything that does not serve the objective.

\begin{attention}[Common mistake — multimedia overload]
A frequent design error is overloading a CAL module with animations, sound effects, colours and long videos "to make it attractive". This violates \cle{cognitive load} principles: the learner's working memory is spent processing decoration instead of the concept. \textbf{Rule of thumb:} every media element must serve a stated learning objective; if it does not, delete it.
\end{attention}

\courssub{Authoring Tools for CAL}

An \cle{authoring tool} is software used to create learning content without necessarily programming from scratch.

\begin{itemize}
\item \cle{Moodle}: a free, open‑source \emph{Learning Management System} (LMS) used to organise courses, activities, quizzes and grades. Widely used in Cameroonian universities (e.g. University of Buea, University of Yaoundé I).
\item \cle{Adobe Captivate}: a commercial tool for interactive simulations and software demonstrations (e.g. teaching how to use Excel).
\item \cle{H5P}: a free, open‑source framework for interactive HTML5 content — quizzes, interactive videos, drag‑and‑drops — that embeds in Moodle and other platforms.
\item \cle{Kolibri}: an open‑source platform designed for \emph{offline} learning, crucial where Internet connectivity is weak or expensive — a realistic constraint in many Cameroonian schools.
\end{itemize}

\begin{propriete}[SCORM]
\cle{SCORM} (\emph{Sharable Content Object Reference Model}) is the international \cle{standard for packaging and tracking e‑learning content}. A SCORM package (a ZIP file with a manifest) can be imported into any SCORM‑compliant platform, and the platform can track completion, scores and time spent. (Knowing the purpose of SCORM is examinable; its internal XML structure is not.)
\end{propriete}

\begin{methode}[Creating a SCORM‑compliant package with H5P for offline deployment on Kolibri]
\begin{enumerate}
\item \textbf{Create content:} build the interactive activity (quiz, interactive video) in H5P, inside an LMS such as Moodle with the H5P plugin.
\item \textbf{Export:} download the activity as an \texttt{.h5p} file, or wrap the whole lesson as a SCORM 1.2 package using an authoring tool.
\item \textbf{Validate:} test the package in a SCORM player to confirm that scores and completion are reported.
\item \textbf{Import to Kolibri:} on the Kolibri server (which can run on a low‑cost local machine), import the package as a lesson resource.
\item \textbf{Deploy offline:} learners connect to the Kolibri local Wi‑Fi network (no Internet needed) and work through the content; progress synchronises when the coach connects.
\end{enumerate}
\end{methode}

\courssub{Multimedia Integration and Gamification}

A CAL module typically combines several media types, each with a pedagogical role:

\begin{itemize}
\item \textbf{Text} — precise definitions and instructions; cheap in bandwidth.
\item \textbf{Audio} — narration supporting weak readers or language learning.
\item \textbf{Video} — demonstrations (e.g. a titration experiment) that would be costly to perform live.
\item \textbf{Interactive simulations} — constructivist exploration (circuits, statistics, maps).
\item \textbf{Gamification elements} — \cle{badges} for achievements, \cle{leaderboards} for motivation, progress bars and levels, which exploit game mechanics to sustain engagement.
\end{itemize}

\begin{attention}[Gamification is a means, not an end]
Leaderboards can motivate competitive students but discourage weaker ones. Use badges for \emph{mastery} (everyone can earn them) rather than only for ranking, and never let points replace actual learning objectives.
\end{attention}

\courssub{Learner Modelling and Adaptive Pathways}

A \cle{learner model} is the system's representation of what the student knows. It drives \emph{adaptive pathways} — different routes through the content for different learners.

\textbf{Rule‑based branching} is the simple approach: the designer writes explicit rules such as "if the quiz score $< 50\%$, send the learner to remedial activity R2; otherwise continue to the next unit." It is transparent and easy to build in Moodle (using conditional activities), but it only adapts in ways the designer anticipated.

\textbf{AI‑driven recommendation} uses data from many learners to predict which resource will help a given student next — similar to how a video platform recommends films. It can discover non‑obvious patterns, but requires large amounts of data and raises privacy questions.

\courssub{Assessment in CAL}

CAL supports three assessment mechanisms:

\begin{itemize}
\item \textbf{Formative quizzes:} short, frequent, low‑stakes tests with immediate feedback, helping the learner self‑regulate.
\item \textbf{Automated marking:} objective questions (multiple choice, true/false, numeric, matching) are marked instantly and consistently, freeing teacher time for essay‑type feedback.
\item \textbf{Learning analytics dashboards:} aggregated data (attempts, time on task, common wrong answers) displayed to the teacher, revealing which concepts the whole class struggles with.
\end{itemize}

\begin{methode}[Using Moodle's competency framework to map CAL activities to GCE ICT syllabus objectives]
\begin{enumerate}
\item \textbf{Define competencies:} in Moodle (Site administration $\rightarrow$ Competencies), create a framework and enter each GCE ICT syllabus objective as a competency (e.g. "Describe the benefits of e‑government").
\item \textbf{Link activities:} attach each quiz, H5P activity or lesson to the competency it assesses.
\item \textbf{Set completion rules:} require a threshold score (e.g. $70\%$) for the competency to be marked "proficient".
\item \textbf{Monitor:} use the competency breakdown report as a learning‑analytics dashboard to see which syllabus objectives each student has mastered.
\item \textbf{Remediate:} assign branching remedial activities to students not yet proficient.
\end{enumerate}
\end{methode}

\courssub{Designing a CAL Lesson: A Systematic Process}

Professional CAL is never improvised. The standard design cycle has five stages:

\begin{popfigure}
\begin{center}
\begin{tikzpicture}[
  box/.style={draw=popblue, thick, fill=popblueL, rounded corners=2pt, minimum width=3.0cm, minimum height=0.85cm, align=center, font=\small},
  arr/.style={-{Stealth[length=2.5mm]}, thick, popdark},
  node distance=0.5cm
]
\node[box, align=center] (needs){\textbf{Needs analysis}\\objectives \& learners};
\node[box, right=of needs, align=center] (story){\textbf{Storyboard}\\screens \& script};
\node[box, right=of story, align=center] (proto){\textbf{Prototype}\\build with tools};
\node[box, below=of proto, align=center] (pilot){\textbf{Pilot}\\test with learners};
\node[box, left=of pilot, align=center] (eval){\textbf{Evaluation}\\analytics \& feedback};
\draw[arr] (needs) -- (story);
\draw[arr] (story) -- (proto);
\draw[arr] (proto) -- (pilot);
\draw[arr] (pilot) -- (eval);
\draw[arr, dashed, poporange] (eval.west) .. controls +(-0.8,0.5) and +(-0.8,-0.5) .. (needs.west) node[midway, left, font=\scriptsize\itshape, text=poporange]{revise};
\end{tikzpicture}
\end{center}
\legende{The CAL lesson design cycle: evaluation feeds back into a revised needs analysis.}
\end{popfigure}

\begin{enumerate}
\item \textbf{Needs analysis:} identify the syllabus objectives, the learners' prior knowledge, and the available infrastructure (electricity, devices, connectivity).
\item \textbf{Storyboard:} sketch every screen — text, media, questions, feedback, navigation — before building anything.
\item \textbf{Prototype:} build a working version with an authoring tool (H5P, Captivate, Moodle).
\item \textbf{Pilot:} test with a small group of real learners; observe where they get stuck.
\item \textbf{Evaluation:} analyse quiz data and feedback, then revise — the cycle repeats.
\end{enumerate}

\begin{exoresolu}[Designing a CAL module on e‑government services]
A teacher wants a 20‑minute CAL module for Upper Sixth on "e‑government services in Cameroon". Propose a design following the five stages.
\tcblower
\textbf{Needs analysis:} Objective — "describe at least three e‑government services and their benefits". Learners have smartphones but unreliable Internet, so the module must work offline on Kolibri.

\textbf{Storyboard:} Screen 1 — short text/video introducing ANTIC and online services; Screen 2 — interactive simulation: "apply for a document online" drag‑and‑drop of steps; Screen 3 — H5P quiz (5 MCQs) with immediate feedback; Screen 4 — badge "E‑Gov Explorer" for scores $\geq 80\%$.

\textbf{Prototype:} build in H5P inside Moodle, export and import into Kolibri.

\textbf{Pilot:} test with 10 students; observe that the video buffers — replace it with narrated slides plus audio to reduce bandwidth.

\textbf{Evaluation:} dashboard shows question 4 (benefits of e‑government) failed by $60\%$; add a remedial branching activity and revise the storyboard. The cycle restarts.
\end{exoresolu}

\begin{exoresolu}[Choosing the right CAL approach for a given topic]
For each of the following topics, state whether a behaviourist drill‑and‑practice, a constructivist simulation, or an ITS would be most appropriate, and justify briefly.
\begin{enumerate}
\item Memorising the seven layers of the OSI model.
\item Understanding how changing resistor values affects current in a parallel circuit.
\item A personalised revision programme for GCE ICT Paper 2 that adapts to each student's weak areas.
\end{enumerate}
\tcblower
\begin{enumerate}
\item \textbf{Behaviourist drill‑and‑practice} — the task is pure recall; immediate feedback on repeated MCQs builds automaticity.
\item \textbf{Constructivist simulation} — the learner must explore cause‑effect relationships; a virtual circuit (e.g. PhET) lets them vary resistors and observe current changes.
\item \textbf{ITS} — requires a learner model to diagnose individual weaknesses and a tutoring model to select the next appropriate question; rule‑based branching would be too rigid for the variety of possible misconceptions.
\end{enumerate}
\end{exoresolu}

\begin{exercice}[CAL concepts and design]
\begin{enumerate}
\item Distinguish clearly between CAL, CBT and ITS, giving one example of each.
\item A designer adds background music, three fonts and a bouncing logo to every screen of a CAL module. Which learning theory does this violate, and why?
\item State the purpose of SCORM and name one tool that can produce SCORM content.
\item Explain the difference between rule‑based branching and AI‑driven recommendation in adaptive learning.
\item List the five stages of the CAL design cycle in order.
\item A school in a rural area has 20 tablets but no Internet connection. Which authoring tool and deployment platform combination would you recommend, and why?
\end{enumerate}
\end{exercice}

\begin{corrige}[Exercise 1]
\begin{enumerate}
\item CAL is the general use of computers to assist learning (e.g. a typing tutor); CBT is a self‑contained training package delivering a full course (e.g. staff safety training on a local kiosk); an ITS is CAL using AI with a learner model to personalise instruction (e.g. a maths tutor that adapts exercises to diagnosed weaknesses).
\item Cognitive load theory: the decorative elements create extraneous cognitive load, consuming limited working memory that should be devoted to the learning content.
\item SCORM is the standard for packaging e‑learning content so it can be shared between platforms and tracked (completion, scores). Tools: H5P (via wrappers), Adobe Captivate, Moodle.
\item Rule‑based branching follows explicit designer‑written rules (e.g. score $< 50\%$ $\rightarrow$ remediation); AI‑driven recommendation learns from data on many learners to predict the best next resource, adapting in ways not explicitly programmed.
\item Needs analysis $\rightarrow$ storyboard $\rightarrow$ prototype $\rightarrow$ pilot $\rightarrow$ evaluation (then revision).
\item H5P for authoring (free, HTML5, works offline once exported) + Kolibri for deployment (designed for offline‑first environments, runs on a local server, syncs when online). This combination respects the infrastructure constraints.
\end{enumerate}
\end{corrige}

\begin{aretenir}[Key points]
\begin{itemize}
\item \cle{CAL} = computer assisting learning; \cle{CBT} = self‑contained training package; \cle{ITS} = AI‑based personalised CAL.
\item CAL design rests on behaviourism (drill‑and‑practice), constructivism (simulations) and \cle{cognitive load theory} (avoid multimedia overload).
\item Tools: Moodle (LMS), Adobe Captivate, H5P (interactive content), Kolibri (offline deployment); \cle{SCORM} standardises packaging and tracking.
\item Adaptivity ranges from simple rule‑based branching to AI‑driven recommendation; assessment uses formative quizzes, automated marking and learning analytics.
\item Design cycle: needs analysis $\rightarrow$ storyboard $\rightarrow$ prototype $\rightarrow$ pilot $\rightarrow$ evaluation.
\end{itemize}
\end{aretenir}

\coursec{E‑Learning Models, Implementation and Evaluation in the Cameroonian Context}

In the previous section we designed Computer‑Assisted Learning resources: tutorials, drill‑and‑practice exercises, simulations and educational games that run \emph{inside} a classroom or a computer laboratory. But what happens when the learner is \emph{not} in the laboratory at all — when a student in Wum follows a lesson prepared in Yaoundé, or when a teacher in Bamenda upgrades her skills in the evening from her phone? We then leave the laboratory and enter the wider world of \cle{e‑learning}: the delivery of learning experiences through electronic networks, mainly the Internet. This final section answers three practical questions that the GCE examiner loves: \emph{which model of e‑learning should a school choose, how should it be implemented, and how do we know it actually works?}

\courssub{Classification of E‑Learning Models}

Imagine four Upper Sixth students. Amina attends a live video lesson every evening with a teacher in Douala. Bih downloads lessons from a platform and studies them whenever she has electricity. Che follows a course where half the work is done in class and half online. Derick, already working, follows a free massive course from a foreign university. All four are doing e‑learning, but in four different \cle{models}.

\begin{definition}[E‑Learning Models]
E‑learning is commonly classified into four models according to \emph{when} and \emph{how} learners and teachers interact:
\begin{itemize}
\item \cle{Synchronous e‑learning}: teacher and learners are connected \emph{at the same time} (live video conferencing, webinars, chat sessions). Example: a Zoom revision class.
\item \cle{Asynchronous e‑learning}: learning materials are placed on a \emph{Learning Management System} (LMS) such as Moodle, and learners access them \emph{at their own time} (recorded videos, forums, quizzes, downloadable notes).
\item \cle{Blended learning}: a deliberate combination of face‑to‑face classroom teaching and online activities.
\item \cle{MOOCs} (Massive Open Online Courses): free or low‑cost online courses open to a very large number of participants, offered by platforms such as Coursera, edX or the francophone platform FUN.
\end{itemize}
\end{definition}

\begin{definition}[Blended Learning]
\cle{Blended learning} is a formal education programme in which a student learns at least in part through online delivery of content and instruction, with some element of student control over time, place, path or pace, and at least in part at a supervised brick‑and‑mortar location away from home. In short: \emph{traditional classroom methods combined with online digital media}, integrated so that each reinforces the other.
\end{definition}

\begin{exemplebox}[Choosing a model for a Cameroonian school]
A Government High School in the North‑West Region has one computer laboratory, an unstable 3G connection and motivated teachers. A purely synchronous model would collapse whenever the network fails. A \textbf{blended model} is chosen: theory notes, quizzes and videos are hosted on a local Moodle server (asynchronous, usable even when the Internet is down), while face‑to‑face sessions in the laboratory are used for practical work and remediation. The model fits the constraint instead of fighting it.
\end{exemplebox}

\courssub{Platform Selection Criteria}

Once the model is chosen, a platform must be selected. This decision is technical \emph{and} pedagogical. The main criteria, in the Cameroonian context, are:

\begin{center}
\begin{tabularx}{\linewidth}{|l|X|}
\hline
\rowcolor{popblueL} \textbf{Criterion} & \textbf{Question to ask} \\
\hline
Bandwidth requirement & Does the platform work on a slow 2G/3G connection, or does it need broadband? \\
\hline
Offline capability & Can content be downloaded and used without Internet (crucial outside Yaoundé and Douala)? \\
\hline
Language support & Does it fully support both French and English interfaces and content? \\
\hline
Cost & Is it free/open‑source (Moodle, Kolibri) or subscription‑based (Canvas)? \\
\hline
Local hosting & Can it be installed on a school server or a low‑cost device (e.g.\ a Raspberry Pi) so it works even with no Internet? \\
\hline
Ease of use & Can an average teacher create a course without advanced training? \\
\hline
\end{tabularx}
\end{center}

\begin{popfigure}
\begin{tikzpicture}
\begin{polaraxis}[
  width=11cm, height=11cm,
  xtick={0,72,144,216,288},
  xticklabels={Low bandwidth, Offline use, FR/EN support, Low cost, Ease of use},
  ytick={2,4,6,8,10},
  ymin=0, ymax=10,
  yticklabel style={font=\tiny, color=popmuted},
  xticklabel style={font=\small},
  grid=both, major grid style={popline!50},
  legend style={at={(1.35,1.02)}, anchor=north west, font=\small}
]
\addplot+[thick, mark=*, color=popblue] coordinates {(0,4) (72,5) (144,9) (216,10) (288,7) (360,4)};
\addlegendentry{Moodle}
\addplot+[thick, mark=square*, color=poporange] coordinates {(0,5) (72,3) (144,7) (216,9) (288,9) (360,5)};
\addlegendentry{Google Classroom}
\addplot+[thick, mark=triangle*, color=popgreen] coordinates {(0,9) (72,10) (144,7) (216,10) (288,6) (360,9)};
\addlegendentry{Kolibri}
\addplot+[thick, mark=diamond*, color=poppurple] coordinates {(0,3) (72,2) (144,8) (216,3) (288,8) (360,3)};
\addlegendentry{Canvas}
\end{polaraxis}
\end{tikzpicture}
\legende{Radar chart comparing four e‑learning platforms on five criteria relevant to Cameroon (scores out of 10, indicative). Kolibri dominates on offline use and low bandwidth; Moodle is the best free all‑rounder; Canvas is powerful but costly and bandwidth‑hungry.}
\end{popfigure}

The radar chart shows why \cle{Kolibri} (an open‑source platform designed for low‑resource environments, installable on a small local server) is often recommended for rural Cameroonian schools, while \cle{Moodle} remains the standard choice for institutions with a server and a technician.

\begin{attention}[Do not copy foreign choices blindly]
A platform that is excellent in a country with fibre‑optic Internet may be unusable in a school where the connection is a shared 3G key. Always test the platform \emph{under real local conditions} before adopting it: measure loading times on the actual school network, in both languages, on the cheapest smartphone students own.
\end{attention}

\courssub{Implementation Steps for E‑Learning in a School}

Deploying e‑learning is a project, not a purchase. A sound implementation follows five ordered steps.

\begin{methode}[Implementing e‑learning in a Cameroonian school]
\begin{enumerate}
\item \textbf{Policy alignment.} Check conformity with national orientations: MINESEC circulars and pedagogical instructions on the use of ICT in education, the school\'s internal regulations, and data‑protection obligations (law on cybersecurity and personal data, with ANTIC as the national authority). Obtain approval from the administration and the Parent‑Teacher Association.
\item \textbf{Teacher training.} Train teachers first on basic ICT use, then on the chosen platform, then on online pedagogy (writing quizzes, giving feedback, animating a forum). No platform survives untrained teachers.
\item \textbf{Content localisation.} Adapt or create content that follows the official MINESEC syllabus, in the correct language of the sub‑system (English for Anglophone schools, French for Francophone schools), using local examples (mobile money, local crops, Cameroonian history).
\item \textbf{Infrastructure provisioning.} Provide the laboratory or local server, a reliable power solution (solar backup where ENEO supply is unstable), student devices or a bring‑your‑own‑device policy, and an Internet bundle negotiated with an ISP (CAMTEL, MTN, Orange).
\item \textbf{Monitoring and continuous improvement.} Collect usage statistics from the LMS, hold termly review meetings with teachers, and adjust content, training and infrastructure each year.
\end{enumerate}
\end{methode}

\courssub{Quality Assurance}

Who guarantees that an online course is \emph{good}? Two complementary instruments exist.

\textbf{The UNESCO ICT Competency Framework for Teachers (ICT‑CFT).} This international reference describes the ICT skills teachers should master, organised in three progressive levels: \emph{Knowledge Acquisition} (teachers use ICT to support standard teaching), \emph{Knowledge Deepening} (teachers use ICT to manage collaborative, problem‑based learning), and \emph{Knowledge Creation} (teachers and students use ICT to innovate and create new knowledge). A school can use it as a checklist to plan teacher training and to certify progress.

\textbf{National accreditation of online courses.} Before an online course is recognised (for example, to count towards a certificate), it should be validated by the competent national bodies (MINESEC for secondary education, MINESUP for higher education). Accreditation verifies syllabus coverage, assessment validity, teacher qualification and data protection.

\begin{aretenir}[Quality assurance in one sentence]
Train teachers against the UNESCO ICT‑CFT levels, and submit courses to national accreditation: quality e‑learning is \emph{people plus validation}, not just software.
\end{aretenir}

\courssub{Evaluating E‑Learning: Kirkpatrick\'s Four Levels}

A headmaster asks: ``Our e‑learning project cost 3 million FCFA. Was it worth it?'' To answer rigorously, evaluators use \cle{Kirkpatrick\'s four‑level model}, originally designed for training evaluation and perfectly applicable to e‑learning.

\begin{definition}[Kirkpatrick\'s Four Levels]
\begin{itemize}
\item \textbf{Level 1 — Reaction:} Did learners \emph{like} the training? (satisfaction questionnaires at the end of the course).
\item \textbf{Level 2 — Learning:} Did learners \emph{acquire} the intended knowledge and skills? (pre‑test vs post‑test scores).
\item \textbf{Level 3 — Behaviour:} Do learners \emph{apply} what they learned in real situations? (classroom observation, practical tasks weeks later).
\item \textbf{Level 4 — Results:} Did the training produce the final \emph{impact} sought? (GCE pass rate, dropout rate, employability).
\end{itemize}
\end{definition}

\begin{methode}[Conducting a Kirkpatrick Level‑2 test with pre‑/post‑quiz on Moodle]
\begin{enumerate}
\item Before the online unit, create a 10‑question Moodle quiz covering the unit\'s objectives and let all students take it as a \textbf{pre‑test} (do not count it in the term mark).
\item Run the e‑learning unit (lessons, videos, forums) for the planned period.
\item After the unit, administer an equivalent \textbf{post‑test} (same objectives, different questions, same difficulty).
\item Export the grades from Moodle and compute, for each student, the \textbf{gain} $=$ post‑test score $-$ pre‑test score, then the class average gain.
\item Interpret: a large positive average gain indicates Level‑2 learning occurred; a near‑zero gain means the content or delivery must be revised — regardless of how much students \emph{enjoyed} the course.
\end{enumerate}
\end{methode}

\begin{exoresolu}[Evaluating a blended ICT course in a Yaoundé college]
Statement: A school introduces a blended e‑learning programme for Upper Sixth ICT. At the end, a satisfaction survey shows $92\%$ of students ``liked the course''. The pre‑test average was $8.2/20$ and the post‑test average $12.6/20$. Three months later, teachers observe that students now use spreadsheets correctly in physics practicals. The school\'s GCE ICT pass rate rises from $61\%$ to $74\%$. Analyse this evaluation using Kirkpatrick\'s model.
\tcblower
\textbf{Solution.}
\begin{itemize}
\item \textbf{Level 1 (Reaction):} $92\%$ satisfaction — positive reaction. Necessary but not sufficient.
\item \textbf{Level 2 (Learning):} average gain $= 12.6 - 8.2 = 4.4$ points out of 20, i.e.\ a $22\%$ relative improvement — real learning occurred.
\item \textbf{Level 3 (Behaviour):} transfer observed in physics practicals — students apply ICT skills outside the ICT class.
\item \textbf{Level 4 (Results):} GCE pass rate up by 13 points — the institutional objective is met.
\end{itemize}
Conclusion: the programme is validated on \emph{all four levels}; the investment is justified.
\end{exoresolu}

\begin{attention}[The Level‑1 trap]
The most common evaluation mistake is to distribute a satisfaction questionnaire, obtain smiles, and declare success. Level 1 measures \emph{feelings}, not \emph{learning}. A course can be loved and useless. Always climb to Level 2 (pre‑/post‑tests) and, for strategic decisions, to Level 4 (examination results such as the GCE pass rate). Conversely, never claim Level‑4 impact from Level‑1 data.
\end{attention}

For Cameroonian schools, indicators can be adapted concretely: Level 3 may be measured by whether students use their skills for real tasks (filling online baccalaureate registration, using mobile money safely), and Level 4 by GCE results, transition rates to university, or success in national competitions.

\courssub{Future Trends in E‑Learning}

E‑learning evolves quickly. Four trends are already reshaping it and will matter for Cameroon:

\begin{itemize}
\item \textbf{AI‑driven learning analytics.} Platforms analyse every click and quiz answer to detect, \emph{before} the examination, which student is likely to fail which topic, and propose personalised remediation automatically.
\item \textbf{Micro‑credentials.} Instead of one big diploma after years of study, learners collect small certified proofs of specific skills.
\item \textbf{Blockchain certificates.} Certificates recorded on a blockchain cannot be forged and can be verified instantly by any employer — a strong weapon against fake diplomas.
\item \textbf{Mobile‑first learning for low‑bandwidth areas.} Since most Cameroonian learners connect through a phone rather than a computer, new platforms are designed for small screens and tiny data bundles first.
\end{itemize}

\begin{propriete}[Micro‑credential]
A \cle{micro‑credential} is a certified, \emph{stackable} digital badge representing verified mastery of a \emph{specific} skill (e.g.\ ``Create a Moodle quiz'' or ``Configure a Wi‑Fi router''). \emph{Stackable} means several micro‑credentials can be accumulated and combined towards a larger qualification. Unlike a diploma, it is granular, quickly earned and digitally verifiable. (The definition is examinable; no proof is required.)
\end{propriete}

\begin{methode}[Designing a low‑bandwidth mobile learning module with PWA technology]
\begin{enumerate}
\item \textbf{Choose the PWA approach.} A \emph{Progressive Web App} is a website that behaves like a mobile app: it installs on the phone\'s home screen and, crucially, works \emph{offline} thanks to cached content.
\item \textbf{Lighten the content.} Replace heavy videos with audio plus slides, compress images, and keep each page under a few hundred kilobytes so it loads on 2G/3G.
\item \textbf{Cache for offline use.} Configure the PWA\'s \emph{service worker} so that lessons, quizzes and diagrams are stored on the phone at first visit and remain available without any connection.
\item \textbf{Synchronise later.} Quiz answers are stored locally and uploaded automatically when the student next finds a connection (e.g.\ in town).
\item \textbf{Test in real conditions.} Verify the module on an entry‑level Android phone, in English and French, with mobile data throttled to 3G, before deploying it to students.
\end{enumerate}
\end{methode}

\begin{savaistu}[E‑learning and Cameroon]
Cameroon\'s universities (such as the Universities of Yaoundé, Buea and Douala) increasingly use Moodle‑based platforms, and national digital‑economy strategies encourage distance education to reach learners in remote areas. During the COVID‑19 school disruptions, radio, television and WhatsApp‑based lessons showed that low‑tech e‑learning — not only sophisticated platforms — can keep learning alive. The lesson for designers: in Cameroon, \emph{resilience beats sophistication}.
\end{savaistu}

\begin{exercice}[Choosing and justifying a model]
A rural Government High School has 40 Android tablets, no Internet in the village, but a teacher who travels to town weekly. (a) Propose an e‑learning model and platform. (b) Justify your choice using at least three selection criteria. (c) State which Kirkpatrick levels you would evaluate in the first year and how.
\end{exercice}

\begin{corrige}[Exercise 1]
(a) An \textbf{asynchronous, offline‑first model} using \textbf{Kolibri} installed on a local server (or a Raspberry Pi) in the school; the travelling teacher downloads new content in town weekly and updates the server. (b) Justification: \emph{offline capability} (Kolibri works fully without Internet), \emph{bandwidth} (no connection needed in the village), \emph{cost} (Kolibri is free and open‑source), \emph{language support} (content can be loaded in English and French). (c) First year: Level 1 with an end‑of‑term satisfaction questionnaire; Level 2 with pre‑/post‑quizzes on the tablets; Level 4 is premature in year one but GCE/sub‑regional results should be tracked from year two.
\end{corrige}

\begin{aretenir}[Chapter summary — what to remember for the GCE]
\begin{itemize}
\item Four e‑learning models: \textbf{synchronous}, \textbf{asynchronous}, \textbf{blended}, \textbf{MOOCs}.
\item Platform choice in Cameroon: bandwidth, offline use, FR/EN support, cost, local hosting.
\item Implementation: policy alignment $\rightarrow$ teacher training $\rightarrow$ content localisation $\rightarrow$ infrastructure $\rightarrow$ monitoring.
\item Quality: UNESCO ICT‑CFT for teachers; national accreditation for courses.
\item Evaluation: Kirkpatrick — Reaction, Learning, Behaviour, Results; never stop at Level 1.
\item Future: AI analytics, micro‑credentials, blockchain certificates, mobile‑first PWA design.
\end{itemize}
\end{aretenir}

\newpage

\coursec{Integration activity}

\courssub{Aims of the activity}

This integration activity brings together everything you have studied in this chapter: the landscape of electronic services (e-commerce, e-health, e-government) from Lesson 42, their applications, benefits and challenges from Lesson 43, and the design of Computer-Assisted Learning (CAL) and e-learning solutions from Lesson 44. By the end of the activity, you should be able to:

\begin{itemize}
\item \textbf{Analyse} a real community context and \textbf{map} the electronic services that exist (or are missing) there;
\item \textbf{Identify} a service gap and \textbf{justify} why an electronic service is an appropriate response;
\item \textbf{Design} a low-bandwidth e-service architecture adapted to local infrastructure (network coverage, devices, electricity, literacy);
\item \textbf{Apply} CAL design principles to produce a training module for the users of the new service;
\item \textbf{Plan} an evaluation of the project using a recognised model (Kirkpatrick's four levels);
\item \textbf{Anticipate} risks (technical, human, financial, ethical/legal) and propose mitigation measures;
\item \textbf{Communicate} a technical proposal clearly, in the form of a structured pitch with diagrams, a storyboard and a risk table.
\end{itemize}

\begin{aretenir}[Skills mobilised]
Analysis of a context $\bullet$ Synthesis of chapter knowledge (Lessons 42--44) $\bullet$ Design (architecture, storyboard) $\bullet$ Critical evaluation $\bullet$ Teamwork and oral presentation. These are exactly the higher-order skills rewarded in the GCE Advanced Level problem-solving questions.
\end{aretenir}

\courssub{Situation}

\textbf{Read the following case carefully. All the data you need are contained in it.}

\begin{experience}[The Mbengwi Rural Health District]
Mbengwi is a rural health district in the North-West Region of Cameroon. The district hospital serves about 85,000 inhabitants spread over 42 villages, many reachable only by unpaved roads that become difficult during the rainy season. The district has one district hospital and 11 integrated health centres, staffed by 38 community health workers (CHWs), most of whom hold a secondary school certificate plus a short nursing auxiliary training.

\textbf{Observed problem.} The district medical officer reports that maternal deaths and complications are frequently linked to \emph{late referral}: a pregnant woman in a remote village develops danger signs (severe headache, bleeding, swollen face), the CHW is unsure whether the case is urgent, and the patient travels to the district hospital too late. There is currently \textbf{no electronic service} connecting CHWs to the district doctor for advice.

\textbf{Available infrastructure (survey results).}
\begin{itemize}
\item Mobile network: 2G voice/SMS coverage in all 42 villages; 3G/4G data coverage in only 9 villages (including Mbengwi town). Average measured bandwidth in 3G zones: about 0.8 Mbit/s, unstable.
\item Devices: 34 of the 38 CHWs own a basic feature phone (calls + SMS + USSD); 11 also own a low-end Android smartphone. The district hospital has 4 desktop computers and a VSAT internet connection.
\item Electricity: the hospital runs on a generator and solar panels; 6 health centres have solar kits; many villages have no grid power. CHWs charge phones at solar kiosks for about 100 FCFA per charge.
\item Mobile money: MTN Mobile Money and Orange Money agents exist in Mbengwi town and 5 other villages.
\item Languages: most CHWs are comfortable in English and Pidgin; many patients speak mainly the local language.
\item Digital skills: CHWs use SMS and WhatsApp daily but have never followed a structured online course.
\end{itemize}

\textbf{Budget and partners.} A local NGO offers a one-year pilot budget of 6,000,000 FCFA, covering software hosting, SMS/USSD costs, training and airtime allowances. The Regional Delegation of Public Health supports the project and requires a written evaluation report after six months.
\end{experience}

Your group of four has been hired as \textbf{young ICT consultants}. You must propose a viable solution and defend it in a 10-minute pitch before a jury composed of the district medical officer, the NGO coordinator and your teacher.

\courssub{Tasks}

\textbf{Task 1 — Mapping existing e-services.}\par
List the electronic services (e-health, e-commerce, e-government) that already exist or are realistically accessible in the district, using the data above. Present your answer as a table with three columns: \emph{Service}, \emph{Category}, \emph{Who can actually use it and why}.

\textbf{Task 2 — Identifying and justifying the gap.}\par
State precisely the service gap you will address. Explain, with at least three arguments drawn from the situation, why an e-health service is appropriate here, and name two expected benefits for each group of stakeholders (patients, CHWs, the district doctor).

\textbf{Task 3 — Designing the service architecture.}\par
Propose a low-bandwidth architecture combining \textbf{USSD} (for feature phones) and a \textbf{lightweight mobile application} (for the 11 smartphones), connected to a server at the district hospital. Your answer must include:
\begin{enumerate}
\item a labelled architecture diagram (users $\rightarrow$ channels $\rightarrow$ server $\rightarrow$ doctor);
\item a description of a complete tele-advice session (from the CHW's first action to the doctor's reply), indicating what happens on each channel;
\item two design decisions justified by the infrastructure data (e.g., why USSD menus rather than video calls; why store-and-forward rather than real-time).
\end{enumerate}

\textbf{Task 4 — Designing the CAL training module.}\par
The 38 CHWs must be trained to use the new service. Design a CAL module (for example built with H5P) that can run on smartphones and on the hospital computers. Provide:
\begin{enumerate}
\item the module's learning objectives (at least three, written with action verbs);
\item a storyboard of at least five screens (title, content, interactivity, feedback);
\item one justification of a CAL principle you applied (immediate feedback, self-pacing, multimedia, or formative assessment).
\end{enumerate}

\textbf{Task 5 — Evaluation plan (Kirkpatrick's model).}\par
For each of Kirkpatrick's four levels (Reaction, Learning, Behaviour, Results), state one indicator you will measure after six months and the tool you will use to measure it.

\textbf{Task 6 — Risk-mitigation table and pitch.}\par
Complete a risk table with at least four risks (technical, human, financial, ethical/legal — think of patient data confidentiality) and a mitigation measure for each. Then prepare your 10-minute pitch: architecture diagram, storyboard, risk table, and a clear division of speaking roles among the four group members.

\courssub{Model answer}

\textbf{Task 1 — Mapping existing e-services.}

\begin{center}
\begin{tabularx}{\linewidth}{|l|l|X|}
\hline
\rowcolor{popblueL} \textbf{Service} & \textbf{Category} & \textbf{Who can use it and why} \\
\hline
SMS / voice calls & Basic e-communication & All 38 CHWs (2G everywhere, feature phones) \\
\hline
USSD mobile money (MTN MoMo, Orange Money) & E-commerce / e-payment & CHWs and patients near the 6 agent villages; works on feature phones \\
\hline
WhatsApp messaging & E-communication & Only the 11 smartphone owners in 3G/4G zones \\
\hline
Hospital VSAT internet + e-mail & E-government / administration & District hospital staff only (4 computers) \\
\hline
Tele-consultation for CHWs & E-health & \textbf{Does not exist} — this is the gap \\
\hline
\end{tabularx}
\end{center}

\textbf{Task 2 — The gap and its justification.}\par
\emph{Gap:} there is no electronic channel allowing a CHW in a remote village to obtain rapid medical advice from the district doctor for a doubtful maternal case.

\emph{Why an e-health service is appropriate:}
\begin{enumerate}
\item 2G coverage reaches all 42 villages, so a USSD/SMS-based service can cover 100\% of CHWs without new infrastructure;
\item the problem is \emph{late referral}, i.e., an information and decision problem, exactly what tele-advice addresses;
\item the pilot budget (6,000,000 FCFA) is compatible with SMS/USSD costs, which are low, unlike video-conferencing equipment.
\end{enumerate}

\emph{Expected benefits.}
\begin{itemize}
\item \textbf{Patients:} faster referral decisions, fewer dangerous journeys.
\item \textbf{CHWs:} professional support, increased competence and confidence.
\item \textbf{District doctor:} better triage, fewer unnecessary transfers, records of cases for planning.
\end{itemize}

\textbf{Task 3 — Service architecture.}

\begin{popfigure}
\begin{tikzpicture}[font=\small, >=stealth, node distance=1.2cm]
\node[draw, rounded corners, fill=popgreenL, text width=2.8cm, align=center] (chw) {CHW feature phone\\ (USSD/SMS)};
\node[draw, rounded corners, fill=popgreenL, text width=2.8cm, align=center, below=of chw] (app) {CHW smartphone\\ (light app)};
\node[draw, rounded corners, fill=popblueL, text width=3.0cm, align=center, right=3.0cm of chw] (gw) {Telecom gateway\\ (USSD/SMS)};
\node[draw, rounded corners, fill=popblueL, text width=3.0cm, align=center, below=of gw] (srv) {District server\\ (case database)};
\node[draw, rounded corners, fill=popblueL, text width=2.8cm, align=center, right=3.0cm of srv] (doc) {Doctor's PC / smartphone\\ (alert + reply)};
\draw[->, thick, popdark] (chw) -- node[above, text width=2.4cm, align=center]{2G USSD session} (gw);
\draw[->, thick, popdark] (app) -- node[above, sloped]{3G/4G or SMS fallback} (srv);
\draw[<->, thick, popdark] (gw) -- (srv);
\draw[<->, thick, popdark] (srv) -- node[above]{VSAT / SMS alert} (doc);
\end{tikzpicture}
\legende{Low-bandwidth tele-advice architecture for the Mbengwi district.}
\end{popfigure}

\emph{A complete tele-advice session:}
\begin{enumerate}
\item The CHW dials the service code (e.g., *123\#) on any phone; a USSD menu asks for patient ID, age, and danger signs (1 = bleeding, 2 = severe headache, 3 = convulsions\ldots).
\item The gateway forwards the structured case to the district server, which stores it and immediately sends an SMS alert to the doctor on duty.
\item The doctor reads the case (on the hospital PC or smartphone) and replies with a coded instruction: \emph{refer urgently}, \emph{observe 24 h}, or \emph{request more information}.
\item The reply reaches the CHW by SMS (or inside the app for smartphone users, with the possibility of attaching a photo when 3G is available).
\item The server logs the whole exchange for follow-up and evaluation.
\end{enumerate}

\emph{Justified design decisions:}
\begin{itemize}
\item \textbf{USSD menus instead of video calls:} video requires stable bandwidth far above the measured 0.8 Mbit/s and smartphones that 27 of the 38 CHWs do not have; USSD works on every feature phone over 2G.
\item \textbf{Store-and-forward instead of real-time:} the doctor may be in theatre or the network unstable; the case is stored on the server and answered as soon as possible, with SMS guaranteeing delivery when data coverage fails.
\end{itemize}

\textbf{Task 4 — CAL training module.}\par
\emph{Learning objectives.} By the end of the module, the CHW will be able to: (1) navigate the USSD menu and submit a complete case; (2) interpret the doctor's coded replies and act accordingly; (3) apply confidentiality rules when handling patient data.

\emph{Storyboard (H5P, runs offline-capable on smartphones and hospital PCs):}

\begin{center}
\begin{tabularx}{\linewidth}{|c|X|X|}
\hline
\rowcolor{popblueL} \textbf{Screen} \& \textbf{Content} \& \textbf{Interactivity / feedback} \\
\hline
1 \& Welcome, objectives, why tele-advice saves mothers \& Tap to continue \\
\hline
2 \& Short illustrated video/animation: a danger-sign scenario \& Pause-and-think question \\
\hline
3 \& Step-by-step USSD simulation (mock menu) \& Learner types the correct codes; immediate correction \\
\hline
4 \& Interactive case study (H5P Branching Scenario): choose the action \& Different feedback per choice, with explanation \\
\hline
5 \& Final quiz (5 questions) + certificate \& Score displayed; retry allowed; summary of errors \\
\hline
\end{tabularx}
\end{center}

\emph{CAL principle applied:} \textbf{immediate feedback} — on screens 3 and 4, every learner action receives an instant, explained correction, which research and classroom experience show to be essential for self-paced learning without a tutor present.

\textbf{Task 5 — Evaluation plan (Kirkpatrick).}

\begin{center}
\begin{tabularx}{\linewidth}{|l|X|X|}
\hline
\rowcolor{popblueL} \textbf{Level} \& \textbf{Indicator} \& \textbf{Tool} \\
\hline
1. Reaction \& \% of CHWs satisfied with the module \& End-of-module questionnaire \\
\hline
2. Learning \& Average score at the final quiz \& H5P quiz results \\
\hline
3. Behaviour \& Number of correct tele-advice cases submitted per month \& Server logs \\
\hline
4. Results \& Average referral delay for maternal emergencies \& District health records (before/after comparison) \\
\hline
\end{tabularx}
\end{center}

\textbf{Task 6 — Risk-mitigation table.}

\begin{center}
\begin{tabularx}{\linewidth}{|l|X|X|}
\hline
\rowcolor{popblueL} \textbf{Type} \& \textbf{Risk} \& \textbf{Mitigation} \\
\hline
Technical \& Network outage or flat phone battery \& SMS fallback; solar charging points; paper referral form as backup \\
\hline
Human \& CHWs resist or misuse the system \& Hands-on CAL training, local champions, refresher sessions \\
\hline
Financial \& SMS/USSD costs exceed the budget \& Negotiate a bundle with the operator; monthly cost monitoring \\
\hline
Ethical / legal \& Breach of patient data confidentiality \& Patient codes instead of names, password-protected server, sensitisation on data protection (in line with ANTIC guidelines) \\
\hline
\end{tabularx}
\end{center}

\emph{Pitch organisation (10 minutes):} Member 1 presents the context and gap (2 min); Member 2 the architecture and a live or simulated USSD demonstration (3 min); Member 3 the CAL module storyboard (2 min); Member 4 the evaluation plan, risks and budget (2 min); 1 minute for the conclusion and jury questions.

\begin{attention}[Common pitfalls to avoid]
Do not propose solutions that ignore the data (video consultations over 2G, smartphones for everyone). Do not confuse the four Kirkpatrick levels. Do not forget confidentiality: health data are sensitive, and a technically brilliant project that exposes patient data is a failed project. Finally, keep every diagram labelled and every number consistent with the situation.
\end{attention}

\newpage

\coursec{Methods and worked examples}

\coursec{Electronic Services Landscape: E-Commerce, E-Health, E-Government}

\begin{savaistu}[Cameroon's digital transformation]
Cameroon's \emph{Stratégie Nationale de Développement 2020-2030} (SND30) identifies digital economy as a pillar. The \cle{Antic} (Agence Nationale des Technologies de l'Information et de la Communication) regulates cybersecurity and promotes trust in electronic services. Mobile Money (MTN MoMo, Orange Money) processes over 80\% of digital transactions, making it the backbone of e-commerce and e-government payments.
\end{savaistu}

\courssub{Concepts and Definitions}

\begin{definition}[Electronic Service (e-Service)]
An \cle{electronic service} is a service delivered \emph{through} a computer network (Internet, intranet, mobile network) where the interaction between provider and user, the processing of the request, and often the payment, are automated without physical contact. Three flagship domains are examined in this syllabus:
\begin{itemize}
\item \cle{E-commerce}: buying and selling goods or services online.
\item \cle{E-health}: use of ICT for healthcare delivery, information and management.
\item \cle{E-government}: online delivery of public services to citizens, businesses and other administrations.
\end{itemize}
\end{definition}

\begin{popfigure}
\begin{tikzpicture}[scale=0.9, every node/.style={font=\small}]
% Three main domains
\node[draw, rounded corners, fill=popblueL, minimum width=3.5cm, minimum height=1.2cm, align=center] (ecom) at (0,0){\textbf{E-Commerce}\\\textit{B2C, B2B, C2C, B2G}};
\node[draw, rounded corners, fill=popgreenL, minimum width=3.5cm, minimum height=1.2cm, align=center] (ehealth) at (5,0){\textbf{E-Health}\\\textit{Telemedicine, EHR, mHealth}};
\node[draw, rounded corners, fill=poporange, minimum width=3.5cm, minimum height=1.2cm, align=center] (egov) at (10,0){\textbf{E-Government}\\\textit{G2C, G2B, G2G}};

% Common infrastructure
\node[draw, rounded corners, fill=poppink, minimum width=12cm, minimum height=1cm] (infra) at (5,-2) {\textbf{Shared Digital Infrastructure: Internet / Mobile Networks (CAMTEL, MTN, Orange) -- Mobile Money -- ANTIC Security Framework -- Digital Identity}};

% Arrows
\draw[->, thick, popblue] (ecom.south) -- (infra.north west);
\draw[->, thick, popgreen] (ehealth.south) -- (infra.north);
\draw[->, thick, poporange] (egov.south) -- (infra.north east);

% Actors
\node[draw, circle, fill=popgold, minimum size=0.8cm] (citizen) at (-2,-3.5) {Citizen};
\node[draw, circle, fill=popgold, minimum size=0.8cm] (business) at (0,-3.5) {Business};
\node[draw, circle, fill=popgold, minimum size=0.8cm] (gov) at (2,-3.5) {Gov't};
\node[draw, circle, fill=popgold, minimum size=0.8cm] (health) at (4,-3.5) {Health};
\node[draw, circle, fill=popgold, minimum size=0.8cm] (admin) at (6,-3.5) {Admin};

\draw[popmuted] (citizen) -- (infra.south west);
\draw[popmuted] (business) -- (infra.south);
\draw[popmuted] (gov) -- (infra.south);
\draw[popmuted] (health) -- (infra.south);
\draw[popmuted] (admin) -- (infra.south east);
\end{tikzpicture}
\legende{The three pillars of electronic services sharing a common digital infrastructure in Cameroon}
\end{popfigure}

\courssub{E-Commerce: Models and Transaction Flow}

\begin{definition}[E-Commerce Models]
Electronic commerce is classified by the nature of the transacting parties:
\begin{center}
\begin{tabularx}{\linewidth}{|l|X|X|}\hline
\rowcolor{popblueL}
\textbf{Model} & \textbf{Full Name} & \textbf{Typical Example in Cameroon} \\ \hline
B2C & Business-to-Consumer & Jumia.cm, Kaymu, restaurant delivery apps (Glovo) \\ \hline
B2B & Business-to-Business & Wholesaler portal for retailers, procurement platforms \\ \hline
C2C & Consumer-to-Consumer & Facebook Marketplace, WhatsApp groups for second-hand goods \\ \hline
B2G / G2B & Business-to-Government & E-procurement (MARCHES PUBLICS ONLINE), tax filing (DGI) \\ \hline
C2G / G2C & Citizen-to-Government & Online birth certificate, passport appointment, land registry \\ \hline
\end{tabularx}
\end{center}
\end{definition}

\begin{methode}[Classifying an electronic service (e-commerce, e-health, e-government)]
When an examination question presents a scenario and asks you to identify the type of electronic service involved, proceed as follows:
\begin{enumerate}
\item \textbf{Identify the actors.} Who provides the service (business, hospital, ministry, school) and who receives it (customer, patient, citizen, learner)?
\item \textbf{Identify the object of the transaction.} Is it a sale of goods or services, a medical act, an administrative procedure, or a learning activity?
\item \textbf{Check the role of ICT.} Confirm that the service is delivered \emph{through} a computer network (website, mobile application, USSD, portal) and not merely supported by a computer used offline.
\item \textbf{Match with the standard categories:}
\begin{itemize}
\item \cle{E-commerce}: buying/selling of goods or services online (B2C, B2B, C2C, B2G).
\item \cle{E-health}: use of ICT for medical consultation, patient records, telemedicine, health information.
\item \cle{E-government}: online delivery of public services to citizens and businesses (G2C, G2B, G2G).
\end{itemize}
\item \textbf{Justify your answer} by quoting one concrete element of the scenario (payment, consultation, certificate, etc.).
\end{enumerate}
\textbf{Reflex:} always name the sub-category (e.g.\ B2C, G2C) — GCE marking guides reward the precise classification.
\end{methode}

\begin{exoresolu}[Classifying services in a Douala scenario]
A trader in Douala uses the following services: (a) she orders stock from a wholesaler's website and pays by MTN Mobile Money; (b) she books an appointment at Laquintinie Hospital through the hospital's online portal; (c) she downloads and fills her taxpayer declaration form on the website of the Directorate General of Taxation (DGI); (d) she sells second-hand phones to individuals through a Facebook page. For each case, state the type of electronic service and its precise sub-category, justifying your answer.
\tcblower
\textbf{Solution.}
\begin{enumerate}
\item \textbf{Case (a):} The actors are two businesses (the trader and the wholesaler) and the object is a purchase of goods paid online. This is \cle{e-commerce}, sub-category \textbf{B2B} (Business-to-Business), because the transaction takes place between two companies. The Mobile Money payment confirms that the transaction is completed electronically.
\item \textbf{Case (b):} The provider is a hospital and the object is a medical service (appointment booking). This is \cle{e-health}: ICT is used to deliver a health-related service to a patient. More precisely it is an online administrative health service (a first step towards telemedicine).
\item \textbf{Case (c):} The provider is a public administration (the DGI) and the user is a citizen/business performing an administrative obligation online. This is \cle{e-government}, sub-category \textbf{G2B} (Government-to-Business), since a tax declaration is filed by her enterprise; if she declared as a private individual it would be \textbf{G2C}.
\item \textbf{Case (d):} Goods are sold by one individual to other individuals via an online platform. This is \cle{e-commerce}, sub-category \textbf{C2C} (Consumer-to-Consumer), the platform playing the role of intermediary.
\end{enumerate}
\textbf{Examiner's remark:} the classification key is always the \emph{nature of the actors} and the \emph{object of the service}, never the device used (phone or computer).
\end{exoresolu}

\courssub{Secure E-Commerce Transaction Flow}

\begin{popfigure}
\begin{tikzpicture}[scale=0.85, every node/.style={font=\small}]
% Steps
\node[draw, rounded corners, fill=popblueL, minimum width=2.8cm, minimum height=1.1cm] (s1) at (0,0) {1. Browse \& Search};
\node[draw, rounded corners, fill=popblueL, minimum width=2.8cm, minimum height=1.1cm] (s2) at (3.5,0) {2. Select \& Cart};
\node[draw, rounded corners, fill=popgreenL, minimum width=2.8cm, minimum height=1.1cm] (s3) at (7,0) {3. Login / Auth};
\node[draw, rounded corners, fill=popgreenL, minimum width=2.8cm, minimum height=1.1cm] (s4) at (10.5,0) {4. Confirm Order};
\node[draw, rounded corners, fill=poporange, minimum width=2.8cm, minimum height=1.1cm] (s5) at (14,0) {5. Payment};
\node[draw, rounded corners, fill=poppurple, minimum width=2.8cm, minimum height=1.1cm] (s6) at (17.5,0) {6. Fulfilment};
\node[draw, rounded corners, fill=poppurple, minimum width=2.8cm, minimum height=1.1cm] (s7) at (21,0) {7. After-Sales};

% Arrows
\draw[->, thick] (s1.east) -- (s2.west);
\draw[->, thick] (s2.east) -- (s3.west);
\draw[->, thick] (s3.east) -- (s4.west);
\draw[->, thick] (s4.east) -- (s5.west);
\draw[->, thick] (s5.east) -- (s6.west);
\draw[->, thick] (s6.east) -- (s7.west);

% Security measures
\node[draw, rounded corners, fill=poppink, minimum width=3cm, minimum height=0.8cm, font=\tiny] (sec1) at (1.75,-2) {HTTPS, Padlock, Verify URL};
\node[draw, rounded corners, fill=poppink, minimum width=3cm, minimum height=0.8cm, font=\tiny] (sec2) at (8.75,-2) {Strong Password, 2FA (OTP)};
\node[draw, rounded corners, fill=poppink, minimum width=3cm, minimum height=0.8cm, font=\tiny] (sec3) at (14,-2) {Encrypted Gateway, MoMo PIN, SMS Confirm};
\node[draw, rounded corners, fill=poppink, minimum width=3cm, minimum height=0.8cm, font=\tiny] (sec4) at (19.25,-2) {Tracking, Digital Receipt, Return Policy};

\draw[dashed, popmuted] (s1.south) -- (sec1.north);
\draw[dashed, popmuted] (s2.south) -- (sec1.north);
\draw[dashed, popmuted] (s3.south) -- (sec2.north);
\draw[dashed, popmuted] (s4.south) -- (sec2.north);
\draw[dashed, popmuted] (s5.south) -- (sec3.north);
\draw[dashed, popmuted] (s6.south) -- (sec4.north);
\draw[dashed, popmuted] (s7.south) -- (sec4.north);

% Pillars
\node[draw, rounded corners, fill=popgold, minimum width=5cm, minimum height=0.7cm, font=\small] (pillars) at (10.5,-3.5) {Three Security Pillars: Confidentiality -- Authentication -- Integrity};
\end{tikzpicture}
\legende{Seven-step e-commerce flow with security checkpoints}
\end{popfigure}

\begin{methode}[Describing the steps of an online purchase and its security]
To answer ``Describe the process of an online purchase and the security measures involved'':
\begin{enumerate}
\item \textbf{List the stages in order:} (1) browsing/catalogue; (2) selection and virtual cart; (3) account creation or login (authentication); (4) order confirmation; (5) electronic payment (card, Mobile Money); (6) order processing and delivery; (7) after-sales service/feedback.
\item \textbf{Attach a security measure to each sensitive stage:}
\begin{itemize}
\item Login: strong password, possibly two-factor authentication (OTP by SMS).
\item Payment: encrypted connection (\cle{HTTPS}/SSL-TLS, visible padlock), trusted payment gateway, confirmation code (MoMo PIN).
\item Data: privacy policy, secure storage, non-disclosure of card/PIN data.
\end{itemize}
\item \textbf{Mention trust signals:} seller identity, customer reviews, return policy.
\item \textbf{Conclude with risks:} phishing sites, fake shops, SIM-swap fraud on mobile money — and the reflex of checking the URL and never sharing a PIN or OTP.
\end{enumerate}
\textbf{Key idea:} security in e-commerce rests on three pillars: \cle{confidentiality} (encryption), \cle{authentication} (who are you?), \cle{integrity} (the order is not modified).
\end{methode}

\begin{exoresolu}[Buying school shoes online]
Manka, a student in Bamenda, wants to buy school shoes from an online shop. (a) List, in order, the six main steps she will follow from search to delivery. (b) At which two steps is she most exposed to fraud, and what precautions should she take at each?
\tcblower
\textbf{Solution.}

\textbf{(a) The six steps:}
\begin{enumerate}
\item \textbf{Search and browsing:} Manka opens the shop's website or app and browses the catalogue of shoes (with photos, sizes and prices).
\item \textbf{Selection and cart:} she chooses her size and adds the shoes to her virtual shopping cart.
\item \textbf{Identification:} she creates an account or logs in with her username and password, and gives her delivery address in Bamenda.
\item \textbf{Order confirmation:} she checks the summary (items, price, delivery fees) and validates the order.
\item \textbf{Electronic payment:} she pays, for example with MTN Mobile Money: she dials the payment code, enters the amount and validates with her secret PIN; she receives a confirmation SMS.
\item \textbf{Processing and delivery:} the shop prepares the parcel and a delivery service brings it to her address; she confirms reception and may leave a review.
\end{enumerate}

\textbf{(b) The two most exposed steps:}
\begin{itemize}
\item \textbf{Step 3 (identification):} risk of password theft through a fake (phishing) website imitating the real shop. \emph{Precautions:} check that the address is exactly the official one and begins with \texttt{https://} with a padlock; use a strong, unique password; never follow a link received by an unsolicited SMS or WhatsApp message.
\item \textbf{Step 5 (payment):} risk of PIN theft or paying a fraudster posing as the shop. \emph{Precautions:} never communicate the Mobile Money PIN or any OTP code to anyone (no agent of MTN will ever ask for it); verify the merchant's name displayed before confirming the transaction; keep the confirmation SMS as proof of payment.
\end{itemize}
\textbf{Examiner's remark:} the answer is structured around the transaction flow; quoting the padlock, HTTPS and the secrecy of the PIN shows mastery of the three security pillars.
\end{exoresolu}

\courssub{E-Health: Scope and Cameroonian Initiatives}

\begin{definition}[E-Health Components]
\cle{E-health} covers:
\begin{itemize}
\item \cle{Telemedicine}: remote consultation (video, phone), diagnosis, prescription.
\item \cle{Electronic Health Records (EHR)}: digital patient history shared across facilities.
\item \cle{mHealth}: mobile-based health services (SMS reminders, USSD appointments, health tips).
\item \cle{Health Information Systems}: district health information software (DHIS2) for reporting.
\item \cle{E-learning for health workers}: continuous professional development via platforms.
\end{itemize}
\end{definition}

\begin{savaistu}[E-Health in Cameroon]
The \emph{Plan National de Développement Sanitaire (PNDS)} integrates e-health. Examples: \textbf{DHIS2} deployed in all health districts for epidemiological surveillance; \textbf{mHealth} pilots for maternal health SMS reminders in rural areas; \textbf{Telemedicine} links between Yaoundé central hospitals and regional hospitals (e.g. Garoua, Bertoua) for specialist consultations. The \cle{Antic} certifies health data hosting providers.
\end{savaistu}

\courssub{E-Government: Models and Cameroon's Portals}

\begin{definition}[E-Government Interaction Models]
\begin{center}
\begin{tabularx}{\linewidth}{|l|X|X|}\hline
\rowcolor{popblueL}
\textbf{Model} & \textbf{Description} & \textbf{Cameroon Example} \\ \hline
G2C & Government to Citizen & Birth/marriage/death certificates (ECAM), passport appointment, voter registration check \\ \hline
G2B & Government to Business & Tax declaration (DGI), customs (CAMCIS), business registration (CFCE), e-procurement \\ \hline
G2G & Government to Government & Integrated Financial Management Information System (IFMIS), inter-ministerial correspondence \\ \hline
G2E & Government to Employee & Payroll (SOLDE), e-recruitment (CONCOURS), e-leave management \\ \hline
\end{tabularx}
\end{center}
\end{definition}

\begin{savaistu}[Key Cameroonian e-Government Portals]
\begin{itemize}
\item \textbf{www.dgi.cm} -- Directorate General of Taxes: online tax declaration and payment.
\item \textbf{www.minfi.gov.cm} -- Ministry of Finance: budget transparency, procurement.
\item \textbf{www.camcis.cm} -- Customs: ASYCUDA World for clearance.
\item \textbf{www.eprocurement.cm} -- Public procurement online.
\item \textbf{www.antic.cm} -- Cybersecurity alerts, .cm domain management, certification.
\end{itemize}
\end{savaistu}

\coursec{Applications, Benefits and Challenges of Electronic Services}

\courssub{The 4P Analysis Grid}

\begin{methode}[Building a balanced benefits/challenges analysis]
GCE questions often ask: ``State three benefits and three challenges of introducing service X in Cameroon.'' Use the \textbf{4P grid} to avoid vague answers:
\begin{enumerate}
\item \textbf{People:} effect on users (convenience, time saved, digital literacy needed, exclusion of non-connected users).
\item \textbf{Purse (cost):} savings on travel and paperwork versus cost of equipment, data bundles, maintenance.
\item \textbf{Process:} speed, transparency, reduction of corruption and queues versus risk of breakdowns, fraud, cybercrime.
\item \textbf{Place (infrastructure):} coverage of electricity and Internet (urban/rural divide), reliability of networks (CAMTEL, MTN, Orange).
\end{enumerate}
Then:
\begin{enumerate}
\item Select the three strongest benefits and the three strongest challenges \emph{adapted to the scenario}.
\item Phrase each point as a complete sentence linking the service to its effect (``reduces\ldots because\ldots'').
\item For challenges, whenever possible suggest a mitigation (training, offline alternatives, ANTIC security guidelines).
\end{enumerate}
\textbf{Reflex:} never write one-word answers (``fast'', ``expensive''); each point must be developed.
\end{methode}

\begin{exoresolu}[E-government at the council level]
The Buea Council wants to allow citizens to apply for and pay for birth certificate duplicates online instead of queueing at the council office. Give three benefits of this project for citizens and three challenges the council is likely to face in Cameroon, suggesting one solution for each challenge.
\tcblower
\textbf{Solution.}

\textbf{Benefits for citizens:}
\begin{enumerate}
\item \textbf{Time and travel savings:} citizens no longer travel to Buea or queue for hours; a person living in Kumba or even abroad can apply from a smartphone, which reduces transport costs and lost working days.
\item \textbf{Transparency and reduced corruption:} an online procedure with fixed published fees and electronic payment (e.g.\ Mobile Money) limits direct cash handling and therefore reduces the risk of unofficial ``facilitation'' payments.
\item \textbf{Availability 24/7:} the portal is accessible at any time, unlike the council office which opens only during working hours, so citizens can apply at their convenience and track the progress of their file.
\end{enumerate}

\textbf{Challenges and possible solutions:}
\begin{enumerate}
\item \textbf{Digital divide:} many citizens, especially in rural areas, lack smartphones, Internet access or digital skills. \emph{Solution:} keep a physical desk at the council and create assisted digital kiosks where an agent helps citizens apply.
\item \textbf{Unreliable electricity and connectivity:} power cuts and network failures can interrupt the service. \emph{Solution:} host the portal on a reliable data centre with backup power and allow applications to be saved and resumed.
\item \textbf{Security and data protection:} civil-status data are sensitive; the system can be attacked or misused. \emph{Solution:} apply the cybersecurity guidelines of ANTIC (the national agency for ICT security), use secure connections (HTTPS), strong authentication and regular backups.
\end{enumerate}
\textbf{Examiner's remark:} notice how each challenge is paired with a realistic, Cameroonian mitigation — this is what earns full marks in the ``discuss'' questions.
\end{exoresolu}

\courssub{Comparative Table: Benefits vs Challenges Across Domains}

\begin{center}
\begin{tabularx}{\linewidth}{|l|X|X|}\hline
\rowcolor{popblueL}
\textbf{Domain} & \textbf{Key Benefits} & \textbf{Key Challenges (Cameroon Context)} \\ \hline
E-Commerce & Wider market, 24/7 sales, lower overhead, Mobile Money integration & Trust deficits (fake goods), logistics/last-mile delivery, payment disputes, cybercrime (phishing, SIM swap) \\ \hline
E-Health & Access to specialists (telemedicine), better records, reduced travel for patients, epidemic surveillance & Equipment cost, bandwidth for video, health worker training, data privacy, legal framework for e-prescription \\ \hline
E-Government & Transparency, reduced corruption, convenience, faster processing, open data & Digital divide, change resistance in administration, legacy systems integration, sustainability of platforms \\ \hline
\end{tabularx}
\end{center}

\coursec{Computer-Assisted Learning (CAL): Concepts, Tools and Design}

\courssub{Definition and Pedagogical Role}

\begin{definition}[Computer-Assisted Learning (CAL)]
\cle{Computer-Assisted Learning (CAL)} refers to the use of computer software \emph{designed specifically for learning} to present content, guide practice, simulate phenomena or assess learners. The computer \emph{assists} the teacher; it does not replace the pedagogical relationship. CAL is typically used \emph{locally} (installed software, CD-ROM, local network) without requiring a continuous Internet connection.
\end{definition}

\begin{aretenir}[CAL vs E-Learning -- The Distinction]
\begin{itemize}
\item \textbf{CAL:} software runs on a standalone machine or LAN; content is pre-packaged; interaction is learner--software; teacher is physically present.
\item \textbf{E-Learning:} learning delivered \emph{through a network} (Internet/intranet); includes communication (teacher-learner, learner-learner), activities and assessment via a platform (LMS); can be fully remote.
\end{itemize}
\textbf{Exam trap:} using a projector to show an educational CD-ROM in class is \textbf{CAL}, not e-learning.
\end{aretenir}

\courssub{Classification of CAL Software}

\begin{popfigure}
\begin{tikzpicture}[scale=0.85, every node/.style={font=\small}]
% Central node
\node[draw, rounded corners, fill=popblue, text=white, minimum width=3cm, minimum height=1.2cm] (cal) at (0,0) {\textbf{CAL Software Families}};

% Families
\node[draw, rounded corners, fill=popblueL, minimum width=3.2cm, minimum height=1.5cm, align=center] (drill) at (-5,3){\textbf{Drill \& Practice}\\\textit{Repetition, immediate feedback}\\\textit{Ex: vocabulary, math facts}};
\node[draw, rounded corners, fill=popgreenL, minimum width=3.2cm, minimum height=1.5cm, align=center] (tutor) at (0,3){\textbf{Tutorial}\\\textit{Step-by-step teaching}\\\textit{Ex: new concept discovery}};
\node[draw, rounded corners, fill=poporange, minimum width=3.2cm, minimum height=1.5cm, align=center] (sim) at (5,3){\textbf{Simulation}\\\textit{Virtual model of reality}\\\textit{Ex: chemistry, circuits}};
\node[draw, rounded corners, fill=poppurple, minimum width=3.2cm, minimum height=1.5cm, align=center] (game) at (-2.5,-2.5){\textbf{Educational Game}\\\textit{Learning through play}\\\textit{Ex: quest, puzzles}};
\node[draw, rounded corners, fill=poppink, minimum width=3.2cm, minimum height=1.5cm, align=center] (assess) at (2.5,-2.5){\textbf{Assessment / Testing}\\\textit{Quizzes, auto-marking}\\\textit{Ex: formative/summative}};

% Arrows
\foreach \target in {drill,tutor,sim,game,assess}
  \draw[->, thick, popmuted] (cal) -- (\target);

% Pedagogical goals
\node[draw, rounded corners, fill=popgold, minimum width=12cm, minimum height=0.8cm, font=\small] (goals) at (0,-4.5) {Pedagogical Goals: Memorisation $\rightarrow$ Drill | New Concepts $\rightarrow$ Tutorial | Danger/Cost/Invisible $\rightarrow$ Simulation | Motivation $\rightarrow$ Game | Evaluation $\rightarrow$ Assessment};
\end{tikzpicture}
\legende{CAL software families and their primary pedagogical goals}
\end{popfigure}

\begin{methode}[Selecting the appropriate type of CAL software]
Computer-Assisted Learning (CAL) software falls into classical families. To choose correctly in an exam:
\begin{enumerate}
\item \textbf{Read the learning objective} in the question (memorise, practise, explore, be assessed, simulate a phenomenon).
\item \textbf{Match it to the family:}
\begin{itemize}
\item \cle{Drill-and-practice}: repetition of exercises with immediate correction $\rightarrow$ memorisation, automatisms (times tables, vocabulary, irregular verbs).
\item \cle{Tutorial}: the computer plays the teacher, presenting content step by step with questions $\rightarrow$ discovering new notions individually.
\item \cle{Simulation}: a virtual model of a real phenomenon $\rightarrow$ experiments too dangerous, slow, costly or impossible in school lab (chemistry reactions, electric circuits, population dynamics).
\item \cle{Educational game}: learning embedded in a game scenario $\rightarrow$ motivation, problem-solving, strategic thinking.
\item \cle{Assessment / testing software}: quizzes with automatic marking $\rightarrow$ evaluation and instant feedback.
\end{itemize}
\item \textbf{Justify} the choice with two arguments: pedagogical (feedback, individual pace, interactivity, safety) and practical (cost, availability of computers, maintenance).
\item \textbf{Mention the teacher's role:} CAL \emph{assists} learning; the teacher selects the software, guides, debriefs and integrates it into the lesson plan.
\end{enumerate}
\textbf{Reflex:} if the situation involves danger, high cost or invisible phenomena $\rightarrow$ simulation. If it involves revision and repetition $\rightarrow$ drill-and-practice.
\end{methode}

\begin{exoresolu}[Equipping a college computer room]
A college in Ngaoundéré has a room of 20 multimedia computers connected to a local network. For each of the following needs, state the most appropriate type of CAL software and justify: (a) Form 3 students must master irregular English verbs; (b) the physics teacher wants students to study the combustion of sodium in chlorine, an experiment too dangerous to perform in the school laboratory; (c) the mathematics teacher wants to test the whole class on fractions and obtain the marks immediately; (d) new students must learn, at their own pace, how the parts of a computer work.
\tcblower
\textbf{Solution.}
\begin{enumerate}
\item \textbf{(a) Irregular verbs $\rightarrow$ drill-and-practice software.} The objective is memorisation through repetition. The software presents verbs, the student types the past tense and past participle, and receives immediate correction; errors are repeated until mastery. Pedagogical advantage: instant, individualised feedback impossible for one teacher with 60 students.
\item \textbf{(b) Combustion of sodium $\rightarrow$ simulation software.} The experiment is dangerous (violent reaction, toxic chlorine gas). A simulation lets students vary the quantities, observe the reaction and the products safely, and repeat it as many times as needed at zero material cost. Advantage: safety, cost and repeatability.
\item \textbf{(c) Test on fractions $\rightarrow$ assessment (testing) software.} The teacher prepares a quiz; each student answers on a computer; the software marks automatically and produces the class results instantly. Advantage: immediate, objective marking and statistics per question, allowing the teacher to detect which notions are not mastered.
\item \textbf{(d) Discovering the computer $\rightarrow$ tutorial software.} The objective is to learn new content individually. A tutorial presents the parts of the computer step by step (text, images, animations), asks check questions after each step, and adapts to the student's pace. Advantage: self-paced learning with built-in verification.
\end{enumerate}
\textbf{Examiner's remark:} each answer names the CAL family, links it to the objective, and gives a pedagogical justification — the three elements expected by the marking guide.
\end{exoresolu}

\courssub{Designing a CAL Lesson: The ADDIE-Inspired Approach}

\begin{methode}[Steps to design a CAL-supported lesson]
When asked to ``describe how a teacher would prepare a lesson using CAL software'':
\begin{enumerate}
\item \textbf{Analyse:} identify the syllabus objective, the learners' prior knowledge, and the available software/hardware.
\item \textbf{Select:} choose the CAL type (drill, tutorial, simulation, etc.) matching the objective.
\item \textbf{Prepare:} install/test the software; create student worksheets or guiding questions; plan grouping (individual, pairs).
\item \textbf{Implement:} (a) brief introduction by teacher (5 min); (b) student activity on software (20-30 min) with teacher circulating; (c) debrief / plenary discussion (10 min) to consolidate.
\item \textbf{Evaluate:} check software logs / quiz scores; ask students for feedback; note technical issues for next session.
\end{enumerate}
\textbf{Key point:} the teacher remains the orchestrator — CAL provides the \emph{individualised practice}, the teacher provides the \emph{meaning-making}.
\end{methode}

\begin{experience}[Guided demonstration: Simulation of an RC circuit]
\textbf{Context:} Upper Sixth Physics, chapter on capacitor charging/discharging. School has 15 computers with PhET simulations installed offline.
\textbf{Procedure:}
\begin{enumerate}
\item Teacher projects the PhET ``Circuit Construction Kit: DC'' simulation.
\item Students predict: ``What happens to the voltage across the capacitor if we double the resistance?''
\item In pairs, students use the simulation: build RC circuit, vary R and C, record voltage-time graphs.
\item Teacher leads discussion: link time constant $\tau = RC$ to observed curves.
\item Students complete a worksheet with three application questions.
\end{enumerate}
\textbf{Value:} invisible phenomenon (current, voltage) made visible; safe; repeatable; zero consumable cost.
\end{experience}

\coursec{E-Learning Models, Implementation and Evaluation in the Cameroonian Context}

\courssub{E-Learning Models: Definitions and Indicators}

\begin{definition}[E-Learning]
\cle{E-Learning} is learning facilitated by digital technologies \emph{delivered through a network} (Internet or intranet), where the learning environment (content, activities, communication, assessment) is hosted on a platform (LMS) and allows interaction between learners and instructors regardless of physical location.
\end{definition}

\begin{popfigure}
\begin{tikzpicture}[scale=0.9, every node/.style={font=\small}]
% Three models
\node[draw, rounded corners, fill=popblueL, minimum width=4cm, minimum height=2.5cm, align=center] (f2f) at (-6,0){\textbf{Face-to-Face (Traditional)}\\\textit{All activities in classroom}\\\textit{Teacher present}\\\textit{No network delivery of course}};
\node[draw, rounded corners, fill=popgreenL, minimum width=4cm, minimum height=2.5cm, align=center] (blended) at (0,0){\textbf{Blended / Hybrid Learning}\\\textit{Mix of F2F + Online}\\\textit{Flipped classroom, rotation}\\\textit{LMS + Classroom}};
\node[draw, rounded corners, fill=poporange, minimum width=4cm, minimum height=2.5cm, align=center] (online) at (6,0){\textbf{Fully Online (Distance)}\\\textit{All via LMS (Moodle)}\\\textit{Sync (Zoom) + Async (Videos, Forums)}\\\textit{Physical only for exams}};

% Spectrum arrow
\draw[<->, very thick, popdark] (-8,-2) -- (8,-2) node[midway, below, font=\small] {Increasing Online Proportion};

% Indicators
\node[text width=11cm, align=left, draw, rounded corners, fill=poppink, minimum width=12cm, minimum height=1.5cm, font=\small] (indicators) at (0,-4) {
\textbf{Key Indicators in a Scenario:}
\begin{itemize}
\item ``Students watch videos at home, do exercises in class'' $\rightarrow$ \textbf{Blended (Flipped Classroom)}
\item ``All courses, tutorials, continuous assessment on platform; only final exam at centre'' $\rightarrow$ \textbf{Fully Online}
\item ``Teacher uses projector to show CD-ROM during lesson'' $\rightarrow$ \textbf{Face-to-Face + CAL (NOT e-learning)}
\item ``WhatsApp group for questions + PDF notes sent by email'' $\rightarrow$ \textbf{Informal supplement, not a structured e-learning model}
\end{itemize}
};
\end{tikzpicture}
\legende{E-Learning models spectrum and scenario indicators}
\end{popfigure}

\begin{methode}[Identifying the e-learning model of a training situation]
\begin{enumerate}
\item \textbf{Ask the key question:} where and when do learners and teacher interact?
\item \textbf{Classify:}
\begin{itemize}
\item \cle{Face-to-face (traditional)}: all activities in the classroom.
\item \cle{Fully online (distance) learning}: all activities through a platform (LMS such as Moodle), no physical meeting; may be \emph{synchronous} (live video class) or \emph{asynchronous} (recorded lessons, forums).
\item \cle{Blended (hybrid) learning}: a deliberate mix of face-to-face sessions and online activities.
\end{itemize}
\item \textbf{Look for indicators in the scenario:} ``students watch videos at home and do exercises in class'' $\rightarrow$ blended (flipped classroom); ``the course is entirely on Moodle with a final exam at a centre'' $\rightarrow$ fully online.
\item \textbf{To choose a model, weigh:} connectivity of learners, cost, need for practical work (labs require presence), autonomy of learners.
\item \textbf{Name the supporting tools:} LMS (Moodle, Google Classroom), videoconferencing (Zoom, Teams), content formats (PDF, video, SCORM quizzes), communication (forum, WhatsApp group as an informal complement).
\end{enumerate}
\textbf{Key idea:} e-learning $\neq$ putting a PDF online; true e-learning includes \emph{interaction, activities and assessment} through the platform.
\end{methode}

\begin{exoresolu}[Three training programmes]
Identify the learning model used in each case and justify: (a) The ICT teacher of GBHS Limbe posts recorded video lessons and quizzes on Moodle; students study them at home, and class time is used for practical work and questions. (b) A university offers a computer science degree where all courses, tutorials and continuous assessments are done on a platform; students only come to a centre in Yaoundé for the final examination. (c) In a village school without Internet, the teacher uses a projector to show an educational CD-ROM during the normal lesson. (d) For situation (b), state two advantages and two difficulties for a Cameroonian student.
\tcblower
\textbf{Solution.}
\begin{enumerate}
\item \textbf{(a) Blended learning, in the ``flipped classroom'' variant.} There is a deliberate combination of online activities (videos, quizzes on Moodle at home) and face-to-face sessions (practical work in class). The online part is asynchronous (recorded content).
\item \textbf{(b) Fully online (distance) learning.} All teaching, interaction and continuous assessment pass through the platform; the only physical presence is the final examination, which does not change the model. It mixes synchronous elements (live tutorials) and asynchronous ones (course materials, forums).
\item \textbf{(c) Face-to-face learning enriched with CAL, not e-learning.} The CD-ROM is a computer-assisted learning resource used \emph{inside} the classroom by the teacher; there is no network-based delivery of the course, so this is traditional teaching supported by ICT, not an e-learning model.
\item \textbf{(d) Advantages:} (i) flexibility — the student can study while keeping a job or staying in his/her region, saving relocation costs to Yaoundé; (ii) self-paced access — recorded lessons can be re-watched until understood. \textbf{Difficulties:} (i) connectivity and cost of data bundles, plus electricity cuts, which can prevent following live sessions; (ii) the need for strong self-discipline and digital skills — without a teacher physically present, isolated students may drop out.
\end{enumerate}
\textbf{Examiner's remark:} case (c) is the classic trap: using a computer in class is CAL, not e-learning. E-learning implies learning \emph{through a network-delivered environment}.
\end{exoresolu}

\courssub{LMS and Technical Infrastructure}

\begin{definition}[Learning Management System (LMS)]
An \cle{LMS} is a software platform that administers, delivers, tracks and reports on e-learning courses. Core functions: course creation, user enrolment, content hosting (SCORM, video, PDF), communication tools (forum, messaging), assessment (quizzes, assignments), gradebook, analytics.
\end{definition}

\begin{center}
\begin{tabularx}{\linewidth}{|l|X|X|}\hline
\rowcolor{popblueL}
\textbf{LMS} & \textbf{Licence / Cost} & \textbf{Suitability for Cameroon} \\ \hline
Moodle & Open source (free), self-hosted & \textbf{Best for schools:} runs on local server (no Internet needed), low hardware needs, large community, French/English \\ \hline
Google Classroom & Free (Google account) & Needs reliable Internet; good for urban schools with connectivity \\ \hline
Canvas / Blackboard & Commercial, expensive & Universities with budget and IT support \\ \hline
Kolibri & Open source, offline-first & Designed for low-resource / offline contexts; pre-loaded content \\ \hline
\end{tabularx}
\end{center}

\begin{savaistu}[Offline-first e-learning in Cameroon]
Many Cameroonian schools (especially rural) have computers but unstable/no Internet. The solution: \textbf{install Moodle on a local server} (a dedicated PC or Raspberry Pi) connected to the school LAN. Students access \texttt{http://moodle.local} from the computer room. Content (videos, PDFs, quizzes) is pre-loaded. For home revision, export packages as \textbf{SCORM} or provide PDF/MP4 on USB drives. This model is used in several GHS in the North-West and Far North regions.
\end{savaistu}

\courssub{Implementation Plan for E-Learning}

\begin{methode}[Writing an implementation plan for e-learning in a school]
Exam questions may ask you to ``propose a plan'' or ``describe the steps'' to introduce e-learning. Use this five-phase framework:
\begin{enumerate}
\item \textbf{Needs analysis:} identify learners (number, level, digital skills), objectives, existing infrastructure (computers, electricity, Internet bandwidth).
\item \textbf{Choice of platform and content:} select an LMS (e.g.\ Moodle, free and installable on a local server), prepare digitised content (notes, videos, quizzes) aligned with the MINESEC syllabus.
\item \textbf{Infrastructure and access:} equip a computer room or provide offline access (local server, downloaded materials); plan for low bandwidth (light files, downloadable videos).
\item \textbf{Training and accompaniment:} train teachers to create courses and learners to use the platform; appoint a referent (e.g.\ the ICT teacher).
\item \textbf{Evaluation and improvement:} track logins, quiz results and completion rates; collect feedback; correct weaknesses each term.
\end{enumerate}
\textbf{Reflexes:} always mention teacher training (the most common cause of failure) and an offline/low-bandwidth solution for the Cameroonian context.
\end{methode}

\begin{exoresolu}[E-learning plan for a rural high school]
GHS Mbengwi (North-West Region) has 30 computers, an unstable Internet connection and one ICT teacher. The principal wants to introduce e-learning for Upper Sixth sciences. Propose a realistic five-step implementation plan, indicating for each step one concrete action adapted to this context.
\tcblower
\textbf{Solution.}
\begin{enumerate}
\item \textbf{Needs analysis.} Survey the 120 Upper Sixth students: how many own a smartphone or have network coverage at home? Test the actual bandwidth of the school connection. Conclusion expected: the platform must work mainly \emph{inside} the school and offline.
\item \textbf{Platform and content.} Install \textbf{Moodle on a local server} in the computer room, accessible through the school network without Internet. The ICT teacher and subject teachers digitise the syllabus: PDF notes, short videos recorded with a smartphone, and self-correcting quizzes per chapter.
\item \textbf{Infrastructure and access.} Organise the 30 computers into a wired/wireless local network connected to the Moodle server; create a timetable so each Upper Sixth class has at least one weekly session in the room; burn content on CDs/flash drives for students who want to revise at home.
\item \textbf{Training.} Organise a two-day workshop: teachers learn to create a Moodle course, upload documents and build quizzes; students receive a one-period induction on logging in, following a course and submitting work. The ICT teacher is appointed platform administrator.
\item \textbf{Evaluation.} Each term, the administrator checks Moodle statistics (number of connections, quiz averages, chapters completed); a short student questionnaire identifies difficulties; the principal and teachers adjust content and timetable accordingly.
\end{enumerate}
\textbf{Examiner's remark:} the strength of this plan is its adaptation to constraints: a \emph{local} Moodle server solves the unstable Internet problem — a detail that distinguishes an excellent script.
\end{exoresolu}

\courssub{Evaluating CAL / E-Learning Effectiveness}

\begin{methode}[Evaluating a computer-assisted learning activity]
To ``evaluate the effectiveness'' of a CAL tool or e-learning course:
\begin{enumerate}
\item \textbf{Define criteria before judging:}
\begin{itemize}
\item \emph{Pedagogical:} are the objectives achieved? Compare results (pre-test vs post-test, or an experimental group vs a control group).
\item \emph{Technical:} does the tool run correctly on the available machines (speed, bugs, compatibility)?
\item \emph{Ergonomic:} is the interface clear, navigation simple, feedback immediate?
\item \emph{Motivational:} do students use it willingly and longer?
\item \emph{Economic:} cost of licences/equipment versus benefits.
\end{itemize}
\item \textbf{Collect evidence:} quiz statistics from the software/LMS, observation of students, questionnaires, comparison of marks.
\item \textbf{Interpret with caution:} an improvement may come from novelty or extra time, not only from the tool.
\item \textbf{Conclude with a decision:} keep, adapt (more training, better content) or abandon.
\end{enumerate}
\textbf{Key formula (gain):} $\text{gain} = \bar{x}_{\text{post}} - \bar{x}_{\text{pre}}$ where $\bar{x}$ is the class average; compare with the gain of a control group.
\end{methode}

\begin{exoresolu}[Did the software improve performance?]
A school introduces a drill-and-practice software for Upper Sixth mathematics. In Term 1, before using it, the class average on a fractions test (marked over 20) is $9.5$. After six weeks of weekly use, the same level test gives an average of $12.3$. A parallel class that did not use the software moved from $9.8$ to $10.4$ on the same tests. (a) Compute the gain of each class. (b) Can we attribute the difference to the software? Justify. (c) State two other criteria the school should check before deciding to buy the software for all classes.
\tcblower
\textbf{Solution.}

\textbf{(a) Gains:}
\begin{align*}
\text{Experimental class:}\quad & \text{gain}_E = 12.3 - 9.5 = 2.8 \text{ points.}\\
\text{Control class:}\quad & \text{gain}_C = 10.4 - 9.8 = 0.6 \text{ point.}
\end{align*}

\textbf{(b) Attribution to the software.} The experimental class gained $2.8$ points while the control class gained only $0.6$ point on the same tests. The \emph{additional gain} attributable to the software is approximately
\[ 2.8 - 0.6 = 2.2 \text{ points over 20.} \]
Since both classes had the same syllabus, the same period and equivalent initial levels ($9.5$ vs $9.8$), the most reasonable explanation of the difference is the use of the software: immediate feedback and repeated practice helped students master fractions. However, we must remain cautious: other factors (different teachers, motivation due to novelty) may have contributed. The conclusion is therefore: the software \emph{probably} contributed significantly, but a longer study would confirm it.

\textbf{(c) Two other criteria:}
\begin{itemize}
\item \textbf{Technical and economic criterion:} does the software run smoothly on the school's computers, and is the licence cost per student affordable compared with the measured benefit?
\item \textbf{Ergonomic and motivational criterion:} do students find the interface clear and remain engaged after the novelty period (checked through a questionnaire and usage statistics), and do teachers find it easy to integrate into their lessons?
\end{itemize}
\textbf{Examiner's remark:} the comparison with a control group is what makes the evaluation rigorous; a simple before/after difference would not prove anything by itself.
\end{exoresolu}

\courssub{Summary: Key Terms for the Examination}

\begin{center}
\begin{tabularx}{\linewidth}{|l|X|}\hline
\rowcolor{popblueL}
\textbf{Term} & \textbf{Definition / Expected Answer} \\ \hline
E-Commerce & Buying/selling goods/services online (B2C, B2B, C2C, B2G) \\ \hline
E-Health & ICT for healthcare: telemedicine, EHR, mHealth, health info systems \\ \hline
E-Government & Online public services: G2C, G2B, G2G, G2E \\ \hline
Mobile Money & Dominant payment rail in Cameroon (MTN MoMo, Orange Money) \\ \hline
ANTIC & National agency for ICT security, .cm domain, certification \\ \hline
CAL & Computer-Assisted Learning: software for learning (drill, tutorial, simulation, game, assessment) \\ \hline
E-Learning & Network-delivered learning via LMS (interaction, activities, assessment) \\ \hline
Blended Learning & Deliberate mix of face-to-face and online (e.g. flipped classroom) \\ \hline
LMS & Learning Management System (Moodle, Google Classroom, Kolibri) \\ \hline
SCORM & Sharable Content Object Reference Model: packaging standard for e-learning content \\ \hline
4P Grid & People, Purse, Process, Place — framework for benefits/challenges analysis \\ \hline
Gain Score & $\bar{x}_{\text{post}} - \bar{x}_{\text{pre}}$; compare experimental vs control group \\ \hline
\end{tabularx}
\end{center}

\begin{exercice}[Practice Questions for GCE Preparation]
\begin{enumerate}
\item A farmer in Dschang uses his phone to: (i) check cocoa prices on a website; (ii) order fertiliser from an agro-dealer's WhatsApp catalogue and pays by Mobile Money; (iii) consult a veterinary officer via video call for his sick cow; (iv) declare his farm revenue on the DGI portal. Classify each as e-commerce, e-health or e-government with sub-category.
\item The Ministry of Public Health wants to deploy a telemedicine platform linking district hospitals to Yaoundé central hospitals. Using the 4P grid, give three benefits and three challenges with a mitigation for each.
\item Describe the seven steps of an online purchase on a Cameroonian e-commerce site using Mobile Money. At which steps are the risks of (a) phishing, (b) SIM-swap fraud highest? State one precaution per risk.
\item A teacher wants students to discover the structure of the heart. The school has no laboratory specimens. Which CAL type is most appropriate? Justify with two pedagogical arguments.
\item A scenario: ``Students download PDF notes from the school website, watch YouTube videos at home, and attend class for discussions.'' Is this e-learning? If yes, which model? If no, why?
\item Propose a 5-step e-learning implementation plan for a school with 40 computers, no Internet, and two ICT teachers. Mention one concrete action per step.
\item An e-learning pilot shows: experimental class gain = 3.1 points, control class gain = 0.9 points. Compute the additional gain. Can we conclude the platform caused the improvement? What other data would strengthen the claim?
\end{enumerate}
\end{exercice}

\begin{corrige}[Exercise 1]
\begin{enumerate}
\item (i) Information search — not a transaction, so not strictly e-commerce (but pre-purchase). (ii) \textbf{E-commerce B2C} (farmer buys from business). (iii) \textbf{E-health} (telemedicine / veterinary teleconsultation). (iv) \textbf{E-government G2C} (citizen declares revenue to administration).
\end{enumerate}
\end{corrige}

\begin{corrige}[Exercise 3]
Steps: 1 Browse, 2 Cart, 3 Login, 4 Confirm, 5 Payment (MoMo), 6 Fulfilment, 7 After-sales.
Risks: (a) Phishing highest at Step 3 (login) — precaution: verify HTTPS/padlock, never click unsolicited links. (b) SIM-swap highest at Step 5 (payment) — precaution: never share PIN/OTP, enable SIM lock with operator.
\end{corrige}

\begin{corrige}[Exercise 4]
\textbf{Simulation software.} Arguments: (1) Safety/ethics — no need for real specimens; (2) Interactivity — students can rotate, zoom, label, animate blood flow; repeatable at zero marginal cost.
\end{corrige}

\begin{corrige}[Exercise 5]
This is \textbf{Blended Learning (Flipped Classroom)}: online content delivery (videos, PDFs) at home + face-to-face interactive session in class. The website/YouTube act as the LMS/content repository.
\end{corrige}

\begin{corrige}[Exercise 7]
Additional gain = $3.1 - 0.9 = 2.2$ points. The difference suggests the platform contributed, but other factors (teacher effect, novelty, extra time) could play a role. Strengthen with: larger sample, random assignment, delayed post-test, qualitative feedback, usage logs correlation.
\end{corrige}

\newpage

\coursec{Exercises}

\courssub{I check my knowledge}

\begin{exercice}[MCQ: E‑service categories]
Which of the following best describes an e‑government service?
\begin{itemize}
\item A) Online marketplace for consumer goods
\item B) Digital platform for filing tax returns
\item C) Telemedicine consultation portal
\item D) Mobile money transfer application
\end{itemize}
\end{exercice}

\begin{exercice}[True/False with justification]
State whether each statement is true or false and justify your answer.
\begin{enumerate}
\item E‑commerce platforms do not require secure payment gateways.
\item E‑health services can reduce the need for physical hospital visits.
\item The digital divide has no impact on the adoption of e‑learning in rural areas.
\item SCORM packages are only used for video streaming.
\end{enumerate}
\end{exercice}

\begin{exercice}[Definitions]
Define the following terms in the context of electronic services:
\begin{enumerate}
\item E‑commerce
\item E‑health
\item E‑government
\item Computer‑Assisted Learning (CAL)
\end{enumerate}
\end{exercice}

\begin{exercice}[Direct calculation: Correlation]
The table below shows Internet penetration (\% of population) and e‑service adoption index (0–100) for five Cameroonian regions.

\begin{center}
\begin{tabular}{|c|c|c|}\hline
Region & Internet penetration (\%) & Adoption index \\ \hline
Centre & 45 & 68 \\ \hline
Littoral & 52 & 74 \\ \hline
West & 38 & 55 \\ \hline
North & 22 & 31 \\ \hline
Far North & 18 & 27 \\ \hline
\end{tabular}
\end{center}

Calculate the Pearson correlation coefficient between Internet penetration and adoption index. (Use the formula $r = \frac{\sum (x_i-\bar{x})(y_i-\bar{y})}{\sqrt{\sum (x_i-\bar{x})^2 \sum (y_i-\bar{y})^2}}$.)
\end{exercice}

\courssub{I practise}

\begin{exercice}[Architecture comparison diagram]
Draw a labelled block diagram comparing the typical architecture of an e‑commerce site (e.g. Jumia) with that of an e‑government portal (e.g. DGI). Indicate at least three common layers and two specific components for each system.
\end{exercice}

\begin{exercice}[Security measures]
For each of the two systems in the previous exercise, explain two security measures that must be implemented to protect user data and transactions. Justify why each measure is critical.
\end{exercice}

\begin{exercice}[SCORM package outline]
Design a SCORM package outline for a CAL lesson on “Spreadsheet Modelling”. Include:
\begin{enumerate}
\item Learning objectives (minimum three)
\item Content chunks (minimum four)
\item Interaction types (e.g. drag‑and‑drop, simulation)
\item Assessment items (minimum two, with question type)
\end{enumerate}
\end{exercice}

\begin{exercice}[MOOC selection criteria]
Using at least three selection criteria from the syllabus (e.g. accessibility, curriculum alignment, assessment quality, offline support), critically assess the suitability of a popular MOOC platform (e.g. Coursera, edX) for Upper Sixth ICT revision in Cameroon.
\end{exercice}

\courssub{I prepare for the GCE}

\begin{exercice}[Essay: Digital divide and e‑learning effectiveness (10 marks)]
Write a short essay (approximately 300 words) on how the digital divide affects the effectiveness of e‑learning in rural Cameroon. Propose two concrete mitigation strategies and discuss their feasibility.
\end{exercice}

\begin{exercice}[Case study: E‑government portal evaluation (12 marks)]
The Directorate General of Taxes (DGI) plans to launch a new e‑filing portal. As an ICT consultant, you are asked to:
\begin{enumerate}
\item Identify four functional requirements and two non‑functional requirements for the portal. (4 marks)
\item Propose a testing strategy covering security, usability, and performance. (4 marks)
\item Explain how you would ensure the portal complies with Cameroon’s data protection regulations (ANTIC guidelines). (4 marks)
\end{enumerate}
\end{exercice}

\begin{exercice}[Project design: CAL module for “Database Design” (13 marks)]
Design a complete Computer‑Assisted Learning module for the topic “Database Design” targeting Upper Sixth students. Your design must include:
\begin{enumerate}
\item A pedagogical scenario (context, target learners, prerequisites). (2 marks)
\item Detailed instructional design using the ADDIE model (Analysis, Design, Development, Implementation, Evaluation). (6 marks)
\item A plan for formative and summative assessment, including at least one interactive simulation. (3 marks)
\item Considerations for deployment in schools with limited Internet connectivity. (2 marks)
\end{enumerate}
\end{exercice}

\newpage

\coursec{Solutions to the exercises}

\begin{corrige}[Exercise 1 — MCQ: E‑service categories]
\textbf{Correct answer: B) Digital platform for filing tax returns.}

\textbf{Justification.} An \cle{e‑government} service is any public administration service delivered to citizens, businesses or other administrations through digital channels. Filing tax returns online (as with the DGI e‑filing portal in Cameroon) is a classic \textbf{G2C/G2B} (government‑to‑citizen / government‑to‑business) service.

The other options belong to different categories:
\begin{itemize}
\item A) An online marketplace for consumer goods is \cle{e‑commerce} (B2C).
\item C) A telemedicine consultation portal is \cle{e‑health}.
\item D) A mobile money transfer application (e.g. MTN MoMo, Orange Money) is a financial/telecom digital service, generally classified under e‑commerce/fintech, not e‑government.
\end{itemize}

\textbf{Examiner's tip.} In MCQs, always ask: \emph{who provides the service?} If it is a public administration (ministry, council, DGI, police), it is e‑government; if it is a private company selling goods, it is e‑commerce; if it concerns health care delivery, it is e‑health.
\end{corrige}

\begin{corrige}[Exercise 2 — True/False with justification]
\begin{enumerate}
\item \textbf{FALSE.} E‑commerce platforms \emph{must} use secure payment gateways (with encryption such as SSL/TLS, and protocols like 3‑D Secure). Without them, customers' card or mobile money credentials could be intercepted, and no buyer would trust the platform. Security is a \emph{prerequisite}, not an option.

\item \textbf{TRUE.} E‑health services such as telemedicine, remote diagnosis and electronic appointment booking allow patients — especially in remote areas like parts of the Far North — to consult a doctor or obtain advice without travelling to a hospital. Physical visits remain necessary for examinations and surgery, but their number can be reduced.

\item \textbf{FALSE.} The \cle{digital divide} (unequal access to devices, electricity and Internet connectivity) is one of the \emph{main obstacles} to e‑learning adoption in rural Cameroon. Where learners lack smartphones, computers, electricity or affordable data, online learning simply cannot reach them, which widens educational inequality.

\item \textbf{FALSE.} \cle{SCORM} (Sharable Content Object Reference Model) is a \emph{packaging and interoperability standard} for e‑learning content. A SCORM package can contain text, images, quizzes, simulations, drag‑and‑drop activities \emph{and} videos, together with tracking data (scores, completion status) exchanged with an LMS. It is not limited to video streaming.
\end{enumerate}
\end{corrige}

\begin{corrige}[Exercise 3 — Definitions]
\begin{enumerate}
\item \textbf{E‑commerce} (electronic commerce): the buying and selling of goods and services, and the transfer of money or data, over electronic networks — mainly the Internet. It includes B2C (e.g. Jumia), B2B, C2C and increasingly mobile commerce via mobile money.

\item \textbf{E‑health}: the use of information and communication technologies to support health services and health management — for example telemedicine, electronic patient records, health information systems, appointment booking and health education platforms.

\item \textbf{E‑government}: the use of ICT, especially the Internet, by public administrations to deliver information and services to citizens, businesses and other government bodies (G2C, G2B, G2G), e.g. online tax filing, e‑visa applications, digital civil status records.

\item \textbf{Computer‑Assisted Learning (CAL)}: the use of computers and digital resources (tutorials, simulations, drill‑and‑practice software, interactive exercises) to support, supplement or deliver teaching and learning, usually with immediate feedback and learner‑paced progression. It is the broader concept of which e‑learning (network/Internet‑based delivery) is a particular form.
\end{enumerate}
\end{corrige}

\begin{corrige}[Exercise 4 — Direct calculation: Correlation]
\textbf{Step 1 — Means.}
\begin{align*}
\bar{x} &= \frac{45+52+38+22+18}{5} = \frac{175}{5} = 35, &
\bar{y} &= \frac{68+74+55+31+27}{5} = \frac{255}{5} = 51.
\end{align*}

\textbf{Step 2 — Table of deviations.}
\begin{center}
\begin{tabular}{|c|c|c|c|c|c|c|}\hline
$x_i$ & $y_i$ & $x_i-\bar{x}$ & $y_i-\bar{y}$ & $(x_i-\bar{x})(y_i-\bar{y})$ & $(x_i-\bar{x})^2$ & $(y_i-\bar{y})^2$ \\ \hline
45 & 68 & $10$ & $17$ & $170$ & $100$ & $289$ \\ \hline
52 & 74 & $17$ & $23$ & $391$ & $289$ & $529$ \\ \hline
38 & 55 & $3$ & $4$ & $12$ & $9$ & $16$ \\ \hline
22 & 31 & $-13$ & $-20$ & $260$ & $169$ & $400$ \\ \hline
18 & 27 & $-17$ & $-24$ & $408$ & $289$ & $576$ \\ \hline
\multicolumn{2}{|c|}{Sums} & & & $1241$ & $856$ & $1810$ \\ \hline
\end{tabular}
\end{center}

\textbf{Step 3 — Apply the formula.}
\begin{align*}
r &= \frac{\sum (x_i-\bar{x})(y_i-\bar{y})}{\sqrt{\sum (x_i-\bar{x})^2 \sum (y_i-\bar{y})^2}}
= \frac{1241}{\sqrt{856 \times 1810}} \\
&= \frac{1241}{\sqrt{1\,549\,360}} \approx \frac{1241}{1244.73} \approx 0.997.
\end{align*}

\textbf{Conclusion.} $r \approx 0.997$, a \textbf{very strong positive linear correlation}: regions with higher Internet penetration show much higher e‑service adoption. This supports the idea that connectivity is a key driver of e‑service uptake (though correlation alone does not prove causation — income and education levels also play a role).
\end{corrige}

\begin{corrige}[Exercise 5 — Architecture comparison diagram]
A correct answer presents the two systems side by side, with three shared layers (presentation, application/logic, data) and specific components for each.

\begin{popfigure}
\begin{tikzpicture}[font=\small, node distance=0.35cm]
\tikzset{layer/.style={draw=popline, rounded corners=2pt, minimum width=5.6cm, minimum height=0.85cm, align=center, fill=popblueL}}
\tikzset{spec/.style={draw=popline, rounded corners=2pt, minimum width=5.6cm, minimum height=0.85cm, align=center, fill=popgreenL}}
% E-commerce column
\node[font=\bfseries, popdark] (t1) at (0,0) {E‑commerce site (Jumia)};
\node[layer, below=of t1, align=center] (p1){Presentation layer\\ web storefront, mobile app};
\node[layer, below=of p1, align=center] (a1){Application layer\\ catalogue, cart, order engine};
\node[layer, below=of a1, align=center] (d1){Data layer\\ product \& customer databases};
\node[spec, below=of d1, align=center] (s1){Specific: payment gateway\\ (MoMo, cards), logistics/tracking};
\node[spec, below=of s1, align=center] (s2){Specific: recommendation \&\\ promotions engine};
% E-government column
\node[font=\bfseries, popdark] (t2) at (7.6,0) {E‑government portal (DGI)};
\node[layer, below=of t2, align=center] (p2){Presentation layer\\ citizen web portal, forms};
\node[layer, below=of p2, align=center] (a2){Application layer\\ tax computation, workflow};
\node[layer, below=of a2, align=center] (d2){Data layer\\ taxpayer registry, filings DB};
\node[spec, below=of d2, align=center] (s3){Specific: digital identity \&\\ authentication (unique taxpayer ID)};
\node[spec, below=of s3, align=center] (s4){Specific: inter‑ministry\\ interoperability gateway (G2G)};
% arrows
\foreach \u/\v in {p1/a1, a1/d1, p2/a2, a2/d2} {\draw[<->, popmuted] (\u) -- (\v);}
\end{tikzpicture}
\legende{Comparison of typical e‑commerce and e‑government architectures.}
\end{popfigure}

\textbf{Three common layers:} (1) \textbf{presentation layer} (user interface: web pages, mobile app); (2) \textbf{application/business logic layer} (processing rules: orders, tax computation); (3) \textbf{data layer} (databases storing products/customers or taxpayers/filings).

\textbf{Two specific components each:} e‑commerce — secure \emph{payment gateway} (mobile money, bank cards) and \emph{order/logistics tracking}; e‑government — \emph{strong digital identity/authentication} (unique taxpayer number, possibly digital certificates) and an \emph{interoperability gateway} exchanging data with other administrations (G2G).

\textbf{Marking guide:} 1 mark per correctly drawn and labelled common layer (max 3), 1 mark per valid specific component (max 4), 1 mark for clarity/labels and correct flow arrows.
\end{corrige}

\begin{corrige}[Exercise 6 — Security measures]
\textbf{E‑commerce site (e.g. Jumia):}
\begin{enumerate}
\item \textbf{SSL/TLS encryption (HTTPS) for all transactions.} Encrypts data in transit so that card numbers, mobile money PINs and passwords cannot be intercepted by attackers on the network. Critical because payment data is the prime target of e‑commerce fraud.
\item \textbf{Strong authentication and access control} (passwords + one‑time codes/2FA, hashed password storage). Prevents account takeover, fraudulent orders and identity theft; critical because a compromised customer account directly causes financial loss and destroys trust.
\end{enumerate}

\textbf{E‑government portal (e.g. DGI):}
\begin{enumerate}
\item \textbf{Digital identity verification and role‑based access control} (unique taxpayer ID, certificates, 2FA for agents). Guarantees that only the legitimate taxpayer can view or modify a tax file and that civil servants only access data needed for their role — critical because tax records are highly confidential and legally protected.
\item \textbf{Audit trails, logging and data integrity controls} (digital signatures on filings, tamper‑evident logs). Every access and modification is recorded, enabling fraud detection, dispute resolution and accountability — critical because falsified or untraceable tax records would undermine state revenue and citizens' trust.
\end{enumerate}

\textbf{Justification principle (to quote in the exam):} a security measure is \emph{critical} when its absence leads to a concrete, severe harm — financial theft, privacy breach, data falsification or loss of public trust — that the system exists precisely to prevent.
\end{corrige}

\begin{corrige}[Exercise 7 — SCORM package outline]
\textbf{Package title:} \emph{Spreadsheet Modelling} — Upper Sixth ICT (SCORM 1.2 package, to run on any LMS such as Moodle).

\textbf{1. Learning objectives} (at the end, the learner will be able to):
\begin{enumerate}
\item Build a spreadsheet model using cell references, formulas and functions (SUM, AVERAGE, IF).
\item Use absolute and relative references correctly when copying formulas.
\item Create and interpret a chart from modelled data and perform a simple what‑if analysis.
\end{enumerate}

\textbf{2. Content chunks (SCOs):}
\begin{enumerate}
\item SCO 1 — What is a model? From a real situation (school canteen budget) to a spreadsheet.
\item SCO 2 — Formulas, functions and cell references (relative vs absolute).
\item SCO 3 — Charts and data presentation.
\item SCO 4 — What‑if analysis and goal seeking (changing inputs, observing outputs).
\item SCO 5 — Summary, glossary and final quiz.
\end{enumerate}

\textbf{3. Interaction types:}
\begin{itemize}
\item \textbf{Drag‑and‑drop:} match each function (SUM, IF, AVERAGE) to its purpose.
\item \textbf{Simulation:} an embedded mini‑spreadsheet where the learner types a formula and sees the result instantly; a what‑if slider changing the price of an item and updating total cost.
\item \textbf{Fill‑in‑the‑blank:} complete the formula \texttt{=B2*\$C\$1} and explain the role of the \$ signs.
\end{itemize}

\textbf{4. Assessment items:}
\begin{enumerate}
\item \textbf{Multiple‑choice question} (auto‑marked): "Which reference stays fixed when a formula is copied? A) A1  B) \$A\$1  C) A\$  D) 1A" — answer B.
\item \textbf{Practical/simulation task} (tracked by score): given a marks table, write a formula computing the class average; the SCO checks the entered formula and reports the score to the LMS.
\end{enumerate}

\textbf{Packaging:} all SCOs are zipped with an \texttt{imsmanifest.xml} file describing the structure, so the package can be imported into any SCORM‑compliant LMS and track completion and scores.
\end{corrige}

\begin{corrige}[Exercise 8 — MOOC selection criteria]
\textbf{Platform assessed:} Coursera (similar reasoning applies to edX) for Upper Sixth ICT revision in Cameroon.

\textbf{Criterion 1 — Accessibility (cost and connectivity).} Courses can be audited free of charge, but certificates are paid, and video streaming consumes large volumes of mobile data, which is expensive for many Cameroonian families. Mobile apps exist, but smooth streaming requires a stable 3G/4G connection, not guaranteed in rural areas. \emph{Verdict: partially suitable.}

\textbf{Criterion 2 — Curriculum alignment.} Coursera content follows university/industry syllabi (mostly American/European), not the MINESEC Upper Sixth ICT programme. Topics like databases, spreadsheets and networking are covered, but the GCE exam format, notation and progression differ. \emph{Verdict: useful as a supplement, not as the main revision tool.}

\textbf{Criterion 3 — Assessment quality.} Auto‑marked quizzes, peer assignments and certificates provide good formative feedback and motivation. However, assessments rarely match GCE question styles (structured essays, practical papers). \emph{Verdict: good for self‑testing concepts, weak for exam technique.}

\textbf{Criterion 4 — Offline support.} The mobile app allows downloading videos for offline viewing, which mitigates connectivity problems; but interactive quizzes and submissions still require Internet. \emph{Verdict: moderate.}

\textbf{Conclusion.} Coursera is a valuable \emph{complementary} resource for motivated Upper Sixth students with smartphones and data access, especially to deepen understanding. But because of cost, curriculum mismatch and connectivity limits, it cannot replace teacher‑guided, syllabus‑aligned revision (past GCE papers, locally designed CAL modules). A critical, criteria‑based judgement like this — with a balanced conclusion — is what earns full marks.
\end{corrige}

\begin{corrige}[Exercise 9 — Essay: Digital divide and e‑learning effectiveness (10 marks)]
\textbf{Indicative mark scheme:} definition and explanation of the digital divide (2 marks); at least three well‑developed effects on e‑learning effectiveness (3 marks); two concrete mitigation strategies with feasibility discussion (4 marks); structure, coherence and language (1 mark).

\textbf{Model essay (approximately 300 words).}

The \emph{digital divide} is the gap between those who have access to digital technologies — devices, electricity, affordable Internet — and those who do not. In Cameroon this divide is strongly geographical: while towns like Yaoundé and Douala enjoy 4G coverage, many rural areas lack reliable electricity and connectivity.

This divide directly undermines the effectiveness of e‑learning in rural Cameroon. First, without smartphones or computers, learners simply cannot access online lessons, videos or quizzes. Second, where mobile data is expensive, families must choose between connectivity and basic needs, so learning platforms are used irregularly. Third, poor electricity supply interrupts study and damages equipment. Finally, rural teachers often receive less ICT training, so even available tools are underexploited. The result is that e‑learning, instead of reducing inequality, risks \emph{widening} the achievement gap between urban and rural candidates at the GCE.

Two mitigation strategies can be proposed. The first is the deployment of \textbf{offline‑first learning solutions}: CAL packages (e.g. SCORM modules on local servers or Raspberry Pi devices) and downloadable content that work without Internet, synchronising only occasionally. This is highly feasible because it requires modest equipment, works with solar power, and content can be aligned with the MINESEC syllabus; its main challenge is the initial funding and teacher training. The second strategy is \textbf{community digital learning centres} — school or council computer rooms with shared devices and subsidised connectivity, possibly supported by universal‑access programmes and operators. Feasibility is good in towns and large villages, but maintenance costs, security of equipment and staffing must be planned for sustainability.

In conclusion, e‑learning can only fulfil its promise in rural Cameroon if access barriers are tackled first. Offline content and shared community centres are realistic, complementary steps towards making digital learning truly inclusive.

\textbf{Examiner expectations:} a clear introduction defining the divide, a developed body linking causes to effects on learning, concrete (not vague) strategies with honest feasibility discussion, and a short conclusion.
\end{corrige}

\begin{corrige}[Exercise 10 — Case study: E‑government portal evaluation (12 marks)]
\textbf{1. Requirements (4 marks — 1 mark per valid requirement).}

\emph{Functional requirements (what the system must do):}
\begin{enumerate}
\item Allow a taxpayer to create an account and authenticate securely (unique taxpayer number + password/OTP).
\item Allow online completion, validation and submission of tax declaration forms, with automatic tax computation.
\item Allow online payment of taxes (mobile money, bank card) and issue an electronic receipt.
\item Allow the taxpayer to track the status of a file and download past declarations; allow DGI agents to process and validate filings.
\end{enumerate}

\emph{Non‑functional requirements (qualities of the system):}
\begin{enumerate}
\item \textbf{Security/confidentiality:} encryption of data in transit and at rest; compliance with data protection rules.
\item \textbf{Performance/availability:} the portal must handle peak loads (declaration deadlines) with response times of a few seconds and high availability (e.g. 99\% uptime); also usability (interface in English and French, accessible to non‑experts) is acceptable.
\end{enumerate}

\textbf{2. Testing strategy (4 marks).}
\begin{itemize}
\item \textbf{Security testing:} penetration testing, vulnerability scans, tests of authentication, session management and resistance to common attacks (SQL injection, XSS); verification that data is encrypted.
\item \textbf{Usability testing:} test sessions with a sample of real taxpayers (including low digital literacy users) performing key tasks (register, file, pay); collection of feedback on clarity, bilingual interface, error messages.
\item \textbf{Performance testing:} load and stress testing simulating thousands of simultaneous users at deadline peaks; measurement of response time and system behaviour under failure.
\item Plus functional/unit and acceptance testing (UAT) with DGI before launch.
\end{itemize}

\textbf{3. Compliance with data protection (ANTIC guidelines) (4 marks).}
\begin{itemize}
\item Collect only the data strictly necessary (data minimisation) and inform users of the purpose of processing (privacy notice, consent).
\item Implement technical and organisational security measures: encryption, access control, audit logs, secure hosting — as promoted by ANTIC, the national agency for ICT security and trust.
\item Guarantee citizens' rights over their data: access, rectification, and defined retention periods.
\item Register/declare the processing with the competent authorities, appoint a data protection officer, and ensure any subcontractors (hosting, payment) respect the same obligations.
\end{itemize}

\textbf{Examiner expectations:} precise, testable requirements (not vague words like "good design"); a testing strategy explicitly covering the three requested dimensions; and concrete, Cameroon‑relevant compliance measures citing ANTIC's role in security and trust of electronic services.
\end{corrige}

\begin{corrige}[Exercise 11 — Project design: CAL module for "Database Design" (13 marks)]
\textbf{1. Pedagogical scenario (2 marks).}
\emph{Context:} a school library in Bamenda wants to computerise its loans. \emph{Target learners:} Upper Sixth ICT students preparing the GCE A‑Level. \emph{Prerequisites:} basic file organisation, spreadsheets, notion of tables and records. \emph{Duration:} 6 one‑hour sessions, blending computer room practice and self‑paced study.

\textbf{2. Instructional design with ADDIE (6 marks).}
\begin{itemize}
\item \textbf{Analysis:} learners' needs (master entities, attributes, keys, relationships, normalisation basics); constraints (limited connectivity, shared computers); syllabus objectives of the Databases chapter.
\item \textbf{Design:} module structured in five units — (U1) from paper records to tables; (U2) entities, attributes, primary keys; (U3) relationships and E‑R diagrams; (U4) building the database in a DBMS; (U5) simple queries. Each unit has objectives, content, activities and a quiz; storyboard of screens; SCORM packaging for tracking.
\item \textbf{Development:} production of interactive content: annotated diagrams, a drag‑and‑drop E‑R builder, a step‑by‑step simulation of table creation, short videos with transcripts, self‑correcting quizzes; integration into the school LMS (e.g. Moodle).
\item \textbf{Implementation:} pilot with one class; teacher trained as facilitator; sessions alternate guided discovery and individual work on the module; technical support plan.
\item \textbf{Evaluation:} learner quiz results and tracking data analysed; satisfaction questionnaire; comparison of pilot class results with previous cohorts; revision of weak units (formative evaluation of the module itself).
\end{itemize}

\textbf{3. Assessment plan (3 marks).}
\emph{Formative:} auto‑marked quizzes after each unit with immediate feedback; an \textbf{interactive simulation} in which the learner drags entities (BOOK, STUDENT, LOAN), adds attributes, chooses primary keys and draws relationship cardinalities — the system validates the diagram and explains errors. \emph{Summative:} a practical project — design and implement the complete library database (tables, keys, relationships, three queries) in a DBMS, marked with a grid (correct keys, integrity, queries), plus a written test in GCE format.

\textbf{4. Deployment with limited connectivity (2 marks).}
\begin{itemize}
\item Provide the module as a downloadable/offline SCORM package on a local school server or USB drives; videos compressed and optional (transcripts available).
\item Use free, lightweight, offline‑capable software; synchronise scores only when a connection is available; plan low‑bandwidth alternatives (printed worksheets mirroring the interactive tasks).
\end{itemize}

\textbf{Indicative mark scheme:} scenario 2 marks (context + learners + prerequisites); ADDIE 6 marks (about 1 mark per phase, requiring concrete, database‑specific content, not generic descriptions); assessment 3 marks (formative + summative + a clearly described simulation); connectivity 2 marks (at least two realistic offline strategies). \textbf{Examiner expectations:} coherence between objectives, activities and assessment, and realistic adaptation to the Cameroonian school context.
\end{corrige}

\newpage

\coursec{Summary sheet}

\begin{popfigure}
\begin{tikzpicture}[
  mindmap,
  grow cyclic,
  every node/.style={concept, font=\small, align=center},
  root concept/.append style={concept color=popblue, font=\bfseries\small, minimum size=3.2cm},
  level 1 concept/.append style={level distance=4.6cm, sibling angle=60, font=\scriptsize},
  level 2 concept/.append style={level distance=3.0cm, font=\tiny},
  concept connection/.append style={color=popmuted}
]
\node[root concept]{Learning with the Assistance of Computers}
  child[concept color=poporange]{ node {E-Commerce}
    child[concept color=poporange!70]{ node {B2B, B2C, C2C, C2B} }
    child[concept color=poporange!70]{ node {Mobile money: MTN MoMo, Orange Money} }
  }
  child[concept color=popgreen]{ node {E-Health}
    child[concept color=popgreen!70]{ node {Telemedicine, EMR} }
    child[concept color=popgreen!70]{ node {Health info systems} }
  }
  child[concept color=poppurple]{ node {E-Government}
    child[concept color=poppurple!70]{ node {G2C, G2B, G2G} }
    child[concept color=poppurple!70]{ node {ANTIC, online tax, e-services} }
  }
  child[concept color=poppink]{ node {Benefits \& Challenges}
    child[concept color=poppink!70]{ node {Speed, cost, access 24/7} }
    child[concept color=poppink!70]{ node {Digital divide, security, literacy} }
  }
  child[concept color=popgold]{ node {CAL}
    child[concept color=popgold!70]{ node {Tutorials, drills, simulations} }
    child[concept color=popgold!70]{ node {CBT, CAI, CMI} }
  }
  child[concept color=popink]{ node {E-Learning}
    child[concept color=popink!70]{ node {LMS: Moodle, Google Classroom} }
    child[concept color=popink!70]{ node {Synchronous vs asynchronous} }
  };
\end{tikzpicture}
\legende{Mind map of the chapter: electronic services and computer-assisted learning.}
\end{popfigure}

\begin{bilan}
\textbf{1. Key definitions}
\begin{itemize}
  \item \cle{Electronic service (e-service)}: any service delivered to users through ICT networks (Internet, mobile networks) instead of face-to-face contact.
  \item \cle{E-commerce}: buying and selling of goods and services, and the transfer of money or data, over an electronic network. Main models: \textbf{B2C} (business to consumer, e.g.\ Jumia), \textbf{B2B} (business to business), \textbf{C2C} (consumer to consumer, e.g.\ online marketplaces), \textbf{C2B} (consumer to business).
  \item \cle{E-health}: the use of ICT to support health services and information: \textbf{telemedicine} (remote consultation), \textbf{electronic medical records (EMR)}, hospital information systems, health education by SMS.
  \item \cle{E-government}: the use of ICT by public administrations to deliver services to citizens (\textbf{G2C}), businesses (\textbf{G2B}) and other administrations (\textbf{G2G}): online tax declaration, e-visa, digital birth certificates, ANTIC services.
  \item \cle{Computer-Assisted Learning (CAL)}: use of the computer as a teaching/learning aid. Related terms: \textbf{CAI} (computer-assisted instruction), \textbf{CBT} (computer-based training), \textbf{CMI} (computer-managed instruction).
  \item \cle{E-learning}: learning delivered wholly or partly through electronic technologies, often via the Internet, using an \textbf{LMS} (Learning Management System) such as Moodle or Google Classroom.
\end{itemize}

\textbf{2. Facts and classifications to know}
\begin{itemize}
  \item CAL software types: \textbf{tutorials} (present content step by step), \textbf{drill and practice} (repetition exercises), \textbf{simulations} (model real situations), \textbf{educational games}, \textbf{problem-solving software}.
  \item E-learning modes: \textbf{synchronous} (live: video-conference, chat) vs \textbf{asynchronous} (any time: recorded lessons, forums, email); \textbf{blended learning} mixes face-to-face and online.
  \item Roles of the computer in CAL: \textbf{tutor} (teaches), \textbf{tool} (word processor, spreadsheet), \textbf{tutee} (the learner programs the computer).
  \item Benefits of e-services: availability 24/7, lower transaction costs, speed, wider reach, transparency, reduced travel (important where distances are long, e.g.\ Far North to Yaound\'e).
  \item Challenges: \textbf{digital divide} (urban/rural, rich/poor), cost of devices and data, unreliable electricity and connectivity, low digital literacy, \textbf{cybersecurity and privacy} risks, fraud in online payments, resistance to change.
\end{itemize}

\textbf{3. Methods to master}
\begin{itemize}
  \item \textbf{Analysing an e-service}: identify the service $\rightarrow$ the technology used $\rightarrow$ the actors $\rightarrow$ the benefits $\rightarrow$ the challenges $\rightarrow$ possible solutions (e.g.\ community telecentres, solar power, training).
  \item \textbf{Designing a simple CAL lesson}: define objectives $\rightarrow$ choose software type (tutorial, drill, simulation) $\rightarrow$ structure content in small units $\rightarrow$ add feedback and assessment $\rightarrow$ test with learners $\rightarrow$ evaluate and improve.
  \item \textbf{Evaluating e-learning}: check achievement of objectives, learner engagement, quality of feedback, technical reliability, accessibility for all learners.
\end{itemize}

\textbf{4. Common mistakes to avoid}
\begin{itemize}
  \item Confusing \textbf{e-commerce} (trade) with \textbf{e-business} (all digital business processes, broader).
  \item Confusing \textbf{CAL} (computer assists learning, any setting) with \textbf{e-learning} (learning delivered through networks); e-learning is a form of CAL but not the only one.
  \item Mixing up \textbf{synchronous} (same time) and \textbf{asynchronous} (different times).
  \item Listing benefits without mentioning the matching challenges (the GCE expects balanced answers).
  \item Writing ``Internet'' where the question asks for the \emph{service} (e.g.\ telemedicine is a service, the Internet is the medium).
\end{itemize}

\textbf{5. GCE tips}
\begin{itemize}
  \item Define the term \emph{first}, then give a \textbf{Cameroonian example} (MTN Mobile Money, Orange Money, Jumia Cameroon, ANTIC, online GCE registration, Moodle at universities) — examples earn marks.
  \item For ``discuss'' questions, structure your answer: \textbf{definition $\rightarrow$ applications $\rightarrow$ benefits $\rightarrow$ challenges $\rightarrow$ solutions/conclusion}.
  \item Learn the acronyms with their full meaning: B2B, B2C, C2C, G2C, G2B, G2G, CAL, CAI, CBT, LMS, EMR.
  \item In tables comparing CAL and traditional learning, use clear criteria: pace, feedback, cost, availability, interaction.
  \item Manage your time: one well-developed point with an example is worth more than three vague points.
\end{itemize}
\end{bilan}

\findecours

\end{document}
